% Complete reviewed Unit 18; editorial master: chapters/ch18/student.md
% Markdown SHA-256: 3d54d8343f68988f385cbed184c6ad0b0d16f3672c5f150ef7691dfebd25e637
% One-use initial assembly followed by complete content and page review.
\newcommand{\runningtitle}{Unit 18. Prose consolidation}
% Layout inherited from the reviewed previous course; set a running title per unit.
% This is a small hand-maintained TeX layout, not a conversion pipeline.
\RequirePackage{fix-cm}
\documentclass[11pt,twoside]{article}
\usepackage{fontspec}
\setmainfont{Cambria}[
  Renderer=HarfBuzz,
  BoldFont=cambriab.ttf,
  ItalicFont=cambriai.ttf,
  BoldItalicFont=cambriaz.ttf
]
\setmonofont{lmmono10-regular.otf}[BoldFont=lmmonolt10-bold.otf]
\usepackage[paperwidth=165mm,paperheight=241mm,inner=21mm,outer=19mm,
  top=22mm,bottom=20mm,headheight=14pt,headsep=7mm,footskip=10mm]{geometry}
\usepackage{xcolor,array,booktabs,longtable,tabularx}
\usepackage{enumitem,fancyhdr}
\usepackage[unicode,hidelinks]{hyperref}
\definecolor{ink}{HTML}{000000}
\definecolor{muted}{HTML}{4C5558}
\setlength{\parindent}{0pt}
\setlength{\parskip}{3.5pt}
\AtBeginDocument{\fontsize{11}{14.3}\selectfont}
\setlength{\tabcolsep}{4pt}
\renewcommand{\arraystretch}{1.10}
\setlist[itemize]{leftmargin=1.3em,label=\textbullet,itemsep=2pt,topsep=3pt,parsep=0pt}
\setlist[enumerate]{leftmargin=1.8em,itemsep=3pt,topsep=4pt,parsep=0pt}
\widowpenalty=10000
\clubpenalty=10000
\emergencystretch=1em
\raggedbottom
\pagestyle{fancy}
\fancyhf{}
\providecommand{\runningtitle}{Old English}
\fancyhead[LE,RO]{\small\itshape \runningtitle}
\fancyhead[LO,RE]{\small Old English}
\fancyfoot[LE,RO]{\small\thepage}
\renewcommand{\headrulewidth}{0.3pt}
\newcommand{\booktitle}[2]{
  {\fontsize{9}{11}\selectfont\color{muted}#1\par}
  \vspace{3pt}
  {\fontsize{20}{24}\selectfont\color{ink}#2\par}
  \vspace{6pt}
}
\newcommand{\source}[1]{\par{\footnotesize\color{muted}#1\par}}
\newcommand{\readingtext}[1]{
  \par\addvspace{5pt}
  \begingroup\leftskip=1.5em
    \fontsize{12}{16}\selectfont #1\par
  \endgroup\addvspace{5pt}
}
\makeatletter
\renewcommand{\section}{\@startsection{section}{1}{\z@}
  {-2.4ex \@plus -0.6ex \@minus -0.2ex}{0.8ex \@plus 0.2ex}
  {\normalfont\fontsize{12}{15}\selectfont\bfseries\color{ink}}}
\renewcommand{\subsection}{\@startsection{subsection}{2}{\z@}
  {-1.8ex \@plus -0.5ex \@minus -0.2ex}{0.5ex \@plus 0.2ex}
  {\normalfont\fontsize{11}{14}\selectfont\bfseries\color{ink}}}
\makeatother

\hyphenpenalty=10000
\exhyphenpenalty=10000
\emergencystretch=3em
\newcommand{\itempara}[2]{\par\addvspace{3pt}\noindent\hangindent=1.8em\hangafter=1\makebox[1.8em][l]{#1.}#2\par}
\newcommand{\reserve}[1]{\par\penalty-100\begingroup\dimen0=\pagegoal\advance\dimen0 by-\pagetotal\ifdim\dimen0<#1\newpage\fi\endgroup}
\begin{document}
\addfontfeatures{Ligatures=TeXOff}
\raggedright
\booktitle{UNIT 18}{Independent prose consolidation}

% mdparagraph:P001:begin
A connected narrative changes speakers, refers back across clauses, and presents some actions as wishes, instructions or reports. Identify each clause’s participants and time relation before joining it to the whole.
% mdparagraph:P001:end

% mdparagraph:P002:begin
The final numbered Cook exercises form a sustained Joseph composition sequence. The original Old English reading continues the banquet scene with the princess’s gifts and the farewells. These are separate texts. The two translated examples and lettered tasks are new teaching material.
% mdparagraph:P002:end

% teaching: NEW-U18-DICTIONARY-REFERENCES
% teaching: NEW-U18-VERBAL-CONTEXT-HELP
% mdparagraph:P003:begin
\textbf{Forms and help.} A dictionary form identifies a word; it may still need inflection. The complete tables teach their named systems and stated meanings. Diagnostic stems are references to those systems, not assignments of unlisted moods. An expression labelled \textbf{Supplied} remains visible during the task: use its unfamiliar forms with the explanation given. Its independent recall is not required. In verb tables, 1, 2 and 3 denote person.
% mdparagraph:P003:end

\reserve{7\baselineskip}
% mdparagraph:P004:begin
\section*{1. Clauses, reference and reporting}
% mdparagraph:P004:end

% mdparagraph:P005:begin
A noun phrase takes its case from its own clause. A relative expression can have a different role from its antecedent. The particle ðe has no visible case ending; identify its verb and complements to establish its role.
% mdparagraph:P005:end

% teaching: NEW-U18-EXAMPLE-01-TEACHING
\reserve{5\baselineskip}
% mdparagraph:P006:begin
\textbf{Se mann ðe se cyning geseah, sang.}
% mdparagraph:P006:end

% mdparagraph:P007:begin
The man whom the king saw sang.
% mdparagraph:P007:end

% mdparagraph:P008:begin
The man is subject of sang, while the relative represents the object of geseah; the king is the relative clause subject.
% mdparagraph:P008:end

% teaching: NEW-U18-EXAMPLE-02-TEACHING
\reserve{5\baselineskip}
% mdparagraph:P009:begin
\textbf{Se cyning hēt ðone mann ðǣm cnihte ðā bōc sęllan.}
% mdparagraph:P009:end

% mdparagraph:P010:begin
The king ordered the man to give the servant the book.
% mdparagraph:P010:end

% mdparagraph:P011:begin
The accusative man performs the infinitive; the servant is its dative recipient and the book its accusative object.
% mdparagraph:P011:end

% mdparagraph:P012:begin
A third-person pronoun follows its referent’s grammatical gender, including that of an inanimate noun. Agreement can exclude a referent without selecting a unique person. There is no general nearest-noun rule.
% mdparagraph:P012:end

% mdparagraph:P013:begin
In direct speech, ic identifies the current speaker. A report revises that reference to fit its reporting clause. A preterite subjunctive under a past report does not itself prove the report false. Establish earlier, simultaneous or prospective time from the whole statement.
% mdparagraph:P013:end

% mdparagraph:P014:begin
With hātan and an infinitive, the accusative actor is the person ordered to act. A dative recipient of that action is a separate complement.
% mdparagraph:P014:end

\reserve{7\baselineskip}
% mdparagraph:P015:begin
\section*{2. Nouns and adjectives}
% mdparagraph:P015:end

% mdparagraph:P016:begin
Choose case from the phrase’s function: nominative subject, accusative direct object, genitive possessor or licensed complement, dative recipient or licensed complement. A preposition or particular verb can determine the case.
% mdparagraph:P016:end

% mdparagraph:P017:begin
Strong masculine nouns normally have genitive singular -es, dative -e, nominative/\allowbreak{}accusative plural -as, genitive plural -a and dative plural -um. Strong neuters share -es, -e, -a and -um; nominative and accusative agree, with plural -u or no ending according to the stem. The selected feminine oblique singular ends in -e. Stem changes and weak assa are given individually.
% mdparagraph:P017:end

% mdparagraph:P018:begin
A singular-only table gives all four cases of the stated mass or abstract meaning, using a familiar class. It does not claim that no plural is possible in other senses.
% mdparagraph:P018:end

\reserve{7\baselineskip}
% mdparagraph:P019:begin
\subsection*{sacc}
% mdparagraph:P019:end

% teaching: NEW-U18-VOC-SACC-U18-SENSE
\par\begin{minipage}{\linewidth}\raggedright
% mdparagraph:P020:begin
\textbf{sacc} — sack.
% mdparagraph:P020:end

% teaching: NEW-U18-NOM-SACC
% mdparagraph:P021:begin
Masculine; stem sacc-. Use the ordinary strong masculine endings.
% mdparagraph:P021:end

% table: NEW-U18-NOM-SACC-DISPLAY
% mdtable:T01:begin
\par\begingroup
\fontsize{10.5}{13.2}\selectfont
\setlength{\LTleft}{0pt}\setlength{\LTright}{\fill}
\noindent
\begin{tabular}{@{}>{\raggedright\arraybackslash}p{39.46mm}>{\raggedright\arraybackslash}p{39.46mm}>{\raggedright\arraybackslash}p{39.46mm}@{}}
\toprule
\textbf{Case} & \textbf{Singular} & \textbf{Plural} \\
\midrule
Nominative & sacc & saccas \\
Genitive & sacces & sacca \\
Dative & sacce & saccum \\
Accusative & sacc & saccas \\
\bottomrule
\end{tabular}
\par\endgroup

% mdtable:T01:end
\end{minipage}\par
\reserve{7\baselineskip}
% mdparagraph:P022:begin
\subsection*{assa}
% mdparagraph:P022:end

% teaching: NEW-U18-VOC-ASSA-U18-SENSE
\par\begin{minipage}{\linewidth}\raggedright
% mdparagraph:P023:begin
\textbf{assa} — ass; donkey.
% mdparagraph:P023:end

% teaching: NEW-U18-NOM-ASSA
% mdparagraph:P024:begin
Masculine; weak stem ass-. All oblique singular cases and nominative/\allowbreak{}accusative plural are assan; genitive plural is assena and dative plural assum.
% mdparagraph:P024:end

% table: NEW-U18-NOM-ASSA-DISPLAY
% mdtable:T02:begin
\par\begingroup
\fontsize{10.5}{13.2}\selectfont
\setlength{\LTleft}{0pt}\setlength{\LTright}{\fill}
\noindent
\begin{tabular}{@{}>{\raggedright\arraybackslash}p{39.46mm}>{\raggedright\arraybackslash}p{39.46mm}>{\raggedright\arraybackslash}p{39.46mm}@{}}
\toprule
\textbf{Case} & \textbf{Singular} & \textbf{Plural} \\
\midrule
Nominative & assa & assan \\
Genitive & assan & assena \\
Dative & assan & assum \\
Accusative & assan & assan \\
\bottomrule
\end{tabular}
\par\endgroup

% mdtable:T02:end
\end{minipage}\par
\reserve{7\baselineskip}
% mdparagraph:P025:begin
\subsection*{hunig}
% mdparagraph:P025:end

% teaching: NEW-U18-VOC-HUNIG-U18-SENSE
\par\begin{minipage}{\linewidth}\raggedright
% mdparagraph:P026:begin
\textbf{hunig} — honey.
% mdparagraph:P026:end

% teaching: NEW-U18-NOM-HUNIG
% mdparagraph:P027:begin
Neuter; stem hunig-. Retain the medial vowel before endings. This table gives the complete singular of the mass meaning honey.
% mdparagraph:P027:end

% table: NEW-U18-NOM-HUNIG-DISPLAY
% mdtable:T03:begin
\par\begingroup
\fontsize{10.5}{13.2}\selectfont
\setlength{\LTleft}{0pt}\setlength{\LTright}{\fill}
\noindent
\begin{tabular}{@{}>{\raggedright\arraybackslash}p{60.59mm}>{\raggedright\arraybackslash}p{60.59mm}@{}}
\toprule
\textbf{Case} & \textbf{Singular} \\
\midrule
Nominative & hunig \\
Genitive & huniges \\
Dative & hunige \\
Accusative & hunig \\
\bottomrule
\end{tabular}
\par\endgroup

% mdtable:T03:end
\end{minipage}\par
\reserve{7\baselineskip}
% mdparagraph:P028:begin
\subsection*{æcern}
% mdparagraph:P028:end

% teaching: NEW-U18-VOC-AECERN-U18-SENSE
\par\begin{minipage}{\linewidth}\raggedright
% mdparagraph:P029:begin
\textbf{æcern} — acorn.
% mdparagraph:P029:end

% teaching: NEW-U18-NOM-AECERN
% mdparagraph:P030:begin
Neuter; stem æcern-. The selected nominative/\allowbreak{}accusative plural is æcernu. Retain the medial vowel in the other cases: final rn makes the second syllable heavy. Preserve unmarked æ.
% mdparagraph:P030:end

% table: NEW-U18-NOM-AECERN-DISPLAY
% mdtable:T04:begin
\par\begingroup
\fontsize{10.5}{13.2}\selectfont
\setlength{\LTleft}{0pt}\setlength{\LTright}{\fill}
\noindent
\begin{tabular}{@{}>{\raggedright\arraybackslash}p{39.46mm}>{\raggedright\arraybackslash}p{39.46mm}>{\raggedright\arraybackslash}p{39.46mm}@{}}
\toprule
\textbf{Case} & \textbf{Singular} & \textbf{Plural} \\
\midrule
Nominative & æcern & æcernu \\
Genitive & æcernes & æcerna \\
Dative & æcerne & æcernum \\
Accusative & æcern & æcernu \\
\bottomrule
\end{tabular}
\par\endgroup

% mdtable:T04:end
\end{minipage}\par
\reserve{7\baselineskip}
% mdparagraph:P031:begin
\subsection*{hwǣte}
% mdparagraph:P031:end

% teaching: NEW-U18-VOC-HWAETE-U18-SENSE
\par\begin{minipage}{\linewidth}\raggedright
% mdparagraph:P032:begin
\textbf{hwǣte} — wheat.
% mdparagraph:P032:end

% teaching: NEW-U18-NOM-HWAETE
% mdparagraph:P033:begin
Masculine; stem hwǣt-. Nominative and accusative singular retain final -e. This table gives the complete singular of wheat.
% mdparagraph:P033:end

% table: NEW-U18-NOM-HWAETE-DISPLAY
% mdtable:T05:begin
\par\begingroup
\fontsize{10.5}{13.2}\selectfont
\setlength{\LTleft}{0pt}\setlength{\LTright}{\fill}
\noindent
\begin{tabular}{@{}>{\raggedright\arraybackslash}p{60.59mm}>{\raggedright\arraybackslash}p{60.59mm}@{}}
\toprule
\textbf{Case} & \textbf{Singular} \\
\midrule
Nominative & hwǣte \\
Genitive & hwǣtes \\
Dative & hwǣte \\
Accusative & hwǣte \\
\bottomrule
\end{tabular}
\par\endgroup

% mdtable:T05:end
\end{minipage}\par
\reserve{7\baselineskip}
% mdparagraph:P034:begin
\subsection*{læfel}
% mdparagraph:P034:end

% teaching: NEW-U18-VOC-LAEFEL-U18-SENSE
\par\begin{minipage}{\linewidth}\raggedright
% mdparagraph:P035:begin
\textbf{læfel} — cup.
% mdparagraph:P035:end

% teaching: NEW-U18-NOM-LAEFEL
% mdparagraph:P036:begin
Masculine; stem læfel-. Retain medial e before endings because the first syllable is short. Preserve unmarked æ.
% mdparagraph:P036:end

% table: NEW-U18-NOM-LAEFEL-DISPLAY
% mdtable:T06:begin
\par\begingroup
\fontsize{10.5}{13.2}\selectfont
\setlength{\LTleft}{0pt}\setlength{\LTright}{\fill}
\noindent
\begin{tabular}{@{}>{\raggedright\arraybackslash}p{39.46mm}>{\raggedright\arraybackslash}p{39.46mm}>{\raggedright\arraybackslash}p{39.46mm}@{}}
\toprule
\textbf{Case} & \textbf{Singular} & \textbf{Plural} \\
\midrule
Nominative & læfel & læfelas \\
Genitive & læfeles & læfela \\
Dative & læfele & læfelum \\
Accusative & læfel & læfelas \\
\bottomrule
\end{tabular}
\par\endgroup

% mdtable:T06:end
\end{minipage}\par
\reserve{7\baselineskip}
% mdparagraph:P037:begin
\subsection*{scēawere}
% mdparagraph:P037:end

% teaching: NEW-U18-VOC-SCEAWERE-U18-SENSE
\par\begin{minipage}{\linewidth}\raggedright
% mdparagraph:P038:begin
\textbf{scēawere} — spy.
% mdparagraph:P038:end

% teaching: NEW-U18-NOM-SCEAWERE
% mdparagraph:P039:begin
Masculine; stem scēawer-. Replace the dictionary form’s final e with the required ending; nominative/\allowbreak{}accusative plural is scēaweras.
% mdparagraph:P039:end

% table: NEW-U18-NOM-SCEAWERE-DISPLAY
% mdtable:T07:begin
\par\begingroup
\fontsize{10.5}{13.2}\selectfont
\setlength{\LTleft}{0pt}\setlength{\LTright}{\fill}
\noindent
\begin{tabular}{@{}>{\raggedright\arraybackslash}p{39.46mm}>{\raggedright\arraybackslash}p{39.46mm}>{\raggedright\arraybackslash}p{39.46mm}@{}}
\toprule
\textbf{Case} & \textbf{Singular} & \textbf{Plural} \\
\midrule
Nominative & scēawere & scēaweras \\
Genitive & scēaweres & scēawera \\
Dative & scēawere & scēawerum \\
Accusative & scēawere & scēaweras \\
\bottomrule
\end{tabular}
\par\endgroup

% mdtable:T07:end
\end{minipage}\par
\reserve{7\baselineskip}
% mdparagraph:P040:begin
\subsection*{ęlcung}
% mdparagraph:P040:end

% teaching: NEW-U18-VOC-ELCUNG-U18-SENSE
\par\begin{minipage}{\linewidth}\raggedright
% mdparagraph:P041:begin
\textbf{ęlcung} — delay.
% mdparagraph:P041:end

% teaching: NEW-U18-NOM-ELCUNG
% mdparagraph:P042:begin
Feminine; stem ęlcung-. This table selects -e in the oblique singular. The -ung class also admits -a variants. Only abstract singular delay is required.
% mdparagraph:P042:end

% table: NEW-U18-NOM-ELCUNG-DISPLAY
% mdtable:T08:begin
\par\begingroup
\fontsize{10.5}{13.2}\selectfont
\setlength{\LTleft}{0pt}\setlength{\LTright}{\fill}
\noindent
\begin{tabular}{@{}>{\raggedright\arraybackslash}p{60.59mm}>{\raggedright\arraybackslash}p{60.59mm}@{}}
\toprule
\textbf{Case} & \textbf{Singular} \\
\midrule
Nominative & ęlcung \\
Genitive & ęlcunge \\
Dative & ęlcunge \\
Accusative & ęlcunge \\
\bottomrule
\end{tabular}
\par\endgroup

% mdtable:T08:end
\end{minipage}\par
\reserve{7\baselineskip}
% mdparagraph:P043:begin
\subsection*{gebēorscipe}
% mdparagraph:P043:end

% teaching: NEW-U18-VOC-GEBEORSCIPE-U18-SENSE
\par\begin{minipage}{\linewidth}\raggedright
% mdparagraph:P044:begin
\textbf{gebēorscipe} — feast; banquet.
% mdparagraph:P044:end

% teaching: NEW-U18-NOM-GEBEORSCIPE
% mdparagraph:P045:begin
Masculine; stem gebēorscip-. Nominative and accusative singular retain -e; the plural follows the ordinary strong masculine pattern.
% mdparagraph:P045:end

% table: NEW-U18-NOM-GEBEORSCIPE-DISPLAY
% mdtable:T09:begin
\par\begingroup
\fontsize{10.5}{13.2}\selectfont
\setlength{\LTleft}{0pt}\setlength{\LTright}{\fill}
\noindent
\begin{tabular}{@{}>{\raggedright\arraybackslash}p{34.82mm}>{\raggedright\arraybackslash}p{41.78mm}>{\raggedright\arraybackslash}p{41.78mm}@{}}
\toprule
\textbf{Case} & \textbf{Singular} & \textbf{Plural} \\
\midrule
Nominative & gebēorscipe & gebēorscipas \\
Genitive & gebēorscipes & gebēorscipa \\
Dative & gebēorscipe & gebēorscipum \\
Accusative & gebēorscipe & gebēorscipas \\
\bottomrule
\end{tabular}
\par\endgroup

% mdtable:T09:end
\end{minipage}\par
\reserve{7\baselineskip}
% mdparagraph:P046:begin
\subsection*{sārnes}
% mdparagraph:P046:end

% teaching: NEW-U18-VOC-SARNES-U18-SENSE
\par\begin{minipage}{\linewidth}\raggedright
% mdparagraph:P047:begin
\textbf{sārnes} — grief.
% mdparagraph:P047:end

% teaching: NEW-U18-NOM-SARNES
% mdparagraph:P048:begin
Feminine; oblique stem sārness-. Double final s and retain the preceding e: sārnesse. Only abstract singular grief is required.
% mdparagraph:P048:end

% table: NEW-U18-NOM-SARNES-DISPLAY
% mdtable:T10:begin
\par\begingroup
\fontsize{10.5}{13.2}\selectfont
\setlength{\LTleft}{0pt}\setlength{\LTright}{\fill}
\noindent
\begin{tabular}{@{}>{\raggedright\arraybackslash}p{60.59mm}>{\raggedright\arraybackslash}p{60.59mm}@{}}
\toprule
\textbf{Case} & \textbf{Singular} \\
\midrule
Nominative & sārnes \\
Genitive & sārnesse \\
Dative & sārnesse \\
Accusative & sārnesse \\
\bottomrule
\end{tabular}
\par\endgroup

% mdtable:T10:end
\end{minipage}\par
\reserve{7\baselineskip}
% mdparagraph:P049:begin
\subsection*{ealdor}
% mdparagraph:P049:end

% teaching: NEW-U18-VOC-EALDOR-CHIEF-U14-COMPONENT-SOURCE
\par\begin{minipage}{\linewidth}\raggedright
% mdparagraph:P050:begin
\textbf{ealdor} — chief; ruler.
% mdparagraph:P050:end

% teaching: NEW-U18-NOM-EALDOR-CHIEF
% mdparagraph:P051:begin
Masculine; ealdor has the oblique stem ealdr-. Lose the unstressed vowel before an ending: ealdres, ealdre, ealdras.
% mdparagraph:P051:end

% table: NEW-U18-NOM-EALDOR-CHIEF-DISPLAY
% mdtable:T11:begin
\par\begingroup
\fontsize{10.5}{13.2}\selectfont
\setlength{\LTleft}{0pt}\setlength{\LTright}{\fill}
\noindent
\begin{tabular}{@{}>{\raggedright\arraybackslash}p{39.46mm}>{\raggedright\arraybackslash}p{39.46mm}>{\raggedright\arraybackslash}p{39.46mm}@{}}
\toprule
\textbf{Case} & \textbf{Singular} & \textbf{Plural} \\
\midrule
Nominative & ealdor & ealdras \\
Genitive & ealdres & ealdra \\
Dative & ealdre & ealdrum \\
Accusative & ealdor & ealdras \\
\bottomrule
\end{tabular}
\par\endgroup

% mdtable:T11:end
\end{minipage}\par
\reserve{7\baselineskip}
% mdparagraph:P052:begin
\subsection*{ðearf}
% mdparagraph:P052:end

% teaching: NEW-U18-VOC-THEARF-U15-SOURCE-SENSE
\par\begin{minipage}{\linewidth}\raggedright
% mdparagraph:P053:begin
\textbf{ðearf} — benefit; good.
% mdparagraph:P053:end

% teaching: NEW-U18-NOM-THEARF
% mdparagraph:P054:begin
Feminine; stem ðearf-. This unit uses the abstract singular benefit or good; the oblique singular is ðearfe.
% mdparagraph:P054:end

% table: NEW-U18-NOM-THEARF-DISPLAY
% mdtable:T12:begin
\par\begingroup
\fontsize{10.5}{13.2}\selectfont
\setlength{\LTleft}{0pt}\setlength{\LTright}{\fill}
\noindent
\begin{tabular}{@{}>{\raggedright\arraybackslash}p{60.59mm}>{\raggedright\arraybackslash}p{60.59mm}@{}}
\toprule
\textbf{Case} & \textbf{Singular} \\
\midrule
Nominative & ðearf \\
Genitive & ðearfe \\
Dative & ðearfe \\
Accusative & ðearfe \\
\bottomrule
\end{tabular}
\par\endgroup

% mdtable:T12:end
\end{minipage}\par
\reserve{7\baselineskip}
% mdparagraph:P055:begin
\subsection*{ðēow}
% mdparagraph:P055:end

% teaching: NEW-U18-VOC-THEOW-SERVANT-U18-SENSE
\par\begin{minipage}{\linewidth}\raggedright
% mdparagraph:P056:begin
\textbf{ðēow} — servant; enslaved person.
% mdparagraph:P056:end

% teaching: NEW-U18-NOM-THEOW-SERVANT
% mdparagraph:P057:begin
Masculine; strong stem ðēow-. Distinguish this noun from the earlier weak ðēowa. Its selected genitive plural is ðēowa, while the weak noun has ðēowena. The complete phrase in the reading has local help.
% mdparagraph:P057:end

% table: NEW-U18-NOM-THEOW-SERVANT-DISPLAY
% mdtable:T13:begin
\par\begingroup
\fontsize{10.5}{13.2}\selectfont
\setlength{\LTleft}{0pt}\setlength{\LTright}{\fill}
\noindent
\begin{tabular}{@{}>{\raggedright\arraybackslash}p{39.46mm}>{\raggedright\arraybackslash}p{39.46mm}>{\raggedright\arraybackslash}p{39.46mm}@{}}
\toprule
\textbf{Case} & \textbf{Singular} & \textbf{Plural} \\
\midrule
Nominative & ðēow & ðēowas \\
Genitive & ðēowes & ðēowa \\
Dative & ðēowe & ðēowum \\
Accusative & ðēow & ðēowas \\
\bottomrule
\end{tabular}
\par\endgroup

% mdtable:T13:end
\end{minipage}\par
\reserve{7\baselineskip}
% mdparagraph:P058:begin
\subsection*{hunger}
% mdparagraph:P058:end

% teaching: NEW-U18-VOC-HUNGER-U18-SENSE
\par\begin{minipage}{\linewidth}\raggedright
% mdparagraph:P059:begin
\textbf{hunger} — famine; hunger.
% mdparagraph:P059:end

% teaching: NEW-U18-NOM-HUNGER
% mdparagraph:P060:begin
Masculine; hunger has the oblique stem hungr-. Retain hunger without an ending; use hungres and hungre before endings. This table gives the mass singular famine or hunger.
% mdparagraph:P060:end

% table: NEW-U18-NOM-HUNGER-DISPLAY
% mdtable:T14:begin
\par\begingroup
\fontsize{10.5}{13.2}\selectfont
\setlength{\LTleft}{0pt}\setlength{\LTright}{\fill}
\noindent
\begin{tabular}{@{}>{\raggedright\arraybackslash}p{60.59mm}>{\raggedright\arraybackslash}p{60.59mm}@{}}
\toprule
\textbf{Case} & \textbf{Singular} \\
\midrule
Nominative & hunger \\
Genitive & hungres \\
Dative & hungre \\
Accusative & hunger \\
\bottomrule
\end{tabular}
\par\endgroup

% mdtable:T14:end
\end{minipage}\par
\reserve{7\baselineskip}
% mdparagraph:P061:begin
\subsection*{sīð}
% mdparagraph:P061:end

% teaching: NEW-U18-VOC-SITH-U18-SENSE
\par\begin{minipage}{\linewidth}\raggedright
% mdparagraph:P062:begin
\textbf{sīð} — time; occasion.
% mdparagraph:P062:end

% teaching: NEW-U18-NOM-SITH
% mdparagraph:P063:begin
Masculine; stem sīð-. These are the four ordinary cases. The second-time expression in 44.7 is instrumental and is supplied separately; its noun is not dative merely because both end in -e.
% mdparagraph:P063:end

% table: NEW-U18-NOM-SITH-DISPLAY
% mdtable:T15:begin
\par\begingroup
\fontsize{10.5}{13.2}\selectfont
\setlength{\LTleft}{0pt}\setlength{\LTright}{\fill}
\noindent
\begin{tabular}{@{}>{\raggedright\arraybackslash}p{39.46mm}>{\raggedright\arraybackslash}p{39.46mm}>{\raggedright\arraybackslash}p{39.46mm}@{}}
\toprule
\textbf{Case} & \textbf{Singular} & \textbf{Plural} \\
\midrule
Nominative & sīð & sīðas \\
Genitive & sīðes & sīða \\
Dative & sīðe & sīðum \\
Accusative & sīð & sīðas \\
\bottomrule
\end{tabular}
\par\endgroup

% mdtable:T15:end
\end{minipage}\par
\reserve{7\baselineskip}
% mdparagraph:P064:begin
\subsection*{lāc}
% mdparagraph:P064:end

% teaching: NEW-U18-VOC-LAC-U18-SENSE
\par\begin{minipage}{\linewidth}\raggedright
% mdparagraph:P065:begin
\textbf{lāc} — gift; present.
% mdparagraph:P065:end

% teaching: NEW-U18-NOM-LAC
% mdparagraph:P066:begin
Neuter; stem lāc-. The long vowel makes the stem heavy, so nominative/\allowbreak{}accusative plural has no ending.
% mdparagraph:P066:end

% table: NEW-U18-NOM-LAC-DISPLAY
% mdtable:T16:begin
\par\begingroup
\fontsize{10.5}{13.2}\selectfont
\setlength{\LTleft}{0pt}\setlength{\LTright}{\fill}
\noindent
\begin{tabular}{@{}>{\raggedright\arraybackslash}p{39.46mm}>{\raggedright\arraybackslash}p{39.46mm}>{\raggedright\arraybackslash}p{39.46mm}@{}}
\toprule
\textbf{Case} & \textbf{Singular} & \textbf{Plural} \\
\midrule
Nominative & lāc & lāc \\
Genitive & lāces & lāca \\
Dative & lāce & lācum \\
Accusative & lāc & lāc \\
\bottomrule
\end{tabular}
\par\endgroup

% mdtable:T16:end
\end{minipage}\par
\reserve{7\baselineskip}
% mdparagraph:P067:begin
\subsection*{fæt}
% mdparagraph:P067:end

% teaching: NEW-U18-VOC-FAET-U18-SENSE
\par\begin{minipage}{\linewidth}\raggedright
% mdparagraph:P068:begin
\textbf{fæt} — vessel; container.
% mdparagraph:P068:end

% teaching: NEW-U18-NOM-FAET
% mdparagraph:P069:begin
Neuter; retain æ throughout the singular and use a throughout the plural: fatu, fata, fatum. The short stem takes -u in nominative/\allowbreak{}accusative plural.
% mdparagraph:P069:end

% table: NEW-U18-NOM-FAET-DISPLAY
% mdtable:T17:begin
\par\begingroup
\fontsize{10.5}{13.2}\selectfont
\setlength{\LTleft}{0pt}\setlength{\LTright}{\fill}
\noindent
\begin{tabular}{@{}>{\raggedright\arraybackslash}p{39.46mm}>{\raggedright\arraybackslash}p{39.46mm}>{\raggedright\arraybackslash}p{39.46mm}@{}}
\toprule
\textbf{Case} & \textbf{Singular} & \textbf{Plural} \\
\midrule
Nominative & fæt & fatu \\
Genitive & fætes & fata \\
Dative & fæte & fatum \\
Accusative & fæt & fatu \\
\bottomrule
\end{tabular}
\par\endgroup

% mdtable:T17:end
\end{minipage}\par
\reserve{7\baselineskip}
% mdparagraph:P070:begin
\subsection*{weorð}
% mdparagraph:P070:end

% teaching: NEW-U18-VOC-WEORTH-PRICE-U18-SENSE
\par\begin{minipage}{\linewidth}\raggedright
% mdparagraph:P071:begin
\textbf{weorð} — price; value.
% mdparagraph:P071:end

% teaching: NEW-U18-NOM-WEORTH-PRICE
% mdparagraph:P072:begin
Neuter; stem weorð-. The table gives the complete singular of price or value.
% mdparagraph:P072:end

% table: NEW-U18-NOM-WEORTH-PRICE-DISPLAY
% mdtable:T18:begin
\par\begingroup
\fontsize{10.5}{13.2}\selectfont
\setlength{\LTleft}{0pt}\setlength{\LTright}{\fill}
\noindent
\begin{tabular}{@{}>{\raggedright\arraybackslash}p{60.59mm}>{\raggedright\arraybackslash}p{60.59mm}@{}}
\toprule
\textbf{Case} & \textbf{Singular} \\
\midrule
Nominative & weorð \\
Genitive & weorðes \\
Dative & weorðe \\
Accusative & weorð \\
\bottomrule
\end{tabular}
\par\endgroup

% mdtable:T18:end
\end{minipage}\par
\reserve{7\baselineskip}
% mdparagraph:P073:begin
\subsection*{feoh}
% mdparagraph:P073:end

% teaching: NEW-U18-VOC-FEOH-U18-SENSE
\par\begin{minipage}{\linewidth}\raggedright
% mdparagraph:P074:begin
\textbf{feoh} — money.
% mdparagraph:P074:end

% teaching: NEW-U18-NOM-FEOH
% mdparagraph:P075:begin
Neuter; money. Final h disappears before endings: genitive fēos and selected dative fēo. The source prints feoh without a macron before h; preserve that spelling. An absent printed macron does not assert that the historical vowel was short.
% mdparagraph:P075:end

% table: NEW-U18-NOM-FEOH-DISPLAY
% mdtable:T19:begin
\par\begingroup
\fontsize{10.5}{13.2}\selectfont
\setlength{\LTleft}{0pt}\setlength{\LTright}{\fill}
\noindent
\begin{tabular}{@{}>{\raggedright\arraybackslash}p{60.59mm}>{\raggedright\arraybackslash}p{60.59mm}@{}}
\toprule
\textbf{Case} & \textbf{Singular} \\
\midrule
Nominative & feoh \\
Genitive & fēos \\
Dative & fēo \\
Accusative & feoh \\
\bottomrule
\end{tabular}
\par\endgroup

% mdtable:T19:end
\end{minipage}\par
\reserve{7\baselineskip}
% mdparagraph:P076:begin
\subsection*{Strong and weak adjectives}
% mdparagraph:P076:end

% mdparagraph:P077:begin
Adjectives agree in gender, number and case. Use the strong pattern in its familiar undetermined uses and the weak pattern after the determiners that require it. The tables give every ordinary case in all three genders and both numbers. Strong masculine and neuter instrumental singular cells are recognition only; they introduce no new instrumental production.
% mdparagraph:P077:end

\reserve{7\baselineskip}
% mdparagraph:P078:begin
\subsection*{scyldig}
% mdparagraph:P078:end

% teaching: NEW-U18-VOC-SCYLDIG-U18-SENSE
\par\begin{minipage}{\linewidth}\raggedright
% mdparagraph:P079:begin
\textbf{scyldig} — guilty; culpable.
% mdparagraph:P079:end

% teaching: NEW-U18-ADJ-SCYLDIG
% mdparagraph:P080:begin
Use the uncontracted stem scyldig-. Strong feminine nominative singular and neuter nominative/\allowbreak{}accusative plural retain -u; contraction is a documented alternative.
% mdparagraph:P080:end

% mdparagraph:P081:begin
\textbf{scyldig — strong}
% mdparagraph:P081:end

% table: NEW-U18-ADJ-SCYLDIG-STRONG-DISPLAY
% mdtable:T20:begin
% course-table-style: 2026-10
% course-panel:T20-P01:begin
\par\begingroup
\fontsize{10.5}{13.2}\selectfont
\setlength{\LTleft}{0pt}\setlength{\LTright}{\fill}
\noindent
\begin{tabular}{@{}>{\raggedright\arraybackslash}p{22.0mm}>{\raggedright\arraybackslash}p{31.188mm}>{\raggedright\arraybackslash}p{31.188mm}>{\raggedright\arraybackslash}p{31.188mm}@{}}
\toprule
\textbf{Case} & \textbf{Masculine singular} & \textbf{Neuter singular} & \textbf{Feminine singular} \\
\midrule
Nominative & scyldig & scyldig & scyldigu \\
Genitive & scyldiges & scyldiges & scyldigre \\
Dative & scyldigum & scyldigum & scyldigre \\
Accusative & scyldigne & scyldig & scyldige \\
Instrumental & scyldige & scyldige & — \\
\bottomrule
\end{tabular}
\par\endgroup
% course-panel:T20-P01:end
\end{minipage}\par
\par\addvspace{5pt}
% course-panel:T20-P02:begin
\par\begingroup
\fontsize{10.5}{13.2}\selectfont
\setlength{\LTleft}{0pt}\setlength{\LTright}{\fill}
\noindent
\begin{tabular}{@{}>{\raggedright\arraybackslash}p{22.0mm}>{\raggedright\arraybackslash}p{31.188mm}>{\raggedright\arraybackslash}p{31.188mm}>{\raggedright\arraybackslash}p{31.188mm}@{}}
\toprule
\textbf{Case} & \textbf{Masculine plural} & \textbf{Neuter plural} & \textbf{Feminine plural} \\
\midrule
Nominative & scyldige & scyldigu & scyldiga \\
Genitive & scyldigra & scyldigra & scyldigra \\
Dative & scyldigum & scyldigum & scyldigum \\
Accusative & scyldige & scyldigu & scyldiga \\
Instrumental & — & — & — \\
\bottomrule
\end{tabular}
\par\endgroup
% course-panel:T20-P02:end
% mdtable:T20:end

\reserve{18\baselineskip}
\reserve{16\baselineskip}
% mdparagraph:P082:begin
\textbf{scyldig — weak}
% mdparagraph:P082:end

% table: NEW-U18-ADJ-SCYLDIG-WEAK-DISPLAY
% mdtable:T21:begin
% course-table-style: 2026-10
% course-panel:T21-P01:begin
\par\begingroup
\fontsize{10.5}{13.2}\selectfont
\setlength{\LTleft}{0pt}\setlength{\LTright}{\fill}
\noindent
\begin{tabular}{@{}>{\raggedright\arraybackslash}p{22.0mm}>{\raggedright\arraybackslash}p{31.188mm}>{\raggedright\arraybackslash}p{31.188mm}>{\raggedright\arraybackslash}p{31.188mm}@{}}
\toprule
\textbf{Case} & \textbf{Masculine singular} & \textbf{Neuter singular} & \textbf{Feminine singular} \\
\midrule
Nominative & scyldiga & scyldige & scyldige \\
Genitive & scyldigan & scyldigan & scyldigan \\
Dative & scyldigan & scyldigan & scyldigan \\
Accusative & scyldigan & scyldige & scyldigan \\
\bottomrule
\end{tabular}
\par\endgroup
% course-panel:T21-P01:end
\par\addvspace{5pt}
% course-panel:T21-P02:begin
\par\begingroup
\fontsize{10.5}{13.2}\selectfont
\setlength{\LTleft}{0pt}\setlength{\LTright}{\fill}
\noindent
\begin{tabular}{@{}>{\raggedright\arraybackslash}p{22.0mm}>{\raggedright\arraybackslash}p{31.188mm}>{\raggedright\arraybackslash}p{31.188mm}>{\raggedright\arraybackslash}p{31.188mm}@{}}
\toprule
\textbf{Case} & \textbf{Masculine plural} & \textbf{Neuter plural} & \textbf{Feminine plural} \\
\midrule
Nominative & scyldigan & scyldigan & scyldigan \\
Genitive & scyldigra & scyldigra & scyldigra \\
Dative & scyldigum & scyldigum & scyldigum \\
Accusative & scyldigan & scyldigan & scyldigan \\
\bottomrule
\end{tabular}
\par\endgroup
% course-panel:T21-P02:end
% mdtable:T21:end

\reserve{7\baselineskip}
% mdparagraph:P083:begin
\subsection*{twīfeald}
% mdparagraph:P083:end

% teaching: NEW-U18-VOC-TWIFEALD-U18-SENSE
\par\begin{minipage}{\linewidth}\raggedright
% mdparagraph:P084:begin
\textbf{twīfeald} — double; twofold.
% mdparagraph:P084:end

% teaching: NEW-U18-ADJ-TWIFEALD
% mdparagraph:P085:begin
The heavy stem twīfeald- takes no -u in strong feminine nominative singular or neuter nominative/\allowbreak{}accusative plural.
% mdparagraph:P085:end

% mdparagraph:P086:begin
\textbf{twīfeald — strong}
% mdparagraph:P086:end

% table: NEW-U18-ADJ-TWIFEALD-STRONG-DISPLAY
% mdtable:T22:begin
% course-table-style: 2026-10
% course-panel:T22-P01:begin
\par\begingroup
\fontsize{10.5}{13.2}\selectfont
\setlength{\LTleft}{0pt}\setlength{\LTright}{\fill}
\noindent
\begin{tabular}{@{}>{\raggedright\arraybackslash}p{22.0mm}>{\raggedright\arraybackslash}p{31.188mm}>{\raggedright\arraybackslash}p{31.188mm}>{\raggedright\arraybackslash}p{31.188mm}@{}}
\toprule
\textbf{Case} & \textbf{Masculine singular} & \textbf{Neuter singular} & \textbf{Feminine singular} \\
\midrule
Nominative & twīfeald & twīfeald & twīfeald \\
Genitive & twīfealdes & twīfealdes & twīfealdre \\
Dative & twīfealdum & twīfealdum & twīfealdre \\
Accusative & twīfealdne & twīfeald & twīfealde \\
Instrumental & twīfealde & twīfealde & — \\
\bottomrule
\end{tabular}
\par\endgroup
% course-panel:T22-P01:end
\end{minipage}\par
\par\addvspace{5pt}
% course-panel:T22-P02:begin
\par\begingroup
\fontsize{10.5}{13.2}\selectfont
\setlength{\LTleft}{0pt}\setlength{\LTright}{\fill}
\noindent
\begin{tabular}{@{}>{\raggedright\arraybackslash}p{22.0mm}>{\raggedright\arraybackslash}p{31.188mm}>{\raggedright\arraybackslash}p{31.188mm}>{\raggedright\arraybackslash}p{31.188mm}@{}}
\toprule
\textbf{Case} & \textbf{Masculine plural} & \textbf{Neuter plural} & \textbf{Feminine plural} \\
\midrule
Nominative & twīfealde & twīfeald & twīfealda \\
Genitive & twīfealdra & twīfealdra & twīfealdra \\
Dative & twīfealdum & twīfealdum & twīfealdum \\
Accusative & twīfealde & twīfeald & twīfealda \\
Instrumental & — & — & — \\
\bottomrule
\end{tabular}
\par\endgroup
% course-panel:T22-P02:end
% mdtable:T22:end

\reserve{18\baselineskip}
\reserve{16\baselineskip}
% mdparagraph:P087:begin
\textbf{twīfeald — weak}
% mdparagraph:P087:end

% table: NEW-U18-ADJ-TWIFEALD-WEAK-DISPLAY
% mdtable:T23:begin
% course-table-style: 2026-10
% course-panel:T23-P01:begin
\par\begingroup
\fontsize{10.5}{13.2}\selectfont
\setlength{\LTleft}{0pt}\setlength{\LTright}{\fill}
\noindent
\begin{tabular}{@{}>{\raggedright\arraybackslash}p{22.0mm}>{\raggedright\arraybackslash}p{31.188mm}>{\raggedright\arraybackslash}p{31.188mm}>{\raggedright\arraybackslash}p{31.188mm}@{}}
\toprule
\textbf{Case} & \textbf{Masculine singular} & \textbf{Neuter singular} & \textbf{Feminine singular} \\
\midrule
Nominative & twīfealda & twīfealde & twīfealde \\
Genitive & twīfealdan & twīfealdan & twīfealdan \\
Dative & twīfealdan & twīfealdan & twīfealdan \\
Accusative & twīfealdan & twīfealde & twīfealdan \\
\bottomrule
\end{tabular}
\par\endgroup
% course-panel:T23-P01:end
\par\addvspace{5pt}
% course-panel:T23-P02:begin
\par\begingroup
\fontsize{10.5}{13.2}\selectfont
\setlength{\LTleft}{0pt}\setlength{\LTright}{\fill}
\noindent
\begin{tabular}{@{}>{\raggedright\arraybackslash}p{22.0mm}>{\raggedright\arraybackslash}p{31.188mm}>{\raggedright\arraybackslash}p{31.188mm}>{\raggedright\arraybackslash}p{31.188mm}@{}}
\toprule
\textbf{Case} & \textbf{Masculine plural} & \textbf{Neuter plural} & \textbf{Feminine plural} \\
\midrule
Nominative & twīfealdan & twīfealdan & twīfealdan \\
Genitive & twīfealdra & twīfealdra & twīfealdra \\
Dative & twīfealdum & twīfealdum & twīfealdum \\
Accusative & twīfealdan & twīfealdan & twīfealdan \\
\bottomrule
\end{tabular}
\par\endgroup
% course-panel:T23-P02:end
% mdtable:T23:end

\reserve{7\baselineskip}
% mdparagraph:P088:begin
\subsection*{gesund}
% mdparagraph:P088:end

% teaching: NEW-U18-VOC-GESUND-U18-SENSE
\par\begin{minipage}{\linewidth}\raggedright
% mdparagraph:P089:begin
\textbf{gesund} — well; in good health.
% mdparagraph:P089:end

% teaching: NEW-U18-ADJ-GESUND
% mdparagraph:P090:begin
The heavy stem gesund- takes no -u in strong feminine nominative singular or neuter nominative/\allowbreak{}accusative plural.
% mdparagraph:P090:end

% mdparagraph:P091:begin
\textbf{gesund — strong}
% mdparagraph:P091:end

% table: NEW-U18-ADJ-GESUND-STRONG-DISPLAY
% mdtable:T24:begin
% course-table-style: 2026-10
% course-panel:T24-P01:begin
\par\begingroup
\fontsize{10.5}{13.2}\selectfont
\setlength{\LTleft}{0pt}\setlength{\LTright}{\fill}
\noindent
\begin{tabular}{@{}>{\raggedright\arraybackslash}p{22.0mm}>{\raggedright\arraybackslash}p{31.188mm}>{\raggedright\arraybackslash}p{31.188mm}>{\raggedright\arraybackslash}p{31.188mm}@{}}
\toprule
\textbf{Case} & \textbf{Masculine singular} & \textbf{Neuter singular} & \textbf{Feminine singular} \\
\midrule
Nominative & gesund & gesund & gesund \\
Genitive & gesundes & gesundes & gesundre \\
Dative & gesundum & gesundum & gesundre \\
Accusative & gesundne & gesund & gesunde \\
Instrumental & gesunde & gesunde & — \\
\bottomrule
\end{tabular}
\par\endgroup
% course-panel:T24-P01:end
\end{minipage}\par
\par\addvspace{5pt}
% course-panel:T24-P02:begin
\par\begingroup
\fontsize{10.5}{13.2}\selectfont
\setlength{\LTleft}{0pt}\setlength{\LTright}{\fill}
\noindent
\begin{tabular}{@{}>{\raggedright\arraybackslash}p{22.0mm}>{\raggedright\arraybackslash}p{31.188mm}>{\raggedright\arraybackslash}p{31.188mm}>{\raggedright\arraybackslash}p{31.188mm}@{}}
\toprule
\textbf{Case} & \textbf{Masculine plural} & \textbf{Neuter plural} & \textbf{Feminine plural} \\
\midrule
Nominative & gesunde & gesund & gesunda \\
Genitive & gesundra & gesundra & gesundra \\
Dative & gesundum & gesundum & gesundum \\
Accusative & gesunde & gesund & gesunda \\
Instrumental & — & — & — \\
\bottomrule
\end{tabular}
\par\endgroup
% course-panel:T24-P02:end
% mdtable:T24:end

\reserve{18\baselineskip}
\reserve{16\baselineskip}
% mdparagraph:P092:begin
\textbf{gesund — weak}
% mdparagraph:P092:end

% table: NEW-U18-ADJ-GESUND-WEAK-DISPLAY
% mdtable:T25:begin
% course-table-style: 2026-10
% course-panel:T25-P01:begin
\par\begingroup
\fontsize{10.5}{13.2}\selectfont
\setlength{\LTleft}{0pt}\setlength{\LTright}{\fill}
\noindent
\begin{tabular}{@{}>{\raggedright\arraybackslash}p{22.0mm}>{\raggedright\arraybackslash}p{31.188mm}>{\raggedright\arraybackslash}p{31.188mm}>{\raggedright\arraybackslash}p{31.188mm}@{}}
\toprule
\textbf{Case} & \textbf{Masculine singular} & \textbf{Neuter singular} & \textbf{Feminine singular} \\
\midrule
Nominative & gesunda & gesunde & gesunde \\
Genitive & gesundan & gesundan & gesundan \\
Dative & gesundan & gesundan & gesundan \\
Accusative & gesundan & gesunde & gesundan \\
\bottomrule
\end{tabular}
\par\endgroup
% course-panel:T25-P01:end
\par\addvspace{5pt}
% course-panel:T25-P02:begin
\par\begingroup
\fontsize{10.5}{13.2}\selectfont
\setlength{\LTleft}{0pt}\setlength{\LTright}{\fill}
\noindent
\begin{tabular}{@{}>{\raggedright\arraybackslash}p{22.0mm}>{\raggedright\arraybackslash}p{31.188mm}>{\raggedright\arraybackslash}p{31.188mm}>{\raggedright\arraybackslash}p{31.188mm}@{}}
\toprule
\textbf{Case} & \textbf{Masculine plural} & \textbf{Neuter plural} & \textbf{Feminine plural} \\
\midrule
Nominative & gesundan & gesundan & gesundan \\
Genitive & gesundra & gesundra & gesundra \\
Dative & gesundum & gesundum & gesundum \\
Accusative & gesundan & gesundan & gesundan \\
\bottomrule
\end{tabular}
\par\endgroup
% course-panel:T25-P02:end
% mdtable:T25:end

\reserve{7\baselineskip}
% mdparagraph:P093:begin
\subsection*{hāl}
% mdparagraph:P093:end

% teaching: NEW-U18-VOC-HAL-U18-HEALTH
\par\begin{minipage}{\linewidth}\raggedright
% mdparagraph:P094:begin
\textbf{hāl} — hale; healthy.
% mdparagraph:P094:end

% teaching: NEW-U18-ADJ-HAL
% mdparagraph:P095:begin
The heavy stem hāl- takes no -u in strong feminine nominative singular or neuter nominative/\allowbreak{}accusative plural. The meaning here is health.
% mdparagraph:P095:end

% mdparagraph:P096:begin
\textbf{hāl — strong}
% mdparagraph:P096:end

% table: NEW-U18-ADJ-HAL-STRONG-DISPLAY
% mdtable:T26:begin
% course-table-style: 2026-10
% course-panel:T26-P01:begin
\par\begingroup
\fontsize{10.5}{13.2}\selectfont
\setlength{\LTleft}{0pt}\setlength{\LTright}{\fill}
\noindent
\begin{tabular}{@{}>{\raggedright\arraybackslash}p{22.0mm}>{\raggedright\arraybackslash}p{31.188mm}>{\raggedright\arraybackslash}p{31.188mm}>{\raggedright\arraybackslash}p{31.188mm}@{}}
\toprule
\textbf{Case} & \textbf{Masculine singular} & \textbf{Neuter singular} & \textbf{Feminine singular} \\
\midrule
Nominative & Hāl & hāl & hāl \\
Genitive & hāles & hāles & hālre \\
Dative & hālum & hālum & hālre \\
Accusative & hālne & hāl & hāle \\
Instrumental & hāle & hāle & — \\
\bottomrule
\end{tabular}
\par\endgroup
% course-panel:T26-P01:end
\end{minipage}\par
\par\addvspace{5pt}
% course-panel:T26-P02:begin
\par\begingroup
\fontsize{10.5}{13.2}\selectfont
\setlength{\LTleft}{0pt}\setlength{\LTright}{\fill}
\noindent
\begin{tabular}{@{}>{\raggedright\arraybackslash}p{22.0mm}>{\raggedright\arraybackslash}p{31.188mm}>{\raggedright\arraybackslash}p{31.188mm}>{\raggedright\arraybackslash}p{31.188mm}@{}}
\toprule
\textbf{Case} & \textbf{Masculine plural} & \textbf{Neuter plural} & \textbf{Feminine plural} \\
\midrule
Nominative & hāle & hāl & hāla \\
Genitive & hālra & hālra & hālra \\
Dative & hālum & hālum & hālum \\
Accusative & hāle & hāl & hāla \\
Instrumental & — & — & — \\
\bottomrule
\end{tabular}
\par\endgroup
% course-panel:T26-P02:end
% mdtable:T26:end

\reserve{18\baselineskip}
\reserve{16\baselineskip}
% mdparagraph:P097:begin
\textbf{hāl — weak}
% mdparagraph:P097:end

% table: NEW-U18-ADJ-HAL-WEAK-DISPLAY
% mdtable:T27:begin
% course-table-style: 2026-10
% course-panel:T27-P01:begin
\par\begingroup
\fontsize{10.5}{13.2}\selectfont
\setlength{\LTleft}{0pt}\setlength{\LTright}{\fill}
\noindent
\begin{tabular}{@{}>{\raggedright\arraybackslash}p{22.0mm}>{\raggedright\arraybackslash}p{31.188mm}>{\raggedright\arraybackslash}p{31.188mm}>{\raggedright\arraybackslash}p{31.188mm}@{}}
\toprule
\textbf{Case} & \textbf{Masculine singular} & \textbf{Neuter singular} & \textbf{Feminine singular} \\
\midrule
Nominative & hāla & hāle & hāle \\
Genitive & hālan & hālan & hālan \\
Dative & hālan & hālan & hālan \\
Accusative & hālan & hāle & hālan \\
\bottomrule
\end{tabular}
\par\endgroup
% course-panel:T27-P01:end
\par\addvspace{5pt}
% course-panel:T27-P02:begin
\par\begingroup
\fontsize{10.5}{13.2}\selectfont
\setlength{\LTleft}{0pt}\setlength{\LTright}{\fill}
\noindent
\begin{tabular}{@{}>{\raggedright\arraybackslash}p{22.0mm}>{\raggedright\arraybackslash}p{31.188mm}>{\raggedright\arraybackslash}p{31.188mm}>{\raggedright\arraybackslash}p{31.188mm}@{}}
\toprule
\textbf{Case} & \textbf{Masculine plural} & \textbf{Neuter plural} & \textbf{Feminine plural} \\
\midrule
Nominative & hālan & hālan & hālan \\
Genitive & hālra & hālra & hālra \\
Dative & hālum & hālum & hālum \\
Accusative & hālan & hālan & hālan \\
\bottomrule
\end{tabular}
\par\endgroup
% course-panel:T27-P02:end
% mdtable:T27:end

\reserve{7\baselineskip}
% mdparagraph:P098:begin
\subsection*{āgen}
% mdparagraph:P098:end

% teaching: NEW-U18-VOC-AGEN-U16-SENSE
\par\begin{minipage}{\linewidth}\raggedright
% mdparagraph:P099:begin
\textbf{āgen} — own.
% mdparagraph:P099:end

% teaching: NEW-U18-ADJ-AGEN
% mdparagraph:P100:begin
Use āgn- before vowel-initial endings but āgen- before consonantal endings: compare āgnes, āgnum, āgna, āgnu with āgenne, āgenre, āgenra.
% mdparagraph:P100:end

% mdparagraph:P101:begin
\textbf{āgen — strong}
% mdparagraph:P101:end

% table: NEW-U18-ADJ-AGEN-STRONG-DISPLAY
% mdtable:T28:begin
% course-table-style: 2026-10
% course-panel:T28-P01:begin
\par\begingroup
\fontsize{10.5}{13.2}\selectfont
\setlength{\LTleft}{0pt}\setlength{\LTright}{\fill}
\noindent
\begin{tabular}{@{}>{\raggedright\arraybackslash}p{22.0mm}>{\raggedright\arraybackslash}p{31.188mm}>{\raggedright\arraybackslash}p{31.188mm}>{\raggedright\arraybackslash}p{31.188mm}@{}}
\toprule
\textbf{Case} & \textbf{Masculine singular} & \textbf{Neuter singular} & \textbf{Feminine singular} \\
\midrule
Nominative & āgen & āgen & āgnu \\
Genitive & āgnes & āgnes & āgenre \\
Dative & āgnum & āgnum & āgenre \\
Accusative & āgenne & āgen & āgne \\
Instrumental & āgne & āgne & — \\
\bottomrule
\end{tabular}
\par\endgroup
% course-panel:T28-P01:end
\end{minipage}\par
\par\addvspace{5pt}
% course-panel:T28-P02:begin
\par\begingroup
\fontsize{10.5}{13.2}\selectfont
\setlength{\LTleft}{0pt}\setlength{\LTright}{\fill}
\noindent
\begin{tabular}{@{}>{\raggedright\arraybackslash}p{22.0mm}>{\raggedright\arraybackslash}p{31.188mm}>{\raggedright\arraybackslash}p{31.188mm}>{\raggedright\arraybackslash}p{31.188mm}@{}}
\toprule
\textbf{Case} & \textbf{Masculine plural} & \textbf{Neuter plural} & \textbf{Feminine plural} \\
\midrule
Nominative & āgne & āgnu & āgna \\
Genitive & āgenra & āgenra & āgenra \\
Dative & āgnum & āgnum & āgnum \\
Accusative & āgne & āgnu & āgna \\
Instrumental & — & — & — \\
\bottomrule
\end{tabular}
\par\endgroup
% course-panel:T28-P02:end
% mdtable:T28:end

\reserve{18\baselineskip}
\reserve{16\baselineskip}
% mdparagraph:P102:begin
\textbf{āgen — weak}
% mdparagraph:P102:end

% table: NEW-U18-ADJ-AGEN-WEAK-DISPLAY
% mdtable:T29:begin
% course-table-style: 2026-10
% course-panel:T29-P01:begin
\par\begingroup
\fontsize{10.5}{13.2}\selectfont
\setlength{\LTleft}{0pt}\setlength{\LTright}{\fill}
\noindent
\begin{tabular}{@{}>{\raggedright\arraybackslash}p{22.0mm}>{\raggedright\arraybackslash}p{31.188mm}>{\raggedright\arraybackslash}p{31.188mm}>{\raggedright\arraybackslash}p{31.188mm}@{}}
\toprule
\textbf{Case} & \textbf{Masculine singular} & \textbf{Neuter singular} & \textbf{Feminine singular} \\
\midrule
Nominative & āgna & āgne & āgne \\
Genitive & āgnan & āgnan & āgnan \\
Dative & āgnan & āgnan & āgnan \\
Accusative & āgnan & āgne & āgnan \\
\bottomrule
\end{tabular}
\par\endgroup
% course-panel:T29-P01:end
\par\addvspace{5pt}
% course-panel:T29-P02:begin
\par\begingroup
\fontsize{10.5}{13.2}\selectfont
\setlength{\LTleft}{0pt}\setlength{\LTright}{\fill}
\noindent
\begin{tabular}{@{}>{\raggedright\arraybackslash}p{22.0mm}>{\raggedright\arraybackslash}p{31.188mm}>{\raggedright\arraybackslash}p{31.188mm}>{\raggedright\arraybackslash}p{31.188mm}@{}}
\toprule
\textbf{Case} & \textbf{Masculine plural} & \textbf{Neuter plural} & \textbf{Feminine plural} \\
\midrule
Nominative & āgnan & āgnan & āgnan \\
Genitive & āgenra & āgenra & āgenra \\
Dative & āgnum & āgnum & āgnum \\
Accusative & āgnan & āgnan & āgnan \\
\bottomrule
\end{tabular}
\par\endgroup
% course-panel:T29-P02:end
% mdtable:T29:end

\reserve{7\baselineskip}
% mdparagraph:P103:begin
\subsection*{clǣne}
% mdparagraph:P103:end

% teaching: NEW-U18-VOC-CLAENE-U18-BLAMELESS
\par\begin{minipage}{\linewidth}\raggedright
% mdparagraph:P104:begin
\textbf{clǣne} — clean; pure; blameless.
% mdparagraph:P104:end

% teaching: NEW-U18-ADJ-CLAENE
% mdparagraph:P105:begin
Use clǣn- before endings. Retain dictionary -e where no other ending is present. In the accusation, the sense is blameless or free from guilt.
% mdparagraph:P105:end

% mdparagraph:P106:begin
\textbf{clǣne — strong}
% mdparagraph:P106:end

% table: NEW-U18-ADJ-CLAENE-STRONG-DISPLAY
% mdtable:T30:begin
% course-table-style: 2026-10
% course-panel:T30-P01:begin
\par\begingroup
\fontsize{10.5}{13.2}\selectfont
\setlength{\LTleft}{0pt}\setlength{\LTright}{\fill}
\noindent
\begin{tabular}{@{}>{\raggedright\arraybackslash}p{22.0mm}>{\raggedright\arraybackslash}p{31.188mm}>{\raggedright\arraybackslash}p{31.188mm}>{\raggedright\arraybackslash}p{31.188mm}@{}}
\toprule
\textbf{Case} & \textbf{Masculine singular} & \textbf{Neuter singular} & \textbf{Feminine singular} \\
\midrule
Nominative & clǣne & clǣne & clǣnu \\
Genitive & clǣnes & clǣnes & clǣnre \\
Dative & clǣnum & clǣnum & clǣnre \\
Accusative & clǣnne & clǣne & clǣne \\
Instrumental & clǣne & clǣne & — \\
\bottomrule
\end{tabular}
\par\endgroup
% course-panel:T30-P01:end
\end{minipage}\par
\par\addvspace{5pt}
% course-panel:T30-P02:begin
\par\begingroup
\fontsize{10.5}{13.2}\selectfont
\setlength{\LTleft}{0pt}\setlength{\LTright}{\fill}
\noindent
\begin{tabular}{@{}>{\raggedright\arraybackslash}p{22.0mm}>{\raggedright\arraybackslash}p{31.188mm}>{\raggedright\arraybackslash}p{31.188mm}>{\raggedright\arraybackslash}p{31.188mm}@{}}
\toprule
\textbf{Case} & \textbf{Masculine plural} & \textbf{Neuter plural} & \textbf{Feminine plural} \\
\midrule
Nominative & clǣne & clǣnu & clǣna \\
Genitive & clǣnra & clǣnra & clǣnra \\
Dative & clǣnum & clǣnum & clǣnum \\
Accusative & clǣne & clǣnu & clǣna \\
Instrumental & — & — & — \\
\bottomrule
\end{tabular}
\par\endgroup
% course-panel:T30-P02:end
% mdtable:T30:end

\reserve{18\baselineskip}
\reserve{16\baselineskip}
% mdparagraph:P107:begin
\textbf{clǣne — weak}
% mdparagraph:P107:end

% table: NEW-U18-ADJ-CLAENE-WEAK-DISPLAY
% mdtable:T31:begin
% course-table-style: 2026-10
% course-panel:T31-P01:begin
\par\begingroup
\fontsize{10.5}{13.2}\selectfont
\setlength{\LTleft}{0pt}\setlength{\LTright}{\fill}
\noindent
\begin{tabular}{@{}>{\raggedright\arraybackslash}p{22.0mm}>{\raggedright\arraybackslash}p{31.188mm}>{\raggedright\arraybackslash}p{31.188mm}>{\raggedright\arraybackslash}p{31.188mm}@{}}
\toprule
\textbf{Case} & \textbf{Masculine singular} & \textbf{Neuter singular} & \textbf{Feminine singular} \\
\midrule
Nominative & clǣna & clǣne & clǣne \\
Genitive & clǣnan & clǣnan & clǣnan \\
Dative & clǣnan & clǣnan & clǣnan \\
Accusative & clǣnan & clǣne & clǣnan \\
\bottomrule
\end{tabular}
\par\endgroup
% course-panel:T31-P01:end
\par\addvspace{5pt}
% course-panel:T31-P02:begin
\par\begingroup
\fontsize{10.5}{13.2}\selectfont
\setlength{\LTleft}{0pt}\setlength{\LTright}{\fill}
\noindent
\begin{tabular}{@{}>{\raggedright\arraybackslash}p{22.0mm}>{\raggedright\arraybackslash}p{31.188mm}>{\raggedright\arraybackslash}p{31.188mm}>{\raggedright\arraybackslash}p{31.188mm}@{}}
\toprule
\textbf{Case} & \textbf{Masculine plural} & \textbf{Neuter plural} & \textbf{Feminine plural} \\
\midrule
Nominative & clǣnan & clǣnan & clǣnan \\
Genitive & clǣnra & clǣnra & clǣnra \\
Dative & clǣnum & clǣnum & clǣnum \\
Accusative & clǣnan & clǣnan & clǣnan \\
\bottomrule
\end{tabular}
\par\endgroup
% course-panel:T31-P02:end
% mdtable:T31:end

\reserve{7\baselineskip}
% mdparagraph:P108:begin
\subsection*{ęlðēodig}
% mdparagraph:P108:end

% teaching: NEW-U18-VOC-ELTHEODIG-U18-FOREIGN
\par\begin{minipage}{\linewidth}\raggedright
% mdparagraph:P109:begin
\textbf{ęlðēodig} — foreign.
% mdparagraph:P109:end

% teaching: NEW-U18-ADJ-ELTHEODIG
% mdparagraph:P110:begin
Use the uncontracted stem ęlðēodig-. Strong feminine nominative singular and neuter nominative/\allowbreak{}accusative plural retain -u.
% mdparagraph:P110:end

% mdparagraph:P111:begin
\textbf{ęlðēodig — strong}
% mdparagraph:P111:end

% table: NEW-U18-ADJ-ELTHEODIG-STRONG-DISPLAY
% mdtable:T32:begin
% course-table-style: 2026-10
% course-panel:T32-P01:begin
\par\begingroup
\fontsize{10.5}{13.2}\selectfont
\setlength{\LTleft}{0pt}\setlength{\LTright}{\fill}
\noindent
\begin{tabular}{@{}>{\raggedright\arraybackslash}p{22.0mm}>{\raggedright\arraybackslash}p{31.188mm}>{\raggedright\arraybackslash}p{31.188mm}>{\raggedright\arraybackslash}p{31.188mm}@{}}
\toprule
\textbf{Case} & \textbf{Masculine singular} & \textbf{Neuter singular} & \textbf{Feminine singular} \\
\midrule
Nominative & ęlðēodig & ęlðēodig & ęlðēodigu \\
Genitive & ęlðēodiges & ęlðēodiges & ęlðēodigre \\
Dative & ęlðēodigum & ęlðēodigum & ęlðēodigre \\
Accusative & ęlðēodigne & ęlðēodig & ęlðēodige \\
Instrumental & ęlðēodige & ęlðēodige & — \\
\bottomrule
\end{tabular}
\par\endgroup
% course-panel:T32-P01:end
\end{minipage}\par
\par\addvspace{5pt}
% course-panel:T32-P02:begin
\par\begingroup
\fontsize{10.5}{13.2}\selectfont
\setlength{\LTleft}{0pt}\setlength{\LTright}{\fill}
\noindent
\begin{tabular}{@{}>{\raggedright\arraybackslash}p{22.0mm}>{\raggedright\arraybackslash}p{31.188mm}>{\raggedright\arraybackslash}p{31.188mm}>{\raggedright\arraybackslash}p{31.188mm}@{}}
\toprule
\textbf{Case} & \textbf{Masculine plural} & \textbf{Neuter plural} & \textbf{Feminine plural} \\
\midrule
Nominative & ęlðēodige & ęlðēodigu & ęlðēodiga \\
Genitive & ęlðēodigra & ęlðēodigra & ęlðēodigra \\
Dative & ęlðēodigum & ęlðēodigum & ęlðēodigum \\
Accusative & ęlðēodige & ęlðēodigu & ęlðēodiga \\
Instrumental & — & — & — \\
\bottomrule
\end{tabular}
\par\endgroup
% course-panel:T32-P02:end
% mdtable:T32:end

\reserve{18\baselineskip}
\reserve{16\baselineskip}
% mdparagraph:P112:begin
\textbf{ęlðēodig — weak}
% mdparagraph:P112:end

% table: NEW-U18-ADJ-ELTHEODIG-WEAK-DISPLAY
% mdtable:T33:begin
% course-table-style: 2026-10
% course-panel:T33-P01:begin
\par\begingroup
\fontsize{10.5}{13.2}\selectfont
\setlength{\LTleft}{0pt}\setlength{\LTright}{\fill}
\noindent
\begin{tabular}{@{}>{\raggedright\arraybackslash}p{22.0mm}>{\raggedright\arraybackslash}p{31.188mm}>{\raggedright\arraybackslash}p{31.188mm}>{\raggedright\arraybackslash}p{31.188mm}@{}}
\toprule
\textbf{Case} & \textbf{Masculine singular} & \textbf{Neuter singular} & \textbf{Feminine singular} \\
\midrule
Nominative & ęlðēodiga & ęlðēodige & ęlðēodige \\
Genitive & ęlðēodigan & ęlðēodigan & ęlðēodigan \\
Dative & ęlðēodigan & ęlðēodigan & ęlðēodigan \\
Accusative & ęlðēodigan & ęlðēodige & ęlðēodigan \\
\bottomrule
\end{tabular}
\par\endgroup
% course-panel:T33-P01:end
\par\addvspace{5pt}
% course-panel:T33-P02:begin
\par\begingroup
\fontsize{10.5}{13.2}\selectfont
\setlength{\LTleft}{0pt}\setlength{\LTright}{\fill}
\noindent
\begin{tabular}{@{}>{\raggedright\arraybackslash}p{22.0mm}>{\raggedright\arraybackslash}p{31.188mm}>{\raggedright\arraybackslash}p{31.188mm}>{\raggedright\arraybackslash}p{31.188mm}@{}}
\toprule
\textbf{Case} & \textbf{Masculine plural} & \textbf{Neuter plural} & \textbf{Feminine plural} \\
\midrule
Nominative & ęlðēodigan & ęlðēodigan & ęlðēodigan \\
Genitive & ęlðēodigra & ęlðēodigra & ęlðēodigra \\
Dative & ęlðēodigum & ęlðēodigum & ęlðēodigum \\
Accusative & ęlðēodigan & ęlðēodigan & ęlðēodigan \\
\bottomrule
\end{tabular}
\par\endgroup
% course-panel:T33-P02:end
% mdtable:T33:end

\reserve{7\baselineskip}
% mdparagraph:P113:begin
\section*{3. Verbal systems}
% mdparagraph:P113:end

% mdparagraph:P114:begin
Each heading states the complete system added or reviewed. Indicative tables distinguish present and preterite. Subjunctive tables give all persons of the named tense; identical forms do not remove person distinctions. Commands distinguish singular and plural addressees. Infinitive and participle pairs add no unlisted finite forms.
% mdparagraph:P114:end

% mdparagraph:P115:begin
A modal takes a plain infinitive. The inflected infinitive requires its stated construction, commonly with tō. A base past participle may occur in the specified passive or result predicate without a new adjective declension. Meaning and complementation remain separate from the form.
% mdparagraph:P115:end

\reserve{7\baselineskip}
% mdparagraph:P116:begin
\subsection*{beðurfan: present and preterite indicative}
% mdparagraph:P116:end

% teaching: NEW-U18-VOC-BETHURFAN-U18-SELECTED
\par\begin{minipage}{\linewidth}\raggedright
% mdparagraph:P117:begin
\textbf{beðurfan} — need; require.
% mdparagraph:P117:end

% teaching: u18-bethurfan-indicative12
% mdparagraph:P118:begin
This preterite-present has present singular beðearf, beðearft, beðearf and plural beðurfon. The dental preterite uses beðorfte- in the singular and beðorft- in the plural. The complete indicative is assigned; the headword alone does not introduce an infinitive construction or subjunctive.
% mdparagraph:P118:end

% mdparagraph:P119:begin
\textbf{Diagnostic forms:} beðearf (present indicative first singular); beðurfon (present indicative plural); beðorfte (preterite indicative first singular); beðorfton (preterite indicative plural).
% mdparagraph:P119:end

% table: u18-bethurfan-indicative12-INDICATIVE-DISPLAY
% mdtable:T34:begin
% course-table-style: 2026-10
% course-panel:T34-P01:begin
\par\begingroup
\fontsize{10.5}{13.2}\selectfont
\setlength{\LTleft}{0pt}\setlength{\LTright}{\fill}
\noindent
\begin{tabular}{@{}>{\raggedright\arraybackslash}p{39.459mm}>{\raggedright\arraybackslash}p{39.459mm}>{\raggedright\arraybackslash}p{39.459mm}@{}}
\toprule
\textbf{Person and number} & \textbf{Present} & \textbf{Preterite} \\
\midrule
1 sg. & beðearf & beðorfte \\
2 sg. & beðearft & beðorftest \\
3 sg. & beðearf & beðorfte \\
1–3 pl. & beðurfon & beðorfton \\
\bottomrule
\end{tabular}
\par\endgroup
% course-panel:T34-P01:end
% mdtable:T34:end
\end{minipage}\par
\reserve{7\baselineskip}
% mdparagraph:P120:begin
\subsection*{oncnāwan: present and preterite indicative}
% mdparagraph:P120:end

% teaching: NEW-U18-VOC-ONCNAWAN-U18-SELECTED
\par\begin{minipage}{\linewidth}\raggedright
% mdparagraph:P121:begin
\textbf{oncnāwan} — recognize; know by recognition.
% mdparagraph:P121:end

% teaching: u18-oncnawan-indicative12
% mdparagraph:P122:begin
The present stem oncnāw- changes ā to ǣ in the second and third singular. Preterite oncnēow- adds -e in the second singular; first and third singular have no ending. The recognized person is accusative.
% mdparagraph:P122:end

% mdparagraph:P123:begin
\textbf{Diagnostic forms:} oncnāwe (present indicative first singular); oncnāwað (present indicative plural); oncnēow (preterite indicative first singular); oncnēowon (preterite indicative plural).
% mdparagraph:P123:end

% table: u18-oncnawan-indicative12-INDICATIVE-DISPLAY
% mdtable:T35:begin
% course-table-style: 2026-10
% course-panel:T35-P01:begin
\par\begingroup
\fontsize{10.5}{13.2}\selectfont
\setlength{\LTleft}{0pt}\setlength{\LTright}{\fill}
\noindent
\begin{tabular}{@{}>{\raggedright\arraybackslash}p{39.459mm}>{\raggedright\arraybackslash}p{39.459mm}>{\raggedright\arraybackslash}p{39.459mm}@{}}
\toprule
\textbf{Person and number} & \textbf{Present} & \textbf{Preterite} \\
\midrule
1 sg. & oncnāwe & oncnēow \\
2 sg. & oncnǣwst & oncnēowe \\
3 sg. & oncnǣwð & oncnēow \\
1–3 pl. & oncnāwað & oncnēowon \\
\bottomrule
\end{tabular}
\par\endgroup
% course-panel:T35-P01:end
% mdtable:T35:end
\end{minipage}\par
\reserve{7\baselineskip}
% mdparagraph:P124:begin
\subsection*{undōn: present and preterite indicative}
% mdparagraph:P124:end

% teaching: NEW-U18-VOC-UNDON-U18-SELECTED
\par\begin{minipage}{\linewidth}\raggedright
% mdparagraph:P125:begin
\textbf{undōn} — undo; open.
% mdparagraph:P125:end

% teaching: u18-undon-indicative12
% mdparagraph:P126:begin
Retain un- with the irregular dōn stems: undō, undēst, undēð, undōð; preterite undyde, undydest, undydon. The accusative object is the tied container opened.
% mdparagraph:P126:end

% mdparagraph:P127:begin
\textbf{Diagnostic forms:} undō (present indicative first singular); undōð (present indicative plural); undyde (preterite indicative first singular); undydon (preterite indicative plural).
% mdparagraph:P127:end

% table: u18-undon-indicative12-INDICATIVE-DISPLAY
% mdtable:T36:begin
% course-table-style: 2026-10
% course-panel:T36-P01:begin
\par\begingroup
\fontsize{10.5}{13.2}\selectfont
\setlength{\LTleft}{0pt}\setlength{\LTright}{\fill}
\noindent
\begin{tabular}{@{}>{\raggedright\arraybackslash}p{39.459mm}>{\raggedright\arraybackslash}p{39.459mm}>{\raggedright\arraybackslash}p{39.459mm}@{}}
\toprule
\textbf{Person and number} & \textbf{Present} & \textbf{Preterite} \\
\midrule
1 sg. & undō & undyde \\
2 sg. & undēst & undydest \\
3 sg. & undēð & undyde \\
1–3 pl. & undōð & undydon \\
\bottomrule
\end{tabular}
\par\endgroup
% course-panel:T36-P01:end
% mdtable:T36:end
\end{minipage}\par
\reserve{7\baselineskip}
% mdparagraph:P128:begin
\subsection*{lęcgean: present and preterite indicative}
% mdparagraph:P128:end

% teaching: NEW-U18-VOC-LECGEAN-U18-SELECTED
\par\begin{minipage}{\linewidth}\raggedright
% mdparagraph:P129:begin
\textbf{lęcgean} — lay; place.
% mdparagraph:P129:end

% teaching: u18-lecgean-indicative12
% mdparagraph:P130:begin
The short weak verb has lęcge, lęgest, lęgeð and plural lęcgeað. The doubled consonant is absent in the second and third singular. Present cg alternates with preterite g: lęgde, lęgdest, lęgdon. These tables select the uncontracted past; lēde/\allowbreak{}lēdon are documented variants. The reading’s imperative lęcgeað is separately supplied.
% mdparagraph:P130:end

% mdparagraph:P131:begin
\textbf{Diagnostic forms:} lęcge (present indicative first singular); lęcgeað (present indicative plural); lęgde (preterite indicative first singular); lęgdon (preterite indicative plural).
% mdparagraph:P131:end

% table: u18-lecgean-indicative12-INDICATIVE-DISPLAY
% mdtable:T37:begin
% course-table-style: 2026-10
% course-panel:T37-P01:begin
\par\begingroup
\fontsize{10.5}{13.2}\selectfont
\setlength{\LTleft}{0pt}\setlength{\LTright}{\fill}
\noindent
\begin{tabular}{@{}>{\raggedright\arraybackslash}p{39.459mm}>{\raggedright\arraybackslash}p{39.459mm}>{\raggedright\arraybackslash}p{39.459mm}@{}}
\toprule
\textbf{Person and number} & \textbf{Present} & \textbf{Preterite} \\
\midrule
1 sg. & lęcge & lęgde \\
2 sg. & lęgest & lęgdest \\
3 sg. & lęgeð & lęgde \\
1–3 pl. & lęcgeað & lęgdon \\
\bottomrule
\end{tabular}
\par\endgroup
% course-panel:T37-P01:end
% mdtable:T37:end
\end{minipage}\par
\reserve{7\baselineskip}
% mdparagraph:P132:begin
\subsection*{foresęcgean: present and preterite indicative}
% mdparagraph:P132:end

% teaching: NEW-U18-VOC-FORESECGEAN-U18-SELECTED
\par\begin{minipage}{\linewidth}\raggedright
% mdparagraph:P133:begin
\textbf{foresęcgean} — mention before; tell of beforehand.
% mdparagraph:P133:end

% teaching: u18-foresecgean-indicative12
% mdparagraph:P134:begin
Add fore- to the irregular sęcgan pattern. Select plural foresęcgeað, second/\allowbreak{}third singular foresægst and foresægð, and preterite foresægde, foresægdest, foresægdon. Retain unmarked æ before gd; contracted sāde is a separate alternative. Here the verb means mention earlier, not predict the future.
% mdparagraph:P134:end

% mdparagraph:P135:begin
\textbf{Diagnostic forms:} foresęcge (present indicative first singular); foresęcgeað (present indicative plural); foresægde (preterite indicative first singular); foresægdon (preterite indicative plural).
% mdparagraph:P135:end

% table: u18-foresecgean-indicative12-INDICATIVE-DISPLAY
% mdtable:T38:begin
% course-table-style: 2026-10
% course-panel:T38-P01:begin
\par\begingroup
\fontsize{10.5}{13.2}\selectfont
\setlength{\LTleft}{0pt}\setlength{\LTright}{\fill}
\noindent
\begin{tabular}{@{}>{\raggedright\arraybackslash}p{39.459mm}>{\raggedright\arraybackslash}p{39.459mm}>{\raggedright\arraybackslash}p{39.459mm}@{}}
\toprule
\textbf{Person and number} & \textbf{Present} & \textbf{Preterite} \\
\midrule
1 sg. & foresęcge & foresægde \\
2 sg. & foresægst & foresægdest \\
3 sg. & foresægð & foresægde \\
1–3 pl. & foresęcgeað & foresægdon \\
\bottomrule
\end{tabular}
\par\endgroup
% course-panel:T38-P01:end
% mdtable:T38:end
\end{minipage}\par
\reserve{7\baselineskip}
% mdparagraph:P136:begin
\subsection*{forstelan: present and preterite indicative}
% mdparagraph:P136:end

% teaching: NEW-U18-VOC-FORSTELAN-U18-SELECTED
\par\begin{minipage}{\linewidth}\raggedright
% mdparagraph:P137:begin
\textbf{forstelan} — steal away; steal.
% mdparagraph:P137:end

% teaching: u18-forstelan-indicative12
% mdparagraph:P138:begin
Strong class IV alternates e, æ, ǣ, o. Present second and third singular have i: forstilst, forstilð. Preterite first/\allowbreak{}third singular is forstæl, second singular forstǣle, and plural forstǣlon. Keep lexical for-; unprefixed stelan is listed separately.
% mdparagraph:P138:end

% mdparagraph:P139:begin
\textbf{Diagnostic forms:} forstele (present indicative first singular); forstelað (present indicative plural); forstæl (preterite indicative first singular); forstǣlon (preterite indicative plural).
% mdparagraph:P139:end

% table: u18-forstelan-indicative12-INDICATIVE-DISPLAY
% mdtable:T39:begin
% course-table-style: 2026-10
% course-panel:T39-P01:begin
\par\begingroup
\fontsize{10.5}{13.2}\selectfont
\setlength{\LTleft}{0pt}\setlength{\LTright}{\fill}
\noindent
\begin{tabular}{@{}>{\raggedright\arraybackslash}p{39.459mm}>{\raggedright\arraybackslash}p{39.459mm}>{\raggedright\arraybackslash}p{39.459mm}@{}}
\toprule
\textbf{Person and number} & \textbf{Present} & \textbf{Preterite} \\
\midrule
1 sg. & forstele & forstæl \\
2 sg. & forstilst & forstǣle \\
3 sg. & forstilð & forstæl \\
1–3 pl. & forstelað & forstǣlon \\
\bottomrule
\end{tabular}
\par\endgroup
% course-panel:T39-P01:end
% mdtable:T39:end
\end{minipage}\par
\reserve{7\baselineskip}
% mdparagraph:P140:begin
\subsection*{cēapian: indicative and infinitives}
% mdparagraph:P140:end

% teaching: NEW-U18-VOC-CEAPIAN-U18-SELECTED
% mdparagraph:P141:begin
\textbf{cēapian} — trade; traffic; buy through trade.
% mdparagraph:P141:end

% teaching: u18-ceapian-indicative12-infinitive2
% mdparagraph:P142:begin
Weak class II has present -ie, -ast, -að, -iað and preterite -ode, -odest, -odon. Retain ēa. The complete infinitive pair is cēapian/\allowbreak{}cēapianne; the special buying-purpose construction remains supplied.
% mdparagraph:P142:end

% mdparagraph:P143:begin
\textbf{Diagnostic forms:} cēapie (present indicative first singular); cēapiað (present indicative plural); cēapode (preterite indicative first singular); cēapodon (preterite indicative plural); cēapian (plain\_infinitive); cēapianne (inflected\_infinitive).
% mdparagraph:P143:end

\reserve{13\baselineskip}
\reserve{11\baselineskip}
% mdparagraph:P144:begin
\textbf{cēapian — indicative}
% mdparagraph:P144:end

% table: u18-ceapian-indicative12-infinitive2-INDICATIVE-DISPLAY
% mdtable:T40:begin
% course-table-style: 2026-10
% course-panel:T40-P01:begin
\par\begingroup
\fontsize{10.5}{13.2}\selectfont
\setlength{\LTleft}{0pt}\setlength{\LTright}{\fill}
\noindent
\begin{tabular}{@{}>{\raggedright\arraybackslash}p{39.459mm}>{\raggedright\arraybackslash}p{39.459mm}>{\raggedright\arraybackslash}p{39.459mm}@{}}
\toprule
\textbf{Person and number} & \textbf{Present} & \textbf{Preterite} \\
\midrule
1 sg. & cēapie & cēapode \\
2 sg. & cēapast & cēapodest \\
3 sg. & cēapað & cēapode \\
1–3 pl. & cēapiað & cēapodon \\
\bottomrule
\end{tabular}
\par\endgroup
% course-panel:T40-P01:end
% mdtable:T40:end
\reserve{9\baselineskip}
\reserve{7\baselineskip}
% mdparagraph:P145:begin
\textbf{cēapian — additional bounded forms}
% mdparagraph:P145:end

% table: u18-ceapian-indicative12-infinitive2-BOUNDED-NONFINITE-DISPLAY
% mdtable:T41:begin
\par\begingroup
\fontsize{10.5}{13.2}\selectfont
\setlength{\LTleft}{0pt}\setlength{\LTright}{\fill}
\noindent
\begin{tabular}{@{}>{\raggedright\arraybackslash}p{71.72mm}>{\raggedright\arraybackslash}p{49.46mm}@{}}
\toprule
\textbf{Type} & \textbf{Form} \\
\midrule
Plain infinitive & cēapian \\
Inflected infinitive & cēapianne \\
\bottomrule
\end{tabular}
\par\endgroup

% mdtable:T41:end
\reserve{7\baselineskip}
% mdparagraph:P146:begin
\subsection*{gedōn: present subjunctive}
% mdparagraph:P146:end

% teaching: NEW-U18-VOC-GEDON-U18-SELECTED
\par\begin{minipage}{\linewidth}\raggedright
% mdparagraph:P147:begin
\textbf{gedōn} — make; cause to be.
% mdparagraph:P147:end

% teaching: u18-gedon-present-subjunctive6
% mdparagraph:P148:begin
Present subjunctive is gedō in all singular persons and gedōn in all plural persons. The causative construction takes an accusative person and an agreeing adjective.
% mdparagraph:P148:end

% mdparagraph:P149:begin
\textbf{Diagnostic forms:} gedō (present subjunctive sg); gedōn (present subjunctive pl).
% mdparagraph:P149:end

% table: u18-gedon-present-subjunctive6-SUBJUNCTIVE-DISPLAY
% mdtable:T42:begin
% course-table-style: 2026-10
% course-panel:T42-P01:begin
\par\begingroup
\fontsize{10.5}{13.2}\selectfont
\setlength{\LTleft}{0pt}\setlength{\LTright}{\fill}
\noindent
\begin{tabular}{@{}>{\raggedright\arraybackslash}p{63.342mm}>{\raggedright\arraybackslash}p{57.846mm}@{}}
\toprule
\textbf{Person and number} & \textbf{Present} \\
\midrule
1 sg. & gedō \\
2 sg. & gedō \\
3 sg. & gedō \\
1–3 pl. & gedōn \\
\bottomrule
\end{tabular}
\par\endgroup
% course-panel:T42-P01:end
% mdtable:T42:end
\end{minipage}\par
\reserve{7\baselineskip}
% mdparagraph:P150:begin
\subsection*{āgiefan: present subjunctive}
% mdparagraph:P150:end

% teaching: NEW-U18-VOC-AGIEFAN-U18-SELECTED
\par\begin{minipage}{\linewidth}\raggedright
% mdparagraph:P151:begin
\textbf{āgiefan} — give back; restore.
% mdparagraph:P151:end

% teaching: u18-agiefan-present-subjunctive6
% mdparagraph:P152:begin
Retain ā- with unchanged gief-. Present subjunctive singular ends in -e and plural in -en. The returned person is accusative and the recipient dative.
% mdparagraph:P152:end

% mdparagraph:P153:begin
\textbf{Diagnostic forms:} āgiefe (present subjunctive sg); āgiefen (present subjunctive pl).
% mdparagraph:P153:end

% table: u18-agiefan-present-subjunctive6-SUBJUNCTIVE-DISPLAY
% mdtable:T43:begin
% course-table-style: 2026-10
% course-panel:T43-P01:begin
\par\begingroup
\fontsize{10.5}{13.2}\selectfont
\setlength{\LTleft}{0pt}\setlength{\LTright}{\fill}
\noindent
\begin{tabular}{@{}>{\raggedright\arraybackslash}p{63.342mm}>{\raggedright\arraybackslash}p{57.846mm}@{}}
\toprule
\textbf{Person and number} & \textbf{Present} \\
\midrule
1 sg. & āgiefe \\
2 sg. & āgiefe \\
3 sg. & āgiefe \\
1–3 pl. & āgiefen \\
\bottomrule
\end{tabular}
\par\endgroup
% course-panel:T43-P01:end
% mdtable:T43:end
\end{minipage}\par
\reserve{7\baselineskip}
% mdparagraph:P154:begin
\subsection*{hlystan: present subjunctive}
% mdparagraph:P154:end

% teaching: NEW-U18-VOC-HLYSTAN-U18-SELECTED
\par\begin{minipage}{\linewidth}\raggedright
% mdparagraph:P155:begin
\textbf{hlystan} — listen to; heed.
% mdparagraph:P155:end

% teaching: u18-hlystan-present-subjunctive6
% mdparagraph:P156:begin
Long weak hlyst- takes present subjunctive -e/\allowbreak{}-en. Speech or content heeded is genitive.
% mdparagraph:P156:end

% mdparagraph:P157:begin
\textbf{Diagnostic forms:} hlyste (present subjunctive sg); hlysten (present subjunctive pl).
% mdparagraph:P157:end

% table: u18-hlystan-present-subjunctive6-SUBJUNCTIVE-DISPLAY
% mdtable:T44:begin
% course-table-style: 2026-10
% course-panel:T44-P01:begin
\par\begingroup
\fontsize{10.5}{13.2}\selectfont
\setlength{\LTleft}{0pt}\setlength{\LTright}{\fill}
\noindent
\begin{tabular}{@{}>{\raggedright\arraybackslash}p{63.342mm}>{\raggedright\arraybackslash}p{57.846mm}@{}}
\toprule
\textbf{Person and number} & \textbf{Present} \\
\midrule
1 sg. & hlyste \\
2 sg. & hlyste \\
3 sg. & hlyste \\
1–3 pl. & hlysten \\
\bottomrule
\end{tabular}
\par\endgroup
% course-panel:T44-P01:end
% mdtable:T44:end
\end{minipage}\par
\reserve{7\baselineskip}
% mdparagraph:P158:begin
\subsection*{ðēowian: present subjunctive}
% mdparagraph:P158:end

% teaching: NEW-U18-VOC-THEOWIAN-U18-SELECTED
\par\begin{minipage}{\linewidth}\raggedright
% mdparagraph:P159:begin
\textbf{ðēowian} — enslave; reduce to servitude.
% mdparagraph:P159:end

% teaching: u18-theowian-present-subjunctive6
% mdparagraph:P160:begin
Weak class II present subjunctive has -ie in the singular and -ien in the plural. The selected sense is enslave, with an accusative person. The source’s -ige variant is the same ending.
% mdparagraph:P160:end

% mdparagraph:P161:begin
\textbf{Diagnostic forms:} ðēowie (present subjunctive sg); ðēowien (present subjunctive pl).
% mdparagraph:P161:end

% table: u18-theowian-present-subjunctive6-SUBJUNCTIVE-DISPLAY
% mdtable:T45:begin
% course-table-style: 2026-10
% course-panel:T45-P01:begin
\par\begingroup
\fontsize{10.5}{13.2}\selectfont
\setlength{\LTleft}{0pt}\setlength{\LTright}{\fill}
\noindent
\begin{tabular}{@{}>{\raggedright\arraybackslash}p{63.342mm}>{\raggedright\arraybackslash}p{57.846mm}@{}}
\toprule
\textbf{Person and number} & \textbf{Present} \\
\midrule
1 sg. & ðēowie \\
2 sg. & ðēowie \\
3 sg. & ðēowie \\
1–3 pl. & ðēowien \\
\bottomrule
\end{tabular}
\par\endgroup
% course-panel:T45-P01:end
% mdtable:T45:end
\end{minipage}\par
\reserve{7\baselineskip}
% mdparagraph:P162:begin
\subsection*{berēafian: present subjunctive}
% mdparagraph:P162:end

% teaching: NEW-U18-VOC-BEREAFIAN-U18-SELECTED
\par\begin{minipage}{\linewidth}\raggedright
% mdparagraph:P163:begin
\textbf{berēafian} — despoil; deprive of.
% mdparagraph:P163:end

% teaching: u18-bereafian-present-subjunctive6
% mdparagraph:P164:begin
Weak class II gives present subjunctive -ie/\allowbreak{}-ien. The deprivation complement is explained below.
% mdparagraph:P164:end

% mdparagraph:P165:begin
\textbf{Diagnostic forms:} berēafie (present subjunctive sg); berēafien (present subjunctive pl).
% mdparagraph:P165:end

% table: u18-bereafian-present-subjunctive6-SUBJUNCTIVE-DISPLAY
% mdtable:T46:begin
% course-table-style: 2026-10
% course-panel:T46-P01:begin
\par\begingroup
\fontsize{10.5}{13.2}\selectfont
\setlength{\LTleft}{0pt}\setlength{\LTright}{\fill}
\noindent
\begin{tabular}{@{}>{\raggedright\arraybackslash}p{63.342mm}>{\raggedright\arraybackslash}p{57.846mm}@{}}
\toprule
\textbf{Person and number} & \textbf{Present} \\
\midrule
1 sg. & berēafie \\
2 sg. & berēafie \\
3 sg. & berēafie \\
1–3 pl. & berēafien \\
\bottomrule
\end{tabular}
\par\endgroup
% course-panel:T46-P01:end
% mdtable:T46:end
\end{minipage}\par
\reserve{7\baselineskip}
% mdparagraph:P166:begin
\subsection*{gemiltsian: present subjunctive}
% mdparagraph:P166:end

% teaching: u18-gemiltsian-present-subjunctive6
% mdparagraph:P167:begin
This extends the familiar weak class II verb to present subjunctive -ie/\allowbreak{}-ien. The person receiving mercy remains dative.
% mdparagraph:P167:end

\reserve{13\baselineskip}
% mdparagraph:P168:begin
\textbf{Diagnostic forms:} gemiltsie (present subjunctive sg); gemiltsien (present subjunctive pl).
% mdparagraph:P168:end

% table: u18-gemiltsian-present-subjunctive6-SUBJUNCTIVE-DISPLAY
% mdtable:T47:begin
% course-table-style: 2026-10
% course-panel:T47-P01:begin
\par\begingroup
\fontsize{10.5}{13.2}\selectfont
\setlength{\LTleft}{0pt}\setlength{\LTright}{\fill}
\noindent
\begin{tabular}{@{}>{\raggedright\arraybackslash}p{63.342mm}>{\raggedright\arraybackslash}p{57.846mm}@{}}
\toprule
\textbf{Person and number} & \textbf{Present} \\
\midrule
1 sg. & gemiltsie \\
2 sg. & gemiltsie \\
3 sg. & gemiltsie \\
1–3 pl. & gemiltsien \\
\bottomrule
\end{tabular}
\par\endgroup
% course-panel:T47-P01:end
% mdtable:T47:end
\reserve{7\baselineskip}
% mdparagraph:P169:begin
\subsection*{gesēon: present subjunctive}
% mdparagraph:P169:end

% teaching: u18-geseon-present-subjunctive6
% mdparagraph:P170:begin
The contracted present subjunctive is gesēo in the singular and gesēon in the plural. Only this mood and tense are added here.
% mdparagraph:P170:end

\reserve{13\baselineskip}
% mdparagraph:P171:begin
\textbf{Diagnostic forms:} gesēo (present subjunctive sg); gesēon (present subjunctive pl).
% mdparagraph:P171:end

% table: u18-geseon-present-subjunctive6-SUBJUNCTIVE-DISPLAY
% mdtable:T48:begin
% course-table-style: 2026-10
% course-panel:T48-P01:begin
\par\begingroup
\fontsize{10.5}{13.2}\selectfont
\setlength{\LTleft}{0pt}\setlength{\LTright}{\fill}
\noindent
\begin{tabular}{@{}>{\raggedright\arraybackslash}p{63.342mm}>{\raggedright\arraybackslash}p{57.846mm}@{}}
\toprule
\textbf{Person and number} & \textbf{Present} \\
\midrule
1 sg. & gesēo \\
2 sg. & gesēo \\
3 sg. & gesēo \\
1–3 pl. & gesēon \\
\bottomrule
\end{tabular}
\par\endgroup
% course-panel:T48-P01:end
% mdtable:T48:end
\reserve{7\baselineskip}
% mdparagraph:P172:begin
\subsection*{wunian: present subjunctive}
% mdparagraph:P172:end

% teaching: u18-wunian-present-subjunctive6
% mdparagraph:P173:begin
The familiar weak class II verb adds present subjunctive -ie/\allowbreak{}-ien. Dwelling in a place uses on with dative.
% mdparagraph:P173:end

\reserve{13\baselineskip}
% mdparagraph:P174:begin
\textbf{Diagnostic forms:} wunie (present subjunctive sg); wunien (present subjunctive pl).
% mdparagraph:P174:end

% table: u18-wunian-present-subjunctive6-SUBJUNCTIVE-DISPLAY
% mdtable:T49:begin
% course-table-style: 2026-10
% course-panel:T49-P01:begin
\par\begingroup
\fontsize{10.5}{13.2}\selectfont
\setlength{\LTleft}{0pt}\setlength{\LTright}{\fill}
\noindent
\begin{tabular}{@{}>{\raggedright\arraybackslash}p{63.342mm}>{\raggedright\arraybackslash}p{57.846mm}@{}}
\toprule
\textbf{Person and number} & \textbf{Present} \\
\midrule
1 sg. & wunie \\
2 sg. & wunie \\
3 sg. & wunie \\
1–3 pl. & wunien \\
\bottomrule
\end{tabular}
\par\endgroup
% course-panel:T49-P01:end
% mdtable:T49:end
\reserve{7\baselineskip}
% mdparagraph:P175:begin
\subsection*{bringan: present and preterite subjunctive}
% mdparagraph:P175:end

% teaching: u18-bringan-subjunctive12
% mdparagraph:P176:begin
Present subjunctive is bringe/\allowbreak{}bringen; preterite subjunctive brōhte/\allowbreak{}brōhten. The irregular dental past changes g to h before the dental group. Each tense has one form for all three persons within a number.
% mdparagraph:P176:end

\reserve{14\baselineskip}
% mdparagraph:P177:begin
\textbf{Diagnostic forms:} bringe (present subjunctive sg); bringen (present subjunctive pl); brōhte (preterite subjunctive sg); brōhten (preterite subjunctive pl).
% mdparagraph:P177:end

% table: u18-bringan-subjunctive12-SUBJUNCTIVE-DISPLAY
% mdtable:T50:begin
% course-table-style: 2026-10
% course-panel:T50-P01:begin
\par\begingroup
\fontsize{10.5}{13.2}\selectfont
\setlength{\LTleft}{0pt}\setlength{\LTright}{\fill}
\noindent
\begin{tabular}{@{}>{\raggedright\arraybackslash}p{39.459mm}>{\raggedright\arraybackslash}p{39.459mm}>{\raggedright\arraybackslash}p{39.459mm}@{}}
\toprule
\textbf{Person and number} & \textbf{Present} & \textbf{Preterite} \\
\midrule
1 sg. & bringe & brōhte \\
2 sg. & bringe & brōhte \\
3 sg. & bringe & brōhte \\
1–3 pl. & bringen & brōhten \\
\bottomrule
\end{tabular}
\par\endgroup
% course-panel:T50-P01:end
% mdtable:T50:end
\reserve{7\baselineskip}
% mdparagraph:P178:begin
\subsection*{lǣdan: preterite subjunctive}
% mdparagraph:P178:end

% teaching: u18-laedan-preterite-subjunctive6
% mdparagraph:P179:begin
The weak dental past has dd: lǣdde/\allowbreak{}lǣdden. These are the preterite subjunctive singular and plural in all persons, used under a past report of desire.
% mdparagraph:P179:end

\reserve{13\baselineskip}
% mdparagraph:P180:begin
\textbf{Diagnostic forms:} lǣdde (preterite subjunctive sg); lǣdden (preterite subjunctive pl).
% mdparagraph:P180:end

% table: u18-laedan-preterite-subjunctive6-SUBJUNCTIVE-DISPLAY
% mdtable:T51:begin
% course-table-style: 2026-10
% course-panel:T51-P01:begin
\par\begingroup
\fontsize{10.5}{13.2}\selectfont
\setlength{\LTleft}{0pt}\setlength{\LTright}{\fill}
\noindent
\begin{tabular}{@{}>{\raggedright\arraybackslash}p{63.342mm}>{\raggedright\arraybackslash}p{57.846mm}@{}}
\toprule
\textbf{Person and number} & \textbf{Preterite} \\
\midrule
1 sg. & lǣdde \\
2 sg. & lǣdde \\
3 sg. & lǣdde \\
1–3 pl. & lǣdden \\
\bottomrule
\end{tabular}
\par\endgroup
% course-panel:T51-P01:end
% mdtable:T51:end
\reserve{7\baselineskip}
% mdparagraph:P181:begin
\subsection*{stelan: preterite subjunctive}
% mdparagraph:P181:end

% teaching: NEW-U18-VOC-STELAN-U18-SELECTED
\par\begin{minipage}{\linewidth}\raggedright
% mdparagraph:P182:begin
\textbf{stelan} — steal.
% mdparagraph:P182:end

% teaching: u18-stelan-preterite-subjunctive6
% mdparagraph:P183:begin
The preterite subjunctive uses the long-vowel plural stem stǣl- in both numbers: stǣle/\allowbreak{}stǣlen. Its use in the accusation concerns earlier theft.
% mdparagraph:P183:end

% mdparagraph:P184:begin
\textbf{Diagnostic forms:} stǣle (preterite subjunctive sg); stǣlen (preterite subjunctive pl).
% mdparagraph:P184:end

% table: u18-stelan-preterite-subjunctive6-SUBJUNCTIVE-DISPLAY
% mdtable:T52:begin
% course-table-style: 2026-10
% course-panel:T52-P01:begin
\par\begingroup
\fontsize{10.5}{13.2}\selectfont
\setlength{\LTleft}{0pt}\setlength{\LTright}{\fill}
\noindent
\begin{tabular}{@{}>{\raggedright\arraybackslash}p{63.342mm}>{\raggedright\arraybackslash}p{57.846mm}@{}}
\toprule
\textbf{Person and number} & \textbf{Preterite} \\
\midrule
1 sg. & stǣle \\
2 sg. & stǣle \\
3 sg. & stǣle \\
1–3 pl. & stǣlen \\
\bottomrule
\end{tabular}
\par\endgroup
% course-panel:T52-P01:end
% mdtable:T52:end
\end{minipage}\par
\reserve{7\baselineskip}
% mdparagraph:P185:begin
\subsection*{gearcian: singular and plural commands}
% mdparagraph:P185:end

% teaching: u18-gearcian-imperative2
% mdparagraph:P186:begin
Weak class II commands have singular -a and plural -iað; these extend the earlier indicative system.
% mdparagraph:P186:end

\reserve{9\baselineskip}
% mdparagraph:P187:begin
\textbf{Diagnostic forms:} gearca (imperative sg); gearciað (imperative pl).
% mdparagraph:P187:end

% table: u18-gearcian-imperative2-IMPERATIVE-DISPLAY
% mdtable:T53:begin
% course-table-style: 2026-10
% course-panel:T53-P01:begin
\par\begingroup
\fontsize{10.5}{13.2}\selectfont
\setlength{\LTleft}{0pt}\setlength{\LTright}{\fill}
\noindent
\begin{tabular}{@{}>{\raggedright\arraybackslash}p{63.342mm}>{\raggedright\arraybackslash}p{57.846mm}@{}}
\toprule
\textbf{Person and number} & \textbf{Command} \\
\midrule
2 sg. & gearca \\
2 pl. & gearciað \\
\bottomrule
\end{tabular}
\par\endgroup
% course-panel:T53-P01:end
% mdtable:T53:end
\reserve{7\baselineskip}
% mdparagraph:P188:begin
\subsection*{āscian: singular and plural commands}
% mdparagraph:P188:end

% teaching: u18-ascian-imperative2
% mdparagraph:P189:begin
Weak class II commands have singular -a and plural -iað. Retain initial ā.
% mdparagraph:P189:end

\reserve{9\baselineskip}
% mdparagraph:P190:begin
\textbf{Diagnostic forms:} āsca (imperative sg); āsciað (imperative pl).
% mdparagraph:P190:end

% table: u18-ascian-imperative2-IMPERATIVE-DISPLAY
% mdtable:T54:begin
% course-table-style: 2026-10
% course-panel:T54-P01:begin
\par\begingroup
\fontsize{10.5}{13.2}\selectfont
\setlength{\LTleft}{0pt}\setlength{\LTright}{\fill}
\noindent
\begin{tabular}{@{}>{\raggedright\arraybackslash}p{63.342mm}>{\raggedright\arraybackslash}p{57.846mm}@{}}
\toprule
\textbf{Person and number} & \textbf{Command} \\
\midrule
2 sg. & āsca \\
2 pl. & āsciað \\
\bottomrule
\end{tabular}
\par\endgroup
% course-panel:T54-P01:end
% mdtable:T54:end
\reserve{7\baselineskip}
% mdparagraph:P191:begin
\subsection*{diernan: plain and inflected infinitives}
% mdparagraph:P191:end

% teaching: NEW-U18-VOC-DIERNAN-U18-SELECTED
\par\begin{minipage}{\linewidth}\raggedright
% mdparagraph:P192:begin
\textbf{diernan} — conceal; hide.
% mdparagraph:P192:end

% teaching: u18-diernan-infinitive2
% mdparagraph:P193:begin
The pair gives plain and inflected infinitives. Select -anne here; -enne is a documented alternative. A modal takes the plain form; the inflected form still requires a licensed construction.
% mdparagraph:P193:end

% mdparagraph:P194:begin
\textbf{Diagnostic forms:} diernan (plain\_infinitive); diernanne (inflected\_infinitive).
% mdparagraph:P194:end

% table: u18-diernan-infinitive2-BOUNDED-NONFINITE-DISPLAY
% mdtable:T55:begin
\par\begingroup
\fontsize{10.5}{13.2}\selectfont
\setlength{\LTleft}{0pt}\setlength{\LTright}{\fill}
\noindent
\begin{tabular}{@{}>{\raggedright\arraybackslash}p{71.72mm}>{\raggedright\arraybackslash}p{49.46mm}@{}}
\toprule
\textbf{Type} & \textbf{Form} \\
\midrule
Plain infinitive & diernan \\
Inflected infinitive & diernanne \\
\bottomrule
\end{tabular}
\par\endgroup

% mdtable:T55:end
\end{minipage}\par
\reserve{7\baselineskip}
% mdparagraph:P195:begin
\subsection*{ęrian: plain and inflected infinitives}
% mdparagraph:P195:end

% teaching: NEW-U18-VOC-ERIAN-U18-SELECTED
\par\begin{minipage}{\linewidth}\raggedright
% mdparagraph:P196:begin
\textbf{ęrian} — plow.
% mdparagraph:P196:end

% teaching: u18-erian-infinitive2
% mdparagraph:P197:begin
The pair gives plain and inflected infinitives, with selected -ianne. The modal construction takes the plain form. Retain printed ę.
% mdparagraph:P197:end

% mdparagraph:P198:begin
\textbf{Diagnostic forms:} ęrian (plain\_infinitive); ęrianne (inflected\_infinitive).
% mdparagraph:P198:end

% table: u18-erian-infinitive2-BOUNDED-NONFINITE-DISPLAY
% mdtable:T56:begin
\par\begingroup
\fontsize{10.5}{13.2}\selectfont
\setlength{\LTleft}{0pt}\setlength{\LTright}{\fill}
\noindent
\begin{tabular}{@{}>{\raggedright\arraybackslash}p{71.72mm}>{\raggedright\arraybackslash}p{49.46mm}@{}}
\toprule
\textbf{Type} & \textbf{Form} \\
\midrule
Plain infinitive & ęrian \\
Inflected infinitive & ęrianne \\
\bottomrule
\end{tabular}
\par\endgroup

% mdtable:T56:end
\end{minipage}\par
\reserve{7\baselineskip}
% mdparagraph:P199:begin
\subsection*{rīpan: plain and inflected infinitives}
% mdparagraph:P199:end

% teaching: NEW-U18-VOC-RIPAN-U18-SELECTED
\par\begin{minipage}{\linewidth}\raggedright
% mdparagraph:P200:begin
\textbf{rīpan} — reap.
% mdparagraph:P200:end

% teaching: u18-ripan-infinitive2
% mdparagraph:P201:begin
The pair gives plain and inflected infinitives, with selected -anne. The modal construction takes the plain form. Retain ī.
% mdparagraph:P201:end

% mdparagraph:P202:begin
\textbf{Diagnostic forms:} rīpan (plain\_infinitive); rīpanne (inflected\_infinitive).
% mdparagraph:P202:end

% table: u18-ripan-infinitive2-BOUNDED-NONFINITE-DISPLAY
% mdtable:T57:begin
\par\begingroup
\fontsize{10.5}{13.2}\selectfont
\setlength{\LTleft}{0pt}\setlength{\LTright}{\fill}
\noindent
\begin{tabular}{@{}>{\raggedright\arraybackslash}p{71.72mm}>{\raggedright\arraybackslash}p{49.46mm}@{}}
\toprule
\textbf{Type} & \textbf{Form} \\
\midrule
Plain infinitive & rīpan \\
Inflected infinitive & rīpanne \\
\bottomrule
\end{tabular}
\par\endgroup

% mdtable:T57:end
\end{minipage}\par
\reserve{7\baselineskip}
% mdparagraph:P203:begin
\subsection*{giernan: plain and inflected infinitives}
% mdparagraph:P203:end

% teaching: u18-giernan-infinitive2
% mdparagraph:P204:begin
The pair extends the earlier finite system to plain and inflected infinitives, with selected -anne. The genitive and desired-content complements are explained separately.
% mdparagraph:P204:end

\reserve{9\baselineskip}
% mdparagraph:P205:begin
\textbf{Diagnostic forms:} giernan (plain\_infinitive); giernanne (inflected\_infinitive).
% mdparagraph:P205:end

% table: u18-giernan-infinitive2-BOUNDED-NONFINITE-DISPLAY
% mdtable:T58:begin
\par\begingroup
\fontsize{10.5}{13.2}\selectfont
\setlength{\LTleft}{0pt}\setlength{\LTright}{\fill}
\noindent
\begin{tabular}{@{}>{\raggedright\arraybackslash}p{71.72mm}>{\raggedright\arraybackslash}p{49.46mm}@{}}
\toprule
\textbf{Type} & \textbf{Form} \\
\midrule
Plain infinitive & giernan \\
Inflected infinitive & giernanne \\
\bottomrule
\end{tabular}
\par\endgroup

% mdtable:T58:end
\reserve{7\baselineskip}
% mdparagraph:P206:begin
\subsection*{āfǣran: present and past base participles}
% mdparagraph:P206:end

% teaching: NEW-U18-VOC-AFAERAN-U18-SELECTED
\par\begin{minipage}{\linewidth}\raggedright
% mdparagraph:P207:begin
\textbf{āfǣran} — frighten.
% mdparagraph:P207:end

% teaching: u18-afaeran-participle2
% mdparagraph:P208:begin
The long weak stem gives present āfǣrende and uncontracted past āfǣred. The base past participle can remain uninflected in the supplied predicate with plural subject.
% mdparagraph:P208:end

% mdparagraph:P209:begin
\textbf{Diagnostic forms:} āfǣrende (present participle); āfǣred (past participle).
% mdparagraph:P209:end

% table: u18-afaeran-participle2-PARTICIPLES-DISPLAY
% mdtable:T59:begin
\par\begingroup
\fontsize{10.5}{13.2}\selectfont
\setlength{\LTleft}{0pt}\setlength{\LTright}{\fill}
\noindent
\begin{tabular}{@{}>{\raggedright\arraybackslash}p{60.59mm}>{\raggedright\arraybackslash}p{60.59mm}@{}}
\toprule
\textbf{Type} & \textbf{Form} \\
\midrule
present & āfǣrende \\
past & āfǣred \\
\bottomrule
\end{tabular}
\par\endgroup

% mdtable:T59:end
\end{minipage}\par
\reserve{7\baselineskip}
% mdparagraph:P210:begin
\subsection*{bēodan: present and past base participles}
% mdparagraph:P210:end

% teaching: NEW-U18-VOC-BEODAN-U18-COMMAND
\par\begin{minipage}{\linewidth}\raggedright
% mdparagraph:P211:begin
\textbf{bēodan} — command; bid.
% mdparagraph:P211:end

% teaching: u18-beodan-participle2
% mdparagraph:P212:begin
Strong class II gives present bēodende and past boden. The selected new sense is command, distinct from offering. The impersonal command passive remains supplied as a whole.
% mdparagraph:P212:end

% mdparagraph:P213:begin
\textbf{Diagnostic forms:} bēodende (present participle); boden (past participle).
% mdparagraph:P213:end

% table: u18-beodan-participle2-PARTICIPLES-DISPLAY
% mdtable:T60:begin
\par\begingroup
\fontsize{10.5}{13.2}\selectfont
\setlength{\LTleft}{0pt}\setlength{\LTright}{\fill}
\noindent
\begin{tabular}{@{}>{\raggedright\arraybackslash}p{60.59mm}>{\raggedright\arraybackslash}p{60.59mm}@{}}
\toprule
\textbf{Type} & \textbf{Form} \\
\midrule
present & bēodende \\
past & boden \\
\bottomrule
\end{tabular}
\par\endgroup

% mdtable:T60:end
\end{minipage}\par
\reserve{7\baselineskip}
% mdparagraph:P214:begin
\subsection*{berēafian: present and past base participles}
% mdparagraph:P214:end

% teaching: u18-bereafian-participle2
% mdparagraph:P215:begin
Weak class II gives present -iende and past -od. Retain ēa. These are base participles; the deprivation complement is taught separately.
% mdparagraph:P215:end

\reserve{9\baselineskip}
% mdparagraph:P216:begin
\textbf{Diagnostic forms:} berēafiende (present participle); berēafod (past participle).
% mdparagraph:P216:end

% table: u18-bereafian-participle2-PARTICIPLES-DISPLAY
% mdtable:T61:begin
\par\begingroup
\fontsize{10.5}{13.2}\selectfont
\setlength{\LTleft}{0pt}\setlength{\LTright}{\fill}
\noindent
\begin{tabular}{@{}>{\raggedright\arraybackslash}p{57.71mm}>{\raggedright\arraybackslash}p{63.48mm}@{}}
\toprule
\textbf{Type} & \textbf{Form} \\
\midrule
present & berēafiende \\
past & berēafod \\
\bottomrule
\end{tabular}
\par\endgroup

% mdtable:T61:end
\reserve{7\baselineskip}
% mdparagraph:P217:begin
\subsection*{ofslēan: present and past base participles}
% mdparagraph:P217:end

% teaching: NEW-U18-VOC-OFSLEAN-U18-SELECTED
\par\begin{minipage}{\linewidth}\raggedright
% mdparagraph:P218:begin
\textbf{ofslēan} — slay; kill.
% mdparagraph:P218:end

% teaching: u18-ofslean-participle2
% mdparagraph:P219:begin
The contracted verb gives present ofslēande and past ofslægen; ofslagen is a documented alternative. Retain of-. These participles do not introduce a finite system.
% mdparagraph:P219:end

% mdparagraph:P220:begin
\textbf{Diagnostic forms:} ofslēande (present participle); ofslægen (past participle).
% mdparagraph:P220:end

% table: u18-ofslean-participle2-PARTICIPLES-DISPLAY
% mdtable:T62:begin
\par\begingroup
\fontsize{10.5}{13.2}\selectfont
\setlength{\LTleft}{0pt}\setlength{\LTright}{\fill}
\noindent
\begin{tabular}{@{}>{\raggedright\arraybackslash}p{60.59mm}>{\raggedright\arraybackslash}p{60.59mm}@{}}
\toprule
\textbf{Type} & \textbf{Form} \\
\midrule
present & ofslēande \\
past & ofslægen \\
\bottomrule
\end{tabular}
\par\endgroup

% mdtable:T62:end
\end{minipage}\par
\reserve{7\baselineskip}
% mdparagraph:P221:begin
\subsection*{ofðyncean: impersonal third-singular indicative/\allowbreak{}subjunctive, present and preterite}
% mdparagraph:P221:end

% teaching: NEW-U18-VOC-OFTHYNCEAN-U18-SELECTED
% mdparagraph:P222:begin
\textbf{ofðyncean} — grieve; displease (impersonal).
% mdparagraph:P222:end

% teaching: u18-ofthyncean-impersonal4
% mdparagraph:P223:begin
This complete impersonal system is limited to third singular: present indicative ofðyncð, present subjunctive ofðynce, and past ofðūhte in both moods. The experiencer is dative and does not trigger personal or plural agreement. A cause/\allowbreak{}content clause gives what grieves that person.
% mdparagraph:P223:end

% mdparagraph:P224:begin
\textbf{Diagnostic forms:} ofðyncð (present indicative impersonal third singular); ofðynce (present subjunctive sg); ofðūhte (preterite indicative impersonal third singular); ofðūhte (preterite subjunctive sg).
% mdparagraph:P224:end

\reserve{8\baselineskip}
\reserve{6\baselineskip}
% mdparagraph:P225:begin
\textbf{ofðyncean — indicative}
% mdparagraph:P225:end

% table: u18-ofthyncean-impersonal4-INDICATIVE-DISPLAY
% mdtable:T63:begin
% course-table-style: 2026-10
% course-panel:T63-P01:begin
\par\begingroup
\fontsize{10.5}{13.2}\selectfont
\setlength{\LTleft}{0pt}\setlength{\LTright}{\fill}
\noindent
\begin{tabular}{@{}>{\raggedright\arraybackslash}p{39.459mm}>{\raggedright\arraybackslash}p{39.459mm}>{\raggedright\arraybackslash}p{39.459mm}@{}}
\toprule
\textbf{Person and number} & \textbf{Present} & \textbf{Preterite} \\
\midrule
3 sg. & ofðyncð & ofðūhte \\
\bottomrule
\end{tabular}
\par\endgroup
% course-panel:T63-P01:end
% mdtable:T63:end
\reserve{8\baselineskip}
\reserve{6\baselineskip}
% mdparagraph:P226:begin
\textbf{ofðyncean — subjunctive}
% mdparagraph:P226:end

% table: u18-ofthyncean-impersonal4-SUBJUNCTIVE-DISPLAY
% mdtable:T64:begin
% course-table-style: 2026-10
% course-panel:T64-P01:begin
\par\begingroup
\fontsize{10.5}{13.2}\selectfont
\setlength{\LTleft}{0pt}\setlength{\LTright}{\fill}
\noindent
\begin{tabular}{@{}>{\raggedright\arraybackslash}p{39.459mm}>{\raggedright\arraybackslash}p{39.459mm}>{\raggedright\arraybackslash}p{39.459mm}@{}}
\toprule
\textbf{Person and number} & \textbf{Present} & \textbf{Preterite} \\
\midrule
3 sg. & ofðynce & ofðūhte \\
\bottomrule
\end{tabular}
\par\endgroup
% course-panel:T64-P01:end
% mdtable:T64:end
\reserve{7\baselineskip}
% mdparagraph:P227:begin
\section*{4. Meanings and complements}
% mdparagraph:P227:end

% teaching: NEW-U18-VOC-GLAED-U18-GRACIOUS
% mdparagraph:P228:begin
\textbf{glæd} — gracious; favourably disposed toward someone with wið.
% mdparagraph:P228:end

\reserve{7\baselineskip}
% mdparagraph:P229:begin
\subsection*{Need a thing; relative complement}
% mdparagraph:P229:end

% teaching: NEW-U18-GRAM-U18-BETHURFAN-GENITIVE-RELATIVE
% mdparagraph:P230:begin
Beðurfan takes a genitive of the thing needed in the selected construction: hwǣtes, of wheat. In ðā ðing ðe wē beðurfon, the things are accusative outside the relative because they are bought. Inside it, invariant ðe represents the genitive complement required by beðurfan. Its unchanged spelling does not transfer the outer accusative into the relative. Other government is attested; the genitive is the target here.
% mdparagraph:P230:end

\reserve{7\baselineskip}
% mdparagraph:P231:begin
\subsection*{Recognize an accusative person}
% mdparagraph:P231:end

% teaching: NEW-U18-GRAM-U18-ONCNAWAN-OBJECT
% mdparagraph:P232:begin
Oncnāwan in the selected recognition sense takes an accusative recognized person. The nominative recognizer controls the verb; a change of recognizer from ic to wē changes person/\allowbreak{}number independently of the object. Present cnāw-/\allowbreak{}cnǣw- and past cnēow- stay distinct.
% mdparagraph:P232:end

\reserve{7\baselineskip}
% mdparagraph:P233:begin
\subsection*{Open an accusative container}
% mdparagraph:P233:end

% teaching: NEW-U18-GRAM-U18-UNDON-OBJECT
% mdparagraph:P234:begin
Undōn means undo/\allowbreak{}open the tied sacks in 45.8 and takes the accusative container. Present dēð and preterite dyde are different stems; plural preterite undydon follows the plural opener, not the number of sacks.
% mdparagraph:P234:end

\reserve{7\baselineskip}
% mdparagraph:P235:begin
\subsection*{Swear oaths and report their content}
% mdparagraph:P235:end

% teaching: NEW-U18-GRAM-U18-OATH-RECIPIENT-CONTENT
% mdparagraph:P236:begin
Swęrian/\allowbreak{}swerian takes accusative āðas as the sworn utterance. A dative recipient mīnum fæder 48.1 needs no separate tō; ætforan ūs 44.2 names the people in whose presence it was sworn. The reported content is a ðæt clause: mōsten gives past reported permission, brōhte a promised bringing viewed prospectively from the former oath. Keep the recipient distinct from those present at the oath.
% mdparagraph:P236:end

\reserve{7\baselineskip}
% mdparagraph:P237:begin
\subsection*{Desire that a state of affairs occur}
% mdparagraph:P237:end

% teaching: NEW-U18-GRAM-U18-GIERNAN-PROPOSITIONAL-DESIRE
% mdparagraph:P238:begin
After giernan, genitive ðæs points forward to a ðæt clause explaining the desired action. The demonstrative is inflected; the following ðæt is a conjunction. In the past report, lǣdden is first-plural preterite subjunctive, describing the action then desired.
% mdparagraph:P238:end

\reserve{7\baselineskip}
% mdparagraph:P239:begin
\subsection*{Make a person be gracious}
% mdparagraph:P239:end

% teaching: NEW-U18-GRAM-U18-GEDON-OBJECT-PREDICATE
% mdparagraph:P240:begin
Gedōn has accusative person hine plus agreeing accusative masculine singular glædne. In 45.2 gedō is third-singular present subjunctive expressing the wish; it is not an imperative addressed to the named Lord. Wið ēow supplies the graciousness target, as the original hint expressly requires.
% mdparagraph:P240:end

\reserve{7\baselineskip}
% mdparagraph:P241:begin
\subsection*{Give back a person}
% mdparagraph:P241:end

% teaching: NEW-U18-GRAM-U18-AGIEFAN-RESTORATION
% mdparagraph:P242:begin
Āgiefan selects the returned person as accusative and the recipient as dative. Ēow and ēowerne brōðor have different roles. Present subjunctive āgiefe is in the final purpose clause, not merely an unexplained alternative ending.
% mdparagraph:P242:end

\reserve{7\baselineskip}
% mdparagraph:P243:begin
\subsection*{Despoil someone of possessions}
% mdparagraph:P243:end

% teaching: NEW-U18-GRAM-U18-BEREAFIAN-DEPRIVATION
% mdparagraph:P244:begin
Berēafian takes an accusative person and genitive possessions in the selected active construction. In the passive, the deprived person becomes nominative while the possessions remain genitive. In 45.6, assena and ðinga are genitive plurals; selfe agrees with accusative ūs in the following clause. Dative and instrumental alternatives are attested elsewhere, but the local target is genitive.
% mdparagraph:P244:end

\reserve{7\baselineskip}
% mdparagraph:P245:begin
\subsection*{Listen to speech}
% mdparagraph:P245:end

% teaching: NEW-U18-GRAM-U18-HLYSTAN-GENITIVE
% mdparagraph:P246:begin
Hlystan governs genitive speech/\allowbreak{}content in 45.7, as the explicit 156 hint directs. Thus ūrre sprǣce is genitive feminine singular, distinct from accusative ūre sprǣce. The people beseech the steward (ðē accusative); ðū is the subject of subordinate hlyste.
% mdparagraph:P246:end

\reserve{7\baselineskip}
% mdparagraph:P247:begin
\subsection*{Enslave an accusative person}
% mdparagraph:P247:end

% teaching: NEW-U18-GRAM-U18-THEOWIAN-ENSLAVE
% mdparagraph:P248:begin
The selected ðēowian sense is enslave, with an accusative person. The more common serve sense with different government remains separate. The selected ðēowie normalizes the documented -ige/\allowbreak{}-ie variation under Cook 119 and does not change tense or mood.
% mdparagraph:P248:end

\reserve{7\baselineskip}
% mdparagraph:P249:begin
\subsection*{Command a dative person}
% mdparagraph:P249:end

% teaching: NEW-U18-GRAM-U18-BEODAN-COMMAND
% mdparagraph:P250:begin
Bēodan takes a dative addressee in 46.7 and 48.8, as 164 expressly directs. It can accompany a direct command or introduce ðæt plus subjunctive commanded action. 47.11 intentionally retains present bīetst (commandest); the lookalike bīetst of bēatan is a different verb here.
% mdparagraph:P250:end

\reserve{7\baselineskip}
% mdparagraph:P251:begin
\subsection*{Lay out things}
% mdparagraph:P251:end

% teaching: NEW-U18-GRAM-U18-LECGEAN-PLACEMENT
% mdparagraph:P252:begin
Lęcgean takes accusative things laid, with forð marking their presentation in 46.3. Its source reading imperative lęcgeað is separately supplied; learning indicative plural lęcgeað does not by itself assign the command use.
% mdparagraph:P252:end

\reserve{7\baselineskip}
% mdparagraph:P253:begin
\subsection*{Mention a person earlier to someone}
% mdparagraph:P253:end

% teaching: NEW-U18-GRAM-U18-FORESECGEAN-PAST-MENTION
% mdparagraph:P254:begin
Foresęcgean here means mention previously, not foretell the future. The invariant ðe is the mentioned person (accusative function), and mē the dative hearer. The English of is explicitly omitted.
% mdparagraph:P254:end

\reserve{7\baselineskip}
% mdparagraph:P255:begin
\subsection*{Be grieved by a stated circumstance}
% mdparagraph:P255:end

% teaching: NEW-U18-GRAM-U18-OFTHYNCEAN-IMPERSONAL
% mdparagraph:P256:begin
Ofðyncean remains third singular with a dative experiencer and a cause/\allowbreak{}content ðæt clause. Present subjunctive ofðynce means let it not grieve in the negative wish. English it needs no separate Old English pronoun here. A plural experiencer does not make the verb plural.
% mdparagraph:P256:end

\reserve{7\baselineskip}
% mdparagraph:P257:begin
\subsection*{Fear with genitive object and reflexive dative}
% mdparagraph:P257:end

% teaching: NEW-U18-GRAM-U18-ONDRAEDAN-GENITIVE
% mdparagraph:P258:begin
Ondrǣdan can take a reflexive dative and a genitive feared thing. In 48.6, ēow names the people themselves and nānes ðinges the feared entity. In 46.1 the reflexive dative occurs without an expressed feared thing. The unchanged spelling of ēow does not make it accusative. This construction supplements the earlier accusative frame.
% mdparagraph:P258:end

\reserve{7\baselineskip}
% mdparagraph:P259:begin
\subsection*{Requite good with evil}
% mdparagraph:P259:end

% teaching: NEW-U18-GRAM-U18-FORGIELDAN-GOOD-EVIL
% mdparagraph:P260:begin
In 46.9 accusative neuter gōd is the good being repaid; mid yfele is the evil given in return. The recipient frame previously used for paying is not by itself this good-for-evil substitution frame. Wolden is third-plural past subjunctive inside the reported why-question; forgieldan is its plain infinitive.
% mdparagraph:P260:end

\reserve{7\baselineskip}
% mdparagraph:P261:begin
\subsection*{Prospective when with subjunctive}
% mdparagraph:P261:end

% teaching: NEW-U18-GRAM-U18-PROSPECTIVE-WHEN
% mdparagraph:P262:begin
46.9 when you have caught them concerns a future condition for the command to ask. The selected ðonne ðū hīe fō uses second-singular present subjunctive, following 196 b. There is no required present-perfect auxiliary. Original English says you, and the singular steward continues as addressee of Ārīs, gā and āsca.
% mdparagraph:P262:end

\reserve{7\baselineskip}
% mdparagraph:P263:begin
\subsection*{Subjunctive within an optative relative}
% mdparagraph:P263:end

% teaching: NEW-U18-GRAM-U18-RELATIVE-SUBJUNCTIVE-ATTRACTION
% mdparagraph:P264:begin
47.3 retains 195: in sīe hē ofslægen mid ðǣm ðe ðū hit finde, finde is second-singular present subjunctive within the relative governed by the hypothetical/\allowbreak{}optative context. Mid takes a dative associate; ðe is invariant. This is not a rule that every relative clause uses subjunctive.
% mdparagraph:P264:end

\reserve{7\baselineskip}
% mdparagraph:P265:begin
\subsection*{Narrative preterite with anterior meaning}
% mdparagraph:P265:end

% teaching: NEW-U18-GRAM-U18-PRETERITE-ANTERIOR-REVIEW
% mdparagraph:P266:begin
Cook 188 permits the simple preterite for English simple, perfect and pluperfect past. In 45.8 cēapodon and ǣr sealdon give the anterior actions; 46.9 fōron ... sumne dǣl ðæs weges conveys had gone without adding a habban perfect. Each clause remains present in the translation.
% mdparagraph:P266:end

\reserve{7\baselineskip}
% mdparagraph:P267:begin
\subsection*{Conceal oneself}
% mdparagraph:P267:end

% teaching: NEW-U18-GRAM-U18-DIERNAN-REFLEXIVE
% mdparagraph:P268:begin
Diernan takes the concealed person as accusative; hine in 48.4 refers back to nominative Iosep. Meahte is past ability and governs the plain infinitive. Longer is an adverb modifying concealment, not an adjective agreeing with Joseph.
% mdparagraph:P268:end

% mdparagraph:P269:begin
The earlier preterite singular funde of findan remains valid; the alternative fand does not invalidate it. Retain the selected spelling unless using an independently justified alternative.
% mdparagraph:P269:end

\reserve{7\baselineskip}
% mdparagraph:P270:begin
\section*{5. Cumulative analysis and controlled composition}
% mdparagraph:P270:end

% mdparagraph:P271:begin
Use the stated earlier or newly taught forms. Keep any explicitly supplied reference visible. Explain roles from the actual sentence; do not supply an unstated referent or event.
% mdparagraph:P271:end

\reserve{8\baselineskip}
% mdparagraph:P272:begin
\subsection*{A01}
% mdparagraph:P272:end

% item: NEW-U18-A01
% mdparagraph:P273:begin
Analyse the two sentences. No further context identifies the reader.
% mdparagraph:P273:end

% mdparagraph:P274:begin
Se cyning sealde ðǣm męnn ðā bōc. Hē hīe rǣdde.
% mdparagraph:P274:end

% output: NEW-U18-A01-O001
% mdparagraph:P275:begin
\itempara{1}{Identify the giver, recipient and thing given.}
% mdparagraph:P275:end

% output: NEW-U18-A01-O002
% mdparagraph:P276:begin
\itempara{2}{Give the referent and case of hīe.}
% mdparagraph:P276:end

% output: NEW-U18-A01-O003
% mdparagraph:P277:begin
\itempara{3}{Which named people can hē refer to? Explain what the forms establish.}
% mdparagraph:P277:end

% output: NEW-U18-A01-O004
% mdparagraph:P278:begin
\itempara{4}{Rewrite only the second sentence so that the recipient is explicitly the reader.}
% mdparagraph:P278:end

\reserve{8\baselineskip}
% mdparagraph:P279:begin
\subsection*{A02}
% mdparagraph:P279:end

% item: NEW-U18-A02
% mdparagraph:P280:begin
Convert this direct statement into a report. Retain the king as the person who saw the man; use sęcgan with a ðæt clause and a preterite subjunctive in that clause.
% mdparagraph:P280:end

% mdparagraph:P281:begin
Se cyning cwæð: “Ic geseah ðone mann.”
% mdparagraph:P281:end

% mdparagraph:P282:begin
\textbf{Reference:} sęcgan.
% mdparagraph:P282:end

% mdparagraph:P283:begin
\textbf{Reference:} ðæt.
% mdparagraph:P283:end

% output: NEW-U18-A02-O001
% mdparagraph:P284:begin
\itempara{1}{Identify the referent of ic and the time of the seeing relative to the quoted statement.}
% mdparagraph:P284:end

% output: NEW-U18-A02-O002
% mdparagraph:P285:begin
\itempara{2}{Write the complete reported statement.}
% mdparagraph:P285:end

% output: NEW-U18-A02-O003
% mdparagraph:P286:begin
\itempara{3}{Explain whether the chosen subjunctive proves that the king’s statement is false.}
% mdparagraph:P286:end

\reserve{8\baselineskip}
% mdparagraph:P287:begin
\subsection*{A03}
% mdparagraph:P287:end

% item: NEW-U18-A03
% mdparagraph:P288:begin
Analyse the sentence, then replace the men by the man and keep the king singular.
% mdparagraph:P288:end

% mdparagraph:P289:begin
Ðā męnn ðe ðone cyning gesāwon, sungon.
% mdparagraph:P289:end

% output: NEW-U18-A03-O001
% mdparagraph:P290:begin
\itempara{1}{Translate the complete sentence.}
% mdparagraph:P290:end

% output: NEW-U18-A03-O002
% mdparagraph:P291:begin
\itempara{2}{Identify the role represented by ðe and explain the agreement of both finite verbs.}
% mdparagraph:P291:end

% output: NEW-U18-A03-O003
% mdparagraph:P292:begin
\itempara{3}{Write the complete singular version.}
% mdparagraph:P292:end

\reserve{8\baselineskip}
% mdparagraph:P293:begin
\subsection*{C01}
% mdparagraph:P293:end

% item: NEW-U18-C01
% mdparagraph:P294:begin
Translate: The man whom the king saw gave the servant the book. Use ðe and sęllan.
% mdparagraph:P294:end

% mdparagraph:P295:begin
\textbf{Reference:} ðe.
% mdparagraph:P295:end

% mdparagraph:P296:begin
\textbf{Reference:} sęllan.
% mdparagraph:P296:end

% output: NEW-U18-C01-O001
% mdparagraph:P297:begin
\itempara{1}{Write one complete Old English sentence.}
% mdparagraph:P297:end

\reserve{8\baselineskip}
% mdparagraph:P298:begin
\subsection*{C02}
% mdparagraph:P298:end

% item: NEW-U18-C02
% mdparagraph:P299:begin
Translate: The king ordered the men to give the servant the books. Use hātan and a plain infinitive of sęllan.
% mdparagraph:P299:end

% mdparagraph:P300:begin
\textbf{Reference:} hātan.
% mdparagraph:P300:end

% mdparagraph:P301:begin
\textbf{Reference:} sęllan.
% mdparagraph:P301:end

% output: NEW-U18-C02-O001
% mdparagraph:P302:begin
\itempara{1}{Write one complete Old English sentence.}
% mdparagraph:P302:end

\reserve{8\baselineskip}
% mdparagraph:P303:begin
\subsection*{F01}
% mdparagraph:P303:end

% item: NEW-U18-F01
% mdparagraph:P304:begin
Translate each sentence, then explain the selected ending contrast. Use the newly taught words for ass, sack and wheat.
% mdparagraph:P304:end

% output: NEW-U18-F01-O001
% mdparagraph:P305:begin
\itempara{1}{The ass has the sack.}
% mdparagraph:P305:end

% output: NEW-U18-F01-O002
% mdparagraph:P306:begin
\itempara{2}{The asses have the sacks.}
% mdparagraph:P306:end

% output: NEW-U18-F01-O003
% mdparagraph:P307:begin
\itempara{3}{The men give wheat to the asses.}
% mdparagraph:P307:end

% output: NEW-U18-F01-O004
% mdparagraph:P308:begin
\itempara{4}{The wheat is in the sacks.}
% mdparagraph:P308:end

% output: NEW-U18-F01-O005
% mdparagraph:P309:begin
\itempara{5}{Compare the plural noun endings in the first three translations and the location phrase.}
% mdparagraph:P309:end

\reserve{8\baselineskip}
% mdparagraph:P310:begin
\subsection*{F02}
% mdparagraph:P310:end

% item: NEW-U18-F02
% mdparagraph:P311:begin
Translate the two sentences using the newly taught verb meaning need. Then explain the governed noun.
% mdparagraph:P311:end

% output: NEW-U18-F02-O001
% mdparagraph:P312:begin
\itempara{1}{The man needs wheat.}
% mdparagraph:P312:end

% output: NEW-U18-F02-O002
% mdparagraph:P313:begin
\itempara{2}{The men needed wheat.}
% mdparagraph:P313:end

% output: NEW-U18-F02-O003
% mdparagraph:P314:begin
\itempara{3}{Identify the tense and number of each finite verb, and explain the case of wheat in both sentences.}
% mdparagraph:P314:end

\reserve{8\baselineskip}
% mdparagraph:P315:begin
\subsection*{F03}
% mdparagraph:P315:end

% item: NEW-U18-F03
% mdparagraph:P316:begin
Translate the four sentences using the newly taught verbs meaning recognize and undo/\allowbreak{}open. Then explain the roles and number changes.
% mdparagraph:P316:end

% output: NEW-U18-F03-O001
% mdparagraph:P317:begin
\itempara{1}{I recognize the king.}
% mdparagraph:P317:end

% output: NEW-U18-F03-O002
% mdparagraph:P318:begin
\itempara{2}{We recognized the king.}
% mdparagraph:P318:end

% output: NEW-U18-F03-O003
% mdparagraph:P319:begin
\itempara{3}{The man opens the sack.}
% mdparagraph:P319:end

% output: NEW-U18-F03-O004
% mdparagraph:P320:begin
\itempara{4}{The men opened the sacks.}
% mdparagraph:P320:end

% output: NEW-U18-F03-O005
% mdparagraph:P321:begin
\itempara{5}{Identify the direct objects and explain which words change when the subject, object or tense changes.}
% mdparagraph:P321:end

\reserve{7\baselineskip}
% mdparagraph:P322:begin
\section*{6. Cook: the Joseph sequence}
% mdparagraph:P322:end

% teaching: NEW-U18-SOURCE-HINTS
% mdparagraph:P323:begin
Translate every numbered prompt into Old English. Parentheses and brackets are Cook’s original controls: a lexical hint selects a word, om. requests omission of a separate word, and an order or number instruction applies where marked. The numbered grammar references are explained beside their prompts. English wording is retained exactly, including commandest in 47.11. Modern hints and supplied expressions remain available during composition. Preserve every participant, clause and time relation.
% mdparagraph:P323:end

\reserve{7\baselineskip}
% mdparagraph:P324:begin
\section*{Cook 44. The return is proposed}
% mdparagraph:P324:end

\par\begin{minipage}{\linewidth}\raggedright
% mdparagraph:P325:begin
\subsection*{Cook 44.1}
% mdparagraph:P325:end

% item: EX-XII-XIII-44-01
% mdparagraph:P326:begin
Jacob (Iacob) said to his sons, Go and buy us a certain amount of food (męte).
% mdparagraph:P326:end

% mdparagraph:P327:begin
\textbf{Hint:} Iacob.
% mdparagraph:P327:end

% mdparagraph:P328:begin
\textbf{Hint:} męte.
% mdparagraph:P328:end

\end{minipage}\par\addvspace{4pt}

% teaching: U18-SUP-IACOB-NOM
% mdparagraph:P329:begin
\textbf{Supplied:} Iacob — Jacob; nominative, masculine, singular.
% mdparagraph:P329:end

% output: EX-XII-XIII-44-01-O001
\par\begin{minipage}{\linewidth}\raggedright
% mdparagraph:P330:begin
\subsection*{Cook 44.2}
% mdparagraph:P330:end

% item: EX-XII-XIII-44-02
% mdparagraph:P331:begin
Then said Judah (Iudas), The chief of the land swore oaths before us that we should not (mōtan) come before him, except we bring our youngest brother with us.
% mdparagraph:P331:end

% mdparagraph:P332:begin
\textbf{Hint:} Iudas.
% mdparagraph:P332:end

% mdparagraph:P333:begin
\textbf{Hint:} mōtan.
% mdparagraph:P333:end

\end{minipage}\par\addvspace{4pt}

% teaching: U18-SUP-IUDAS-NOM
% mdparagraph:P334:begin
\textbf{Supplied:} Iudas — Judah; nominative, masculine, singular.
% mdparagraph:P334:end

% mdparagraph:P335:begin
\textbf{Supplied:} ætforan — before; in front of.
% mdparagraph:P335:end

% mdparagraph:P336:begin
\textbf{Supplied:} ætforan — before; in front of.
% mdparagraph:P336:end

% teaching: U18-SUP-CONDITIONAL
% mdparagraph:P337:begin
\textbf{Supplied:} būtan — unless; except if.
% mdparagraph:P337:end

% teaching: U18-SUP-GINGSTAN-ACC-M-SG
% mdparagraph:P338:begin
\textbf{Supplied:} gingstan — youngest; accusative, masculine, singular.
% mdparagraph:P338:end

% teaching: NEW-U18-SUP-AETFORAN-DATIVE
% mdparagraph:P339:begin
\textbf{Supplied construction:} \textbf{ætforan ūs} — Before us. Ūs is dative plural after ætforan.
% mdparagraph:P339:end

% mdparagraph:P340:begin
\textbf{Supplied construction:} \textbf{cuman ætforan him} — Come before him. Him is dative singular, referring to the chief. Ætforan retains dative government here even with cuman; copy the whole supplied phrase.
% mdparagraph:P340:end

% output: EX-XII-XIII-44-02-O001
\par\begin{minipage}{\linewidth}\raggedright
% mdparagraph:P341:begin
\subsection*{Cook 44.3}
% mdparagraph:P341:end

% item: EX-XII-XIII-44-03
% mdparagraph:P342:begin
If thou wilt send him with us (verb last), we (om. will) go together and buy the things that (87. d) we need (beðurfan, 131).
% mdparagraph:P342:end

% mdparagraph:P343:begin
\textbf{Hint:} beðurfan.
% mdparagraph:P343:end

% mdparagraph:P344:begin
\textbf{Supplied:} ætgædere — together; at once.
% mdparagraph:P344:end

% mdparagraph:P345:begin
\textbf{(verb last)} Put the finite verb last in this specified clause. This local composition target is not a universal Old English order rule.
% mdparagraph:P345:end

% mdparagraph:P346:begin
\textbf{(om. will)} Omit a separate Old English word for the marked English expression will. Preserve its grammatical relation through the surrounding selected construction.
% mdparagraph:P346:end

% mdparagraph:P347:begin
\textbf{(87. d)} Invariant relative ðe does not display case endings. The head ðā ðing is object of buy, while within the relative ðe supplies the complement required by beðurfan; its spelling does not transfer the head’s accusative to that internal role.
% mdparagraph:P347:end

% mdparagraph:P348:begin
\textbf{(beðurfan, 131)} Beðurfan is a preterite-present: its present singular uses beðearf, plural beðurfon; weak dental past uses beðorfte/\allowbreak{}beðorfton. The complete twelve-cell indicative system is displayed, with the selected genitive frame separately taught.
% mdparagraph:P348:end

% output: EX-XII-XIII-44-03-O001
\end{minipage}\par\addvspace{4pt}

\par\begin{minipage}{\linewidth}\raggedright
% mdparagraph:P349:begin
\subsection*{Cook 44.4}
% mdparagraph:P349:end

% item: EX-XII-XIII-44-04
% mdparagraph:P350:begin
Ye dealt ill (earmlice) with (wið) me, when ye said to him that ye had more (77) brethren (154).
% mdparagraph:P350:end

% mdparagraph:P351:begin
\textbf{Hint:} earmlice.
% mdparagraph:P351:end

% mdparagraph:P352:begin
\textbf{Hint:} wið.
% mdparagraph:P352:end

\end{minipage}\par\addvspace{4pt}

% teaching: NEW-U18-GRAM-U18-H-44-04
% teaching: NEW-U18-VOC-DON-U18-DEAL-WITH
% teaching: U18-SUP-EARMLICE-INVARIANT
% teaching: U18-SUP-TOWARD
% mdparagraph:P353:begin
\textbf{Supplied:} Gē earmlice wið mē dydon, — Complete interpersonal clause: you behaved/\allowbreak{}dealt badly toward me. Wið names the affected counterpart; it is not a physical battle assertion. Dydon agrees with plural gē.
% mdparagraph:P353:end

% mdparagraph:P354:begin
\textbf{Supplied:} earmlice — cruelly; miserably.
% mdparagraph:P354:end

% mdparagraph:P355:begin
\textbf{Supplied:} wið — toward; in relation to a person.
% mdparagraph:P355:end

% mdparagraph:P356:begin
\textbf{Supplied:} dydon — deal; behave toward someone; indicative, preterite, second person, plural.
% mdparagraph:P356:end

% mdparagraph:P357:begin
\textbf{(77)} Mā expresses greater quantity, more; it does not agree as a comparative adjective with brothers. Its complement is the genitive plural brōðra.
% mdparagraph:P357:end

% mdparagraph:P358:begin
\textbf{(154)} Partitive genitive after a quantity: mā brōðra means more brothers. The brothers are the set whose quantity is compared.
% mdparagraph:P358:end

% output: EX-XII-XIII-44-04-O001
\par\begin{minipage}{\linewidth}\raggedright
% mdparagraph:P359:begin
\subsection*{Cook 44.5}
% mdparagraph:P359:end

% item: EX-XII-XIII-44-05
% mdparagraph:P360:begin
The man asked (āscian) us all whether our father was living [lived] and whether we had brethren, and we answered him; we knew not that he would (willan) desire (giernan) that (156. a), that we should bring (lǣdan, 194. b) our brother thither.
% mdparagraph:P360:end

% mdparagraph:P361:begin
\textbf{Hint:} āscian.
% mdparagraph:P361:end

% mdparagraph:P362:begin
\textbf{Hint:} willan.
% mdparagraph:P362:end

% mdparagraph:P363:begin
\textbf{Hint:} giernan.
% mdparagraph:P363:end

% mdparagraph:P364:begin
\textbf{Hint:} lǣdan.
% mdparagraph:P364:end

% mdparagraph:P365:begin
\textbf{[lived]} Use the bracketed English as the explicit target wording, preserving the complete sentence and every participant: lived.
% mdparagraph:P365:end

% mdparagraph:P366:begin
\textbf{(156. a)} Giernan takes genitive ðæs here; the following ðæt clause explains the thing desired. Retain both words: ðæs is the genitive demonstrative, ðæt the content conjunction.
% mdparagraph:P366:end

% mdparagraph:P367:begin
\textbf{(lǣdan, 194. b)} The desired bringing is expressed by past subjunctive lǣdden under the former desire. The brothers are plural subject; the youngest brother is accusative object.
% mdparagraph:P367:end

% output: EX-XII-XIII-44-05-O001
\end{minipage}\par\addvspace{4pt}

\par\begin{minipage}{\linewidth}\raggedright
% mdparagraph:P368:begin
\subsection*{Cook 44.6}
% mdparagraph:P368:end

% item: EX-XII-XIII-44-06
% mdparagraph:P369:begin
Judah said, Send the boy with me, that we may (mugan) buy the things that we need, that we perish (forweorðan) not; except I bring him back home, may I be (193. a) guilty (scyldig).
% mdparagraph:P369:end

% mdparagraph:P370:begin
\textbf{Hint:} mugan.
% mdparagraph:P370:end

% mdparagraph:P371:begin
\textbf{Hint:} forweorðan.
% mdparagraph:P371:end

% mdparagraph:P372:begin
\textbf{Hint:} scyldig.
% mdparagraph:P372:end

% mdparagraph:P373:begin
\textbf{Supplied:} Iudas — Judah; nominative, masculine, singular.
% mdparagraph:P373:end

% mdparagraph:P374:begin
\textbf{Supplied:} būtan — unless; except if.
% mdparagraph:P374:end

% mdparagraph:P375:begin
\textbf{(193. a)} Sīe ic scyldig is an independent present-subjunctive wish or self-imprecation: let me be guilty. It is not a present indicative statement of guilt.
% mdparagraph:P375:end

% output: EX-XII-XIII-44-06-O001
\end{minipage}\par\addvspace{4pt}

\par\begin{minipage}{\linewidth}\raggedright
% mdparagraph:P376:begin
\subsection*{Cook 44.7}
% mdparagraph:P376:end

% item: EX-XII-XIII-44-07
% mdparagraph:P377:begin
Were it not for this delay [if this delay were not], we had [were, 188] come (om. a) second time (sīð, 176).
% mdparagraph:P377:end

% mdparagraph:P378:begin
\textbf{Hint:} sīð.
% mdparagraph:P378:end

\end{minipage}\par\addvspace{4pt}

% teaching: NEW-U18-GRAM-U18-H-44-07
% teaching: NEW-U18-VOC-WESAN-U18-ANTERIOR-COUNTERFACTUAL
% teaching: U18-SUP-INS-SG
% mdparagraph:P379:begin
\textbf{Supplied:} Gif ðēos ęlcung nǣre, wē wǣren ōðre sīðe cumen. — Complete contrary condition and anterior result: if this delay were not, we would already have come a second time. Nǣre is ne+wǣre, third-singular past subjunctive; wǣren is first-plural past subjunctive. Cumen is a past participle with intransitive wesan, not the identically spelled present-subjunctive plural. Ōðre sīðe is adverbial instrumental singular, expressly 176; the sīð hint names the lemma and does not suppress its ending. The entire combined construction and contracted form are supplied; no general counterfactual-perfect mastery is inferred.
% mdparagraph:P379:end

% teaching: NEW-U18-GRAM-U18-HELP-SECOND-TIME
% mdparagraph:P380:begin
\textbf{Supplied:} ōðre sīðe — A second time: both visible instrumental singular forms in the complete hypothetical clause of 44.7. The source headword sīð is not an instruction to leave the noun uninflected.
% mdparagraph:P380:end

% mdparagraph:P381:begin
\textbf{Supplied:} nǣre — be: present and past states/\allowbreak{}identification; subjunctive, preterite, third person, singular.
% mdparagraph:P381:end

% mdparagraph:P382:begin
\textbf{Supplied:} wǣren — be as auxiliary of counterfactual anterior coming; subjunctive, preterite, first person, plural.
% mdparagraph:P382:end

% mdparagraph:P383:begin
\textbf{Supplied:} ōðre — second; instrumental, masculine, singular, strong.
% mdparagraph:P383:end

% mdparagraph:P384:begin
\textbf{Supplied:} sīðe — time; occasion; instrumental, masculine, singular.
% mdparagraph:P384:end

% mdparagraph:P385:begin
\textbf{Supplied:} cumen — come; past, participle.
% mdparagraph:P385:end

% mdparagraph:P386:begin
\textbf{[if this delay were not]} Use the bracketed English as the explicit target wording, preserving the complete sentence and every participant: if this delay were not.
% mdparagraph:P386:end

% mdparagraph:P387:begin
\textbf{[were, 188]} Use the bracketed English as the explicit target wording, preserving the complete sentence and every participant: were, 188. The source deliberately says were, pointing to intransitive wesan plus past participle cumen. In the complete supplied contrary condition, past subjunctive wǣren makes the coming anterior to the imagined result; do not replace it with unassigned habban plus an intransitive participle.
% mdparagraph:P387:end

% mdparagraph:P388:begin
\textbf{(om. a)} Omit a separate Old English word for the marked English expression a. Preserve its grammatical relation through the surrounding selected construction.
% mdparagraph:P388:end

% mdparagraph:P389:begin
\textbf{(sīð, 176)} Section 176 is adverbial instrumental. A second time is ōðre sīðe: ōðre instrumental singular, sīðe the noun’s dative/\allowbreak{}instrumental singular. The sīð hint is a lemma, not a command to suppress the noun ending.
% mdparagraph:P389:end

% output: EX-XII-XIII-44-07-O001
\par\begin{minipage}{\linewidth}\raggedright
% mdparagraph:P390:begin
\subsection*{Cook 44.8}
% mdparagraph:P390:end

% item: EX-XII-XIII-44-08
% mdparagraph:P391:begin
Then said Jacob, If ye needs (nīede, 176) must do so, do as ye will; take of (of) your best fruits in vessels, and carry (bringan) presents (om. to) the man, honey (hunig, 47) and acorns (æcern, 47).
% mdparagraph:P391:end

% mdparagraph:P392:begin
\textbf{Hint:} nīede.
% mdparagraph:P392:end

% mdparagraph:P393:begin
\textbf{Hint:} of.
% mdparagraph:P393:end

% mdparagraph:P394:begin
\textbf{Hint:} bringan.
% mdparagraph:P394:end

% mdparagraph:P395:begin
\textbf{Hint:} hunig.
% mdparagraph:P395:end

% mdparagraph:P396:begin
\textbf{Hint:} æcern.
% mdparagraph:P396:end

% mdparagraph:P397:begin
\textbf{Supplied:} Iacob — Jacob; nominative, masculine, singular.
% mdparagraph:P397:end

\end{minipage}\par\addvspace{4pt}

% teaching: U18-SUP-NIEDE-INVARIANT
% mdparagraph:P398:begin
\textbf{Supplied:} nīede — necessarily; of necessity.
% mdparagraph:P398:end

% mdparagraph:P399:begin
\textbf{(nīede, 176)} Nīede is the supplied adverbial instrumental expressing necessity, necessarily/\allowbreak{}by necessity. With sculon it conveys must needs; do not treat it as a finite form of need.
% mdparagraph:P399:end

% mdparagraph:P400:begin
\textbf{(om. to)} Omit a separate Old English word for the marked English expression to. Preserve its grammatical relation through the surrounding selected construction.
% mdparagraph:P400:end

% mdparagraph:P401:begin
\textbf{(hunig, 47)} The number refers to the noun’s declension. Use the complete hunig table in section 2 and its stated stem rule.
% mdparagraph:P401:end

% mdparagraph:P402:begin
\textbf{(æcern, 47)} The number refers to the noun’s declension. Use the complete æcern table in section 2 and its stated stem rule.
% mdparagraph:P402:end

% output: EX-XII-XIII-44-08-O001
\reserve{7\baselineskip}
% mdparagraph:P403:begin
\section*{Cook 45. The brothers return to Egypt}
% mdparagraph:P403:end

\par\begin{minipage}{\linewidth}\raggedright
% mdparagraph:P404:begin
\subsection*{Cook 45.1}
% mdparagraph:P404:end

% item: EX-XII-XIII-45-01
% mdparagraph:P405:begin
Take your (plur.) brother and go to the man.
% mdparagraph:P405:end

% mdparagraph:P406:begin
\textbf{(plur.)} Address more than one person and select plural personal and finite forms.
% mdparagraph:P406:end

% output: EX-XII-XIII-45-01-O001
\end{minipage}\par\addvspace{4pt}

\par\begin{minipage}{\linewidth}\raggedright
% mdparagraph:P407:begin
\subsection*{Cook 45.2}
% mdparagraph:P407:end

% item: EX-XII-XIII-45-02
% mdparagraph:P408:begin
And my Lord make (gedōn) him gracious (glæd) toward (wið) you, that he may give (āgiefan, 196. f) you your brother, whom he hath with him; I am now bereft [despoiled] of my children (156).
% mdparagraph:P408:end

% mdparagraph:P409:begin
\textbf{Hint:} gedōn.
% mdparagraph:P409:end

% mdparagraph:P410:begin
\textbf{Hint:} glæd.
% mdparagraph:P410:end

% mdparagraph:P411:begin
\textbf{Hint:} wið.
% mdparagraph:P411:end

% mdparagraph:P412:begin
\textbf{Hint:} āgiefan.
% mdparagraph:P412:end

% mdparagraph:P413:begin
\textbf{Supplied:} wið — toward; in relation to a person.
% mdparagraph:P413:end

% mdparagraph:P414:begin
\textbf{(āgiefan, 196. f)} Āgiefe is third-singular present subjunctive in the purpose clause, so that he may restore the brother to you. Ēow is dative recipient and ēowerne brōðor accusative returned person.
% mdparagraph:P414:end

% mdparagraph:P415:begin
\textbf{[despoiled]} Use the bracketed English as the explicit target wording, preserving the complete sentence and every participant: despoiled.
% mdparagraph:P415:end

% mdparagraph:P416:begin
\textbf{(156)} The genitive denotes what someone has been deprived of: mīnra bearna remains genitive after passive eom berēafod. In the active frame the deprived person is accusative; the genitive possessions do not become the passive subject.
% mdparagraph:P416:end

% output: EX-XII-XIII-45-02-O001
\end{minipage}\par\addvspace{4pt}

\par\begin{minipage}{\linewidth}\raggedright
% mdparagraph:P417:begin
\subsection*{Cook 45.3}
% mdparagraph:P417:end

% item: EX-XII-XIII-45-03
% mdparagraph:P418:begin
Then they took the presents and double (twīfeald) money (feoh, 47), and went to the land of the Egyptians to Joseph (Iosep).
% mdparagraph:P418:end

% mdparagraph:P419:begin
\textbf{Hint:} twīfeald.
% mdparagraph:P419:end

% mdparagraph:P420:begin
\textbf{Hint:} feoh.
% mdparagraph:P420:end

% mdparagraph:P421:begin
\textbf{Hint:} Iosep.
% mdparagraph:P421:end

\end{minipage}\par\addvspace{4pt}

% teaching: U18-SUP-EGYPTA-GEN-PL
% mdparagraph:P422:begin
\textbf{Supplied:} Egypta — Egyptians; genitive, masculine, plural.
% mdparagraph:P422:end

% teaching: U18-SUP-IOSEP-DAT
% mdparagraph:P423:begin
\textbf{Supplied:} Iosepe — Joseph; dative, masculine, singular.
% mdparagraph:P423:end

% mdparagraph:P424:begin
\textbf{(feoh, 47)} The number refers to the noun’s declension. Use the complete feoh table in section 2 and its stated stem rule.
% mdparagraph:P424:end

% output: EX-XII-XIII-45-03-O001
\par\begin{minipage}{\linewidth}\raggedright
% mdparagraph:P425:begin
\subsection*{Cook 45.4}
% mdparagraph:P425:end

% item: EX-XII-XIII-45-04
% mdparagraph:P426:begin
When he saw them, then said he to (tō) his steward (gerēfa, 53), Lead in these men, and prepare our feast, that they may (mugan) eat with me at noon (tō middes dæges, 158; 166. 1).
% mdparagraph:P426:end

% mdparagraph:P427:begin
\textbf{Hint:} tō.
% mdparagraph:P427:end

% mdparagraph:P428:begin
\textbf{Hint:} gerēfa.
% mdparagraph:P428:end

% mdparagraph:P429:begin
\textbf{Hint:} mugan.
% mdparagraph:P429:end

% mdparagraph:P430:begin
\textbf{Hint:} tō middes dæges.
% mdparagraph:P430:end

% mdparagraph:P431:begin
\textbf{Supplied:} in — in; inward.
% mdparagraph:P431:end

\end{minipage}\par\addvspace{4pt}

% teaching: U18-SUP-MIDDES-INVARIANT
% mdparagraph:P432:begin
\textbf{Supplied:} middes — midday, in tō middes dæges.
% mdparagraph:P432:end

% mdparagraph:P433:begin
\textbf{(gerēfa, 53)} Gerēfa follows the earlier weak masculine pattern: nominative singular -a; other singular cases and nominative/\allowbreak{}accusative plural -an; genitive plural -ena; dative plural -um.
% mdparagraph:P433:end

% mdparagraph:P434:begin
\textbf{(tō middes dæges, 158; 166. 1)} Tō middes dæges is the source-supplied time phrase at midday. Here tō governs a genitive; middes agrees with genitive dæges. Section 166.1 supplies the parallel principle of adjective agreement within the prepositional phrase, not a direction to change this printed genitive to dative.
% mdparagraph:P434:end

% output: EX-XII-XIII-45-04-O001
\par\begin{minipage}{\linewidth}\raggedright
% mdparagraph:P435:begin
\subsection*{Cook 45.5}
% mdparagraph:P435:end

% item: EX-XII-XIII-45-05
% mdparagraph:P436:begin
And he did as he was commanded [was commanded him; cf. p. 177, l. 10].
% mdparagraph:P436:end

\end{minipage}\par\addvspace{4pt}

% teaching: NEW-U18-GRAM-U18-H-45-05
% mdparagraph:P437:begin
\textbf{Supplied:} hē dyde swā him boden wæs. — Impersonal command passive: he did as was commanded him. Him is dative addressee; boden wæs is singular without promoting him to nominative. Cook 177.10 swā hiere beboden wæs is the explicitly cited parallel; replacing be-bēodan with independently attested command bēodan is a declared lexical adaptation.
% mdparagraph:P437:end

% mdparagraph:P438:begin
\textbf{[was commanded him; cf. p. 177, l. 10]} Use the bracketed English as the explicit target wording, preserving the complete sentence and every participant: was commanded him; cf. p. 177, l. 10. Cook 177.10 has swā hiere beboden wæs. The supplied expression uses bēodan: swā him boden wæs, as was commanded him. Him is a dative addressee in an impersonal passive, not a nominative subject.
% mdparagraph:P438:end

% output: EX-XII-XIII-45-05-O001
\par\begin{minipage}{\linewidth}\raggedright
% mdparagraph:P439:begin
\subsection*{Cook 45.6}
% mdparagraph:P439:end

% item: EX-XII-XIII-45-06
% mdparagraph:P440:begin
When they were there (ðǣr), then they became greatly frightened, and said to one another (betwēonan him), Because of (for) the money that we found in our sacks (sacc, 43), we (om. have) been led (62) in hither, that he may despoil us of our asses, and of our things, and enslave us (our)selves.
% mdparagraph:P440:end

% mdparagraph:P441:begin
\textbf{Hint:} ðǣr.
% mdparagraph:P441:end

% mdparagraph:P442:begin
\textbf{Hint:} betwēonan him.
% mdparagraph:P442:end

% mdparagraph:P443:begin
\textbf{Hint:} for.
% mdparagraph:P443:end

% mdparagraph:P444:begin
\textbf{Hint:} sacc.
% mdparagraph:P444:end

\end{minipage}\par\addvspace{4pt}

% teaching: NEW-U18-GRAM-U18-H-45-06-PASSIVE
% mdparagraph:P445:begin
\textbf{Supplied:} wē sindon hider in lǣded, — We have been led in hither. Present wesan plus the past participle is the selected passive/\allowbreak{}result expression; uninflected lǣded with plural wē is expressly allowed by 62. Inflected plural lǣdde is not needed to satisfy the prompt and is not silently treated as taught.
% mdparagraph:P445:end

% teaching: NEW-U18-GRAM-U18-HELP-FOR-CAUSE
% teaching: U18-SUP-CAUSE
% mdparagraph:P446:begin
\textbf{Supplied:} For ðǣm fēo — Because of that money /\allowbreak{} on account of that money. Ðǣm and fēo are dative neuter singular.
% mdparagraph:P446:end

% teaching: U18-SUP-BETWEONAN-INVARIANT
% mdparagraph:P447:begin
\textbf{Supplied:} betwēonan — between; toward one another.
% mdparagraph:P447:end

% mdparagraph:P448:begin
\textbf{Supplied:} For — because of; on account of; governs dative.
% mdparagraph:P448:end

% mdparagraph:P449:begin
\textbf{Supplied:} in — in; inward.
% mdparagraph:P449:end

% mdparagraph:P450:begin
\textbf{(sacc, 43)} The number refers to the noun’s declension. Use the complete sacc table in section 2 and its stated stem rule.
% mdparagraph:P450:end

% mdparagraph:P451:begin
\textbf{(om. have)} Omit a separate Old English word for the marked English expression have. Preserve its grammatical relation through the surrounding selected construction.
% mdparagraph:P451:end

% mdparagraph:P452:begin
\textbf{(62)} A past participle used predicatively may be uninflected. Thus plural wē can have base lǣded in wē sindon hider in lǣded. No new plural participial ending is required by the prompt.
% mdparagraph:P452:end

% mdparagraph:P453:begin
\textbf{(our)} The parenthesis (our) belongs to the English expression (our)selves, reinforcing the preceding us. In the supplied clause, ūs selfe means ourselves: selfe agrees with accusative plural ūs. This parenthesis does not introduce a possessive modifying the things.
% mdparagraph:P453:end

% output: EX-XII-XIII-45-06-O001
\par\begin{minipage}{\linewidth}\raggedright
% mdparagraph:P454:begin
\subsection*{Cook 45.7}
% mdparagraph:P454:end

% item: EX-XII-XIII-45-07
% mdparagraph:P455:begin
Then said they to (tō) the steward, O sir, we pray (biddan) thee that thou hear (hlystan, 156) our speech; we were here a while ago and bought wheat (hwǣte, 44) for ourselves (161; 184. a).
% mdparagraph:P455:end

% mdparagraph:P456:begin
\textbf{Hint:} tō.
% mdparagraph:P456:end

% mdparagraph:P457:begin
\textbf{Hint:} biddan.
% mdparagraph:P457:end

% mdparagraph:P458:begin
\textbf{Hint:} hlystan.
% mdparagraph:P458:end

% mdparagraph:P459:begin
\textbf{Hint:} hwǣte.
% mdparagraph:P459:end

% mdparagraph:P460:begin
\textbf{(hlystan, 156)} Hlystan governs genitive speech/\allowbreak{}content, so ūrre sprǣce is genitive feminine singular. Hlyste is second-singular present subjunctive after the plea.
% mdparagraph:P460:end

% mdparagraph:P461:begin
\textbf{(hwǣte, 44)} The number refers to the noun’s declension. Use the complete hwǣte table in section 2 and its stated stem rule.
% mdparagraph:P461:end

% mdparagraph:P462:begin
\textbf{(161; 184. a)} Ūs is the benefactive/\allowbreak{}reflexive dative: we bought wheat for ourselves. It refers to the buying subject, while wheat is the accusative thing purchased. Retain the phrase; the prompt gives no omission instruction.
% mdparagraph:P462:end

\end{minipage}\par\addvspace{4pt}

% teaching: NEW-U18-SUP-LEOF-ADDRESS
% mdparagraph:P463:begin
\textbf{Supplied construction:} \textbf{Lēof} — Sir. Lēof is a substantival use of the strong masculine nominative singular adjective in direct address. This respectful address sense is supplied here; the previously taught meaning dear alone does not establish it.
% mdparagraph:P463:end

% teaching: NEW-U18-SUP-HWILUM-AER
% mdparagraph:P464:begin
\textbf{Supplied construction:} \textbf{hwīlum ǣr} — A while ago; formerly. Use hwīlum ǣr as a complete earlier-time expression. Hwīlum here is part of this supplied expression, rather than the frequency meaning sometimes previously learned.
% mdparagraph:P464:end

% output: EX-XII-XIII-45-07-O001
\par\begin{minipage}{\linewidth}\raggedright
% mdparagraph:P465:begin
\subsection*{Cook 45.8}
% mdparagraph:P465:end

% item: EX-XII-XIII-45-08
% mdparagraph:P466:begin
When we had trafficked (cēapian) and were on the way home (hāmweard wǣron), then we undid (undōn) our sacks; then we found the money therein which we had given (see ǣr, adv.).
% mdparagraph:P466:end

% mdparagraph:P467:begin
\textbf{Hint:} cēapian.
% mdparagraph:P467:end

% mdparagraph:P468:begin
\textbf{Hint:} hāmweard wǣron.
% mdparagraph:P468:end

% mdparagraph:P469:begin
\textbf{Hint:} undōn.
% mdparagraph:P469:end

% mdparagraph:P470:begin
\textbf{Hint:} ǣr.
% mdparagraph:P470:end

\end{minipage}\par\addvspace{4pt}

% teaching: U18-SUP-HAMWEARD-INVARIANT
% mdparagraph:P471:begin
\textbf{Supplied:} hāmweard — homeward.
% mdparagraph:P471:end

% teaching: U18-SUP-THAERON-INVARIANT
% mdparagraph:P472:begin
\textbf{Supplied:} ðǣron — therein; in it.
% mdparagraph:P472:end

% output: EX-XII-XIII-45-08-O001
\par\begin{minipage}{\linewidth}\raggedright
% mdparagraph:P473:begin
\subsection*{Cook 45.9}
% mdparagraph:P473:end

% item: EX-XII-XIII-45-09
% mdparagraph:P474:begin
Now we have brought it back, and also other silver to traffic with [with to traffic]; and we know not who put (dōn) it into (on, 172) our sacks.
% mdparagraph:P474:end

% mdparagraph:P475:begin
\textbf{Hint:} dōn.
% mdparagraph:P475:end

% mdparagraph:P476:begin
\textbf{Hint:} on.
% mdparagraph:P476:end

\end{minipage}\par\addvspace{4pt}

% teaching: NEW-U18-GRAM-U18-H-45-09
% teaching: U18-SUP-INSTRUMENT-STRANDED
% mdparagraph:P477:begin
\textbf{Supplied:} and ēac ōðer seolfor mid tō cēapianne; — The earlier we brought governs the coordinated silver-object; mid marks silver as the instrument of trading and is stranded before the purposive inflected infinitive. This does not require a new finite verb or make the silver a human trading partner. Supply the whole compressed phrase.
% mdparagraph:P477:end

% mdparagraph:P478:begin
\textbf{Supplied:} mid — with; by means of.
% mdparagraph:P478:end

% mdparagraph:P479:begin
\textbf{[with to traffic]} Use the bracketed English as the explicit target wording, preserving the complete sentence and every participant: with to traffic.
% mdparagraph:P479:end

% mdparagraph:P480:begin
\textbf{(on, 172)} On plus accusative marks putting something into a container. Ūre saccas is accusative plural after on in the clause about who put the money there; static location would instead take a dative.
% mdparagraph:P480:end

% output: EX-XII-XIII-45-09-O001
\reserve{7\baselineskip}
% mdparagraph:P481:begin
\section*{Cook 46. The welcome and the cup}
% mdparagraph:P481:end

\par\begin{minipage}{\linewidth}\raggedright
% mdparagraph:P482:begin
\subsection*{Cook 46.1}
% mdparagraph:P482:end

% item: EX-XII-XIII-46-01
% mdparagraph:P483:begin
Then said the steward, Peace be with you, fear (ondrǣdan, 184) ye not; your God and your father’s God (om. hath) given you treasure in your sacks.
% mdparagraph:P483:end

% mdparagraph:P484:begin
\textbf{Hint:} ondrǣdan.
% mdparagraph:P484:end

\end{minipage}\par\addvspace{4pt}

% teaching: NEW-U18-GRAM-U18-H-46-01
% mdparagraph:P485:begin
\textbf{Supplied:} ēower God and ēowres fæder God sealde ēow goldhord — The one deity is described as your God and your father’s God, explaining singular sealde. Ēowres agrees with the genitive father; second ēow is dative recipient. The expression does not assert two gods.
% mdparagraph:P485:end

% mdparagraph:P486:begin
\textbf{(ondrǣdan, 184)} Ondrǣdað ēow uses a reflexive dative: do not fear. The plural addressees are subject, and ēow refers back to them. It is not the object feared; no feared object is expressed in this item.
% mdparagraph:P486:end

% mdparagraph:P487:begin
\textbf{(om. hath)} Omit a separate Old English word for the marked English expression hath. Preserve its grammatical relation through the surrounding selected construction.
% mdparagraph:P487:end

% output: EX-XII-XIII-46-01-O001
\par\begin{minipage}{\linewidth}\raggedright
% mdparagraph:P488:begin
\subsection*{Cook 46.2}
% mdparagraph:P488:end

% item: EX-XII-XIII-46-02
% mdparagraph:P489:begin
And they washed (101) their feet, and he gave their asses provender [fodder].
% mdparagraph:P489:end

% mdparagraph:P490:begin
\textbf{(101)} Ðwēan is a contracted strong verb. Select the earlier past plural ðwōgon for they washed; the headword’s contracted vowel does not supply the past stem by itself.
% mdparagraph:P490:end

% mdparagraph:P491:begin
\textbf{[fodder]} Use the bracketed English as the explicit target wording, preserving the complete sentence and every participant: fodder.
% mdparagraph:P491:end

% output: EX-XII-XIII-46-02-O001
\end{minipage}\par\addvspace{4pt}

\par\begin{minipage}{\linewidth}\raggedright
% mdparagraph:P492:begin
\subsection*{Cook 46.3}
% mdparagraph:P492:end

% item: EX-XII-XIII-46-03
% mdparagraph:P493:begin
They set (lęcgean) forth their presents, against that (ongēan ðætte) Joseph (om. should) enter.
% mdparagraph:P493:end

% mdparagraph:P494:begin
\textbf{Hint:} lęcgean.
% mdparagraph:P494:end

% mdparagraph:P495:begin
\textbf{Hint:} ongēan ðætte.
% mdparagraph:P495:end

\end{minipage}\par\addvspace{4pt}

% teaching: NEW-U18-GRAM-U18-H-46-03
% teaching: U18-SUP-IOSEP-NOM
% teaching: U18-SUP-TEMPORAL-CONJUNCTION-COMPONENT
% mdparagraph:P496:begin
\textbf{Supplied:} ongēan ðætte Iosep inēode. — Against/\allowbreak{}in readiness for Joseph’s entering. The gifts are laid out in preparation for this prospective arrival; inēode is third-singular past subjunctive, an already productive ingān cell. The whole source conjunction is supplied, not equated with spatial against.
% mdparagraph:P496:end

% mdparagraph:P497:begin
\textbf{Supplied:} ongēan — by the time that.
% mdparagraph:P497:end

% mdparagraph:P498:begin
\textbf{Supplied:} ðætte — that (content conjunction).
% mdparagraph:P498:end

% mdparagraph:P499:begin
\textbf{Supplied:} Iosep — Joseph; nominative, masculine, singular.
% mdparagraph:P499:end

% mdparagraph:P500:begin
\textbf{(om. should)} Omit a separate Old English word for the marked English expression should. Preserve its grammatical relation through the surrounding selected construction.
% mdparagraph:P500:end

% output: EX-XII-XIII-46-03-O001
\par\begin{minipage}{\linewidth}\raggedright
% mdparagraph:P501:begin
\subsection*{Cook 46.4}
% mdparagraph:P501:end

% item: EX-XII-XIII-46-04
% mdparagraph:P502:begin
Then Joseph recognized them, and asked them whether their father was hale, or whether he was living [lived].
% mdparagraph:P502:end

% mdparagraph:P503:begin
\textbf{Supplied:} Iosep — Joseph; nominative, masculine, singular.
% mdparagraph:P503:end

% mdparagraph:P504:begin
\textbf{[lived]} Use the bracketed English as the explicit target wording, preserving the complete sentence and every participant: lived.
% mdparagraph:P504:end

% output: EX-XII-XIII-46-04-O001
\end{minipage}\par\addvspace{4pt}

\par\begin{minipage}{\linewidth}\raggedright
% mdparagraph:P505:begin
\subsection*{Cook 46.5}
% mdparagraph:P505:end

% item: EX-XII-XIII-46-05
% mdparagraph:P506:begin
Then said they, Thy servant (ðēow) our father is well (gesund); he is yet alive [liveth yet].
% mdparagraph:P506:end

% mdparagraph:P507:begin
\textbf{Hint:} ðēow.
% mdparagraph:P507:end

% mdparagraph:P508:begin
\textbf{Hint:} gesund.
% mdparagraph:P508:end

% mdparagraph:P509:begin
\textbf{[liveth yet]} Use the bracketed English as the explicit target wording, preserving the complete sentence and every participant: liveth yet.
% mdparagraph:P509:end

% output: EX-XII-XIII-46-05-O001
\end{minipage}\par\addvspace{4pt}

\par\begin{minipage}{\linewidth}\raggedright
% mdparagraph:P510:begin
\subsection*{Cook 46.6}
% mdparagraph:P510:end

% item: EX-XII-XIII-46-06
% mdparagraph:P511:begin
When Joseph saw his brother Benjamin, then said he, Is this (neut.) the lad, (om. of) whom (87. d) ye told me before (foresęcgean)? And he said, God have compassion on thee, my son.
% mdparagraph:P511:end

% mdparagraph:P512:begin
\textbf{Hint:} foresęcgean.
% mdparagraph:P512:end

% mdparagraph:P513:begin
\textbf{Supplied:} Iosep — Joseph; nominative, masculine, singular.
% mdparagraph:P513:end

\end{minipage}\par\addvspace{4pt}

% teaching: U18-SUP-BENIAMIN-ACC
% mdparagraph:P514:begin
\textbf{Supplied:} Beniamin — Benjamin; accusative, masculine, singular.
% mdparagraph:P514:end

% mdparagraph:P515:begin
\textbf{(neut.)} Use a neuter demonstrative for identifying this, even though the identified boy is a person; the complete identifying clause remains visible in its support.
% mdparagraph:P515:end

% mdparagraph:P516:begin
\textbf{(om. of)} Omit a separate Old English word for the marked English expression of. Preserve its grammatical relation through the surrounding selected construction.
% mdparagraph:P516:end

% mdparagraph:P517:begin
\textbf{(87. d)} The relative pronoun ðe is invariant. It represents the boy as the person previously mentioned; mē is the dative person told. Omit English of as explicitly instructed instead of adding an unsupported preposition.
% mdparagraph:P517:end

% output: EX-XII-XIII-46-06-O001
\par\begin{minipage}{\linewidth}\raggedright
% mdparagraph:P518:begin
\subsection*{Cook 46.7}
% mdparagraph:P518:end

% item: EX-XII-XIII-46-07
% mdparagraph:P519:begin
Then Joseph commanded (bēodan, 164) his steward, and said, Fill their sacks with wheat, and put the money of each of them into (on, 172) his own sack.
% mdparagraph:P519:end

% mdparagraph:P520:begin
\textbf{Hint:} bēodan.
% mdparagraph:P520:end

% mdparagraph:P521:begin
\textbf{Hint:} on.
% mdparagraph:P521:end

% mdparagraph:P522:begin
\textbf{Supplied:} Iosep — Joseph; nominative, masculine, singular.
% mdparagraph:P522:end

% mdparagraph:P523:begin
\textbf{(bēodan, 164)} The person commanded by bēodan is dative: his gerēfan. The direct imperative that follows gives the commanded act. This is the newly assigned command sense, distinct from the earlier offer sense.
% mdparagraph:P523:end

% mdparagraph:P524:begin
\textbf{(on, 172)} On plus accusative gives the destination into his own sack: on his āgenne sacc. The adjective and noun agree as accusative masculine singular.
% mdparagraph:P524:end

\end{minipage}\par\addvspace{4pt}

% teaching: NEW-U18-SUP-GEFYLLAN-MATERIAL
% mdparagraph:P525:begin
\textbf{Supplied construction:} \textbf{Gefyl hiera saccas mid hwǣte} — Fill their sacks with wheat. Gefyl is a singular command; hiera saccas is the accusative plural container, and mid hwǣte names the material with dative singular hwǣte. This full material-complement frame is supplied here.
% mdparagraph:P525:end

% output: EX-XII-XIII-46-07-O001
\par\begin{minipage}{\linewidth}\raggedright
% mdparagraph:P526:begin
\subsection*{Cook 46.8}
% mdparagraph:P526:end

% item: EX-XII-XIII-46-08
% mdparagraph:P527:begin
And take my silver cup (læfel, 43), and the price (weorð, 47) of the wheat, which he gave thee, and put into the sack of the youngest; and he did so.
% mdparagraph:P527:end

% mdparagraph:P528:begin
\textbf{Hint:} læfel.
% mdparagraph:P528:end

% mdparagraph:P529:begin
\textbf{Hint:} weorð.
% mdparagraph:P529:end

\end{minipage}\par\addvspace{4pt}

% teaching: NEW-U18-GRAM-U18-H-46-08
% teaching: U18-SUP-GINGSTAN-GEN-M-SG
% teaching: U18-SUP-SEOLFRENNE-ACC-M-SG
% mdparagraph:P530:begin
\textbf{Supplied:} nim mīnne seolfrenne læfel, and ðæs hwǣtes weorð ðe hē ðē sealde, and dō on ðæs gingstan sacc; — Supplied seolfrenne is the strong accusative masculine singular adjective agreeing with the cup; ðæs hwǣtes specifies the wheat price, and the price is what he gave thee. The coordinated put shares the cup-and-price objects; no unprinted singular it is required. The youngest is the substantival weak adjective identifying the sack’s possessor. The exact silver adjective is supplied; its remaining forms are not assigned here.
% mdparagraph:P530:end

% teaching: NEW-U18-GRAM-U18-HELP-YOUNGEST
% mdparagraph:P531:begin
\textbf{Supplied:} ðæs gingstan — The youngest and eldest refer to individual brothers. Supply the masculine singular human referents and each case, rather than inferring substantive person uses from familiar adjective headwords.
% mdparagraph:P531:end

% teaching: NEW-U18-GRAM-U18-HELP-SILVER-CUP
% mdparagraph:P532:begin
\textbf{Supplied:} mīnne seolfrenne læfel, — My silver cup; the entire strong accusative adjective is supplied. Earlier proposed læfel seolfres is withdrawn because a direct material-genitive frame was not established by the checked evidence.
% mdparagraph:P532:end

% mdparagraph:P533:begin
\textbf{Supplied:} seolfrenne — silver; made of silver; accusative, masculine, singular.
% mdparagraph:P533:end

% mdparagraph:P534:begin
\textbf{Supplied:} gingstan — youngest; genitive, masculine, singular.
% mdparagraph:P534:end

% mdparagraph:P535:begin
\textbf{(læfel, 43)} The number refers to the noun’s declension. Use the complete læfel table in section 2 and its stated stem rule.
% mdparagraph:P535:end

% mdparagraph:P536:begin
\textbf{(weorð, 47)} The number refers to the noun’s declension. Use the complete weorð table in section 2 and its stated stem rule.
% mdparagraph:P536:end

% output: EX-XII-XIII-46-08-O001
\par\begin{minipage}{\linewidth}\raggedright
% mdparagraph:P537:begin
\subsection*{Cook 46.9}
% mdparagraph:P537:end

% item: EX-XII-XIII-46-09
% mdparagraph:P538:begin
In (om. the) morning, as (ðā) they journeyed, and they were without (om. the) city, and had gone (faran) some part of the way, then said Joseph to (tō) his steward, Arise and go after these men, and when you have (196. b) overtaken [caught] them, then ask them why they would (willan) requite good with (mid) evil.
% mdparagraph:P538:end

% mdparagraph:P539:begin
\textbf{Hint:} ðā.
% mdparagraph:P539:end

% mdparagraph:P540:begin
\textbf{Hint:} faran.
% mdparagraph:P540:end

% mdparagraph:P541:begin
\textbf{Hint:} tō.
% mdparagraph:P541:end

% mdparagraph:P542:begin
\textbf{Hint:} willan.
% mdparagraph:P542:end

% mdparagraph:P543:begin
\textbf{Hint:} mid.
% mdparagraph:P543:end

\end{minipage}\par\addvspace{4pt}

% teaching: NEW-U18-GRAM-U18-HELP-OUTSIDE-CITY
% teaching: U18-SUP-OUTSIDE
% mdparagraph:P544:begin
\textbf{Supplied:} būtan byrig — Outside the city. Byrig is dative singular. The preposition here gives a location, rather than the separate without/\allowbreak{}absence or unless uses.
% mdparagraph:P544:end

% teaching: U18-SUP-MORNING
% mdparagraph:P545:begin
\textbf{Supplied:} On — in; at a time.
% mdparagraph:P545:end

% mdparagraph:P546:begin
\textbf{Supplied:} būtan — outside; governs dative.
% mdparagraph:P546:end

% mdparagraph:P547:begin
\textbf{Supplied:} Iosep — Joseph; nominative, masculine, singular.
% mdparagraph:P547:end

% teaching: U18-SUP-ABSTRACT-GOOD-ACC-SG
% mdparagraph:P548:begin
\textbf{Supplied:} gōd — good; benefit; kindness; accusative, neuter, singular.
% mdparagraph:P548:end

% mdparagraph:P549:begin
\textbf{Supplied:} yfele — evil; dative, neuter, singular.
% mdparagraph:P549:end

% mdparagraph:P550:begin
\textbf{(om. the)} Omit a separate Old English word for the marked English expression the. Preserve its grammatical relation through the surrounding selected construction.
% mdparagraph:P550:end

% mdparagraph:P551:begin
\textbf{(om. the)} Omit a separate Old English word for the marked English expression the. Preserve its grammatical relation through the surrounding selected construction.
% mdparagraph:P551:end

% mdparagraph:P552:begin
\textbf{(196. b)} The catching is prospective from the command: when you have caught them, then ask. Cook 196 b permits present subjunctive fō in that future temporal condition; the source’s you is singular here because Joseph addresses one steward. An English perfect does not force an OE perfect auxiliary.
% mdparagraph:P552:end

% mdparagraph:P553:begin
\textbf{[caught]} Use the bracketed English as the explicit target wording, preserving the complete sentence and every participant: caught.
% mdparagraph:P553:end

% teaching: NEW-U18-SUP-ON-MORGEN
% mdparagraph:P554:begin
\textbf{Supplied construction:} \textbf{On morgen} — In the morning. Morgen is accusative singular in this fixed temporal expression. Keep the whole phrase available; it does not introduce a general rule for every time expression with on.
% mdparagraph:P554:end

% output: EX-XII-XIII-46-09-O001
\reserve{7\baselineskip}
% mdparagraph:P555:begin
\section*{Cook 47. The accusation and the plea}
% mdparagraph:P555:end

\par\begin{minipage}{\linewidth}\raggedright
% mdparagraph:P556:begin
\subsection*{Cook 47.1}
% mdparagraph:P556:end

% item: EX-XII-XIII-47-01
% mdparagraph:P557:begin
The cup, which ye stole (forstelan, 105), was very precious to my lord (hlāford).
% mdparagraph:P557:end

% mdparagraph:P558:begin
\textbf{Hint:} forstelan.
% mdparagraph:P558:end

% mdparagraph:P559:begin
\textbf{Hint:} hlāford.
% mdparagraph:P559:end

% mdparagraph:P560:begin
\textbf{(forstelan, 105)} Forstelan belongs to strongIV: present e/\allowbreak{}i, past singular æ, past plural ǣ. Second-plural forstǣlon therefore differs from third-singular forstæl in 47.8. The complete twelve-cell indicative system is displayed.
% mdparagraph:P560:end

% output: EX-XII-XIII-47-01-O001
\end{minipage}\par\addvspace{4pt}

\par\begin{minipage}{\linewidth}\raggedright
% mdparagraph:P561:begin
\subsection*{Cook 47.2}
% mdparagraph:P561:end

% item: EX-XII-XIII-47-02
% mdparagraph:P562:begin
He did as was commanded him, and when they were overtaken [caught], they said, The money that we found in our sacks we brought to thee from our country; thinkest (wēnan) thou that we (om. have) stolen (stelan) thy lord’s gold or his silver?
% mdparagraph:P562:end

% mdparagraph:P563:begin
\textbf{Hint:} wēnan.
% mdparagraph:P563:end

% mdparagraph:P564:begin
\textbf{Hint:} stelan.
% mdparagraph:P564:end

% mdparagraph:P565:begin
\textbf{Supplied:} Hē dyde swā him boden wæs, — Impersonal command passive: he did as was commanded him. Him is dative addressee; boden wæs is singular without promoting him to nominative. Cook 177.10 swā hiere beboden wæs is the explicitly cited parallel; replacing be-bēodan with independently attested command bēodan is a declared lexical adaptation.
% mdparagraph:P565:end

% mdparagraph:P566:begin
\textbf{[caught]} Use the bracketed English as the explicit target wording, preserving the complete sentence and every participant: caught.
% mdparagraph:P566:end

% mdparagraph:P567:begin
\textbf{(om. have)} Omit a separate Old English word for the marked English expression have. Preserve its grammatical relation through the surrounding selected construction.
% mdparagraph:P567:end

% output: EX-XII-XIII-47-02-O001
\end{minipage}\par\addvspace{4pt}

\par\begin{minipage}{\linewidth}\raggedright
% mdparagraph:P568:begin
\subsection*{Cook 47.3}
% mdparagraph:P568:end

% item: EX-XII-XIII-47-03
% mdparagraph:P569:begin
Look for [seek out] it; and let him be slain (ofslēan) with whom thou findest (195) it.
% mdparagraph:P569:end

% mdparagraph:P570:begin
\textbf{Hint:} ofslēan.
% mdparagraph:P570:end

% mdparagraph:P571:begin
\textbf{[seek out]} Use the bracketed English as the explicit target wording, preserving the complete sentence and every participant: seek out.
% mdparagraph:P571:end

% mdparagraph:P572:begin
\textbf{(195)} Finde is second-singular present subjunctive in the relative under the optative/\allowbreak{}hypothetical sīe hē ofslægen. Section 195 describes this subjunctive attraction; it does not require subjunctive in every relative clause.
% mdparagraph:P572:end

% output: EX-XII-XIII-47-03-O001
\end{minipage}\par\addvspace{4pt}

\par\begin{minipage}{\linewidth}\raggedright
% mdparagraph:P573:begin
\subsection*{Cook 47.4}
% mdparagraph:P573:end

% item: EX-XII-XIII-47-04
% mdparagraph:P574:begin
Then said he, With whomsoever I find it, let him be my servant, and let the others be blameless (clǣne).
% mdparagraph:P574:end

% mdparagraph:P575:begin
\textbf{Hint:} clǣne.
% mdparagraph:P575:end

\end{minipage}\par\addvspace{4pt}

% teaching: NEW-U18-GRAM-U18-H-47-04
% mdparagraph:P576:begin
\textbf{Supplied:} Mid swā hwǣm swā ic hit finde, sīe hē mīn ðēow, and sīen ðā ōðre clǣne. — The free-choice dative person is governed by mid within the indefinite condition; hē resumes that person as subject of the singular wish. The other people are plural subject of sīen. Clǣne means clear of blame here. Preserve the scope contrast between the one offender and the other people.
% mdparagraph:P576:end

% output: EX-XII-XIII-47-04-O001
\par\begin{minipage}{\linewidth}\raggedright
% mdparagraph:P577:begin
\subsection*{Cook 47.5}
% mdparagraph:P577:end

% item: EX-XII-XIII-47-05
% mdparagraph:P578:begin
And he sought from (fram) the eldest unto [until] the youngest; and he found the cup in Benjamin’s sack; and they turned back to the city.
% mdparagraph:P578:end

% mdparagraph:P579:begin
\textbf{Hint:} fram.
% mdparagraph:P579:end

\end{minipage}\par\addvspace{4pt}

% teaching: U18-SUP-OLDEST-DAT-M-SG-SUBSTANTIVE
% mdparagraph:P580:begin
\textbf{Supplied:} ðǣm ieldestan — The youngest and eldest refer to individual brothers. Supply the masculine singular human referents and each case, rather than inferring substantive person uses from familiar adjective headwords.
% mdparagraph:P580:end

% mdparagraph:P581:begin
\textbf{Supplied:} ðone gingstan; — The youngest and eldest refer to individual brothers. Supply the masculine singular human referents and each case, rather than inferring substantive person uses from familiar adjective headwords.
% mdparagraph:P581:end

% mdparagraph:P582:begin
\textbf{Supplied:} ieldestan — oldest person; eldest brother; dative, masculine, singular.
% mdparagraph:P582:end

% mdparagraph:P583:begin
\textbf{Supplied:} gingstan — youngest; accusative, masculine, singular.
% mdparagraph:P583:end

% teaching: U18-SUP-BENIAMIN-GEN
% mdparagraph:P584:begin
\textbf{Supplied:} Beniamines — Benjamin; genitive, masculine, singular.
% mdparagraph:P584:end

% mdparagraph:P585:begin
\textbf{[until]} Use the bracketed English as the explicit target wording, preserving the complete sentence and every participant: until.
% mdparagraph:P585:end

% output: EX-XII-XIII-47-05-O001
\par\begin{minipage}{\linewidth}\raggedright
% mdparagraph:P586:begin
\subsection*{Cook 47.6}
% mdparagraph:P586:end

% item: EX-XII-XIII-47-06
% mdparagraph:P587:begin
And Judah went in first (fyrmest) with his brethren to Joseph, and they fell all together before Joseph.
% mdparagraph:P587:end

% mdparagraph:P588:begin
\textbf{Hint:} fyrmest.
% mdparagraph:P588:end

% mdparagraph:P589:begin
\textbf{Supplied:} Iudas — Judah; nominative, masculine, singular.
% mdparagraph:P589:end

\end{minipage}\par\addvspace{4pt}

% teaching: U18-SUP-FYRMEST-INVARIANT
% mdparagraph:P590:begin
\textbf{Supplied:} fyrmest — first of all; at the outset.
% mdparagraph:P590:end

% mdparagraph:P591:begin
\textbf{Supplied:} Iosepe — Joseph; dative, masculine, singular.
% mdparagraph:P591:end

% mdparagraph:P592:begin
\textbf{Supplied:} ætgædere — together; at once.
% mdparagraph:P592:end

% mdparagraph:P593:begin
\textbf{Supplied:} Iosepe — Joseph; dative, masculine, singular.
% mdparagraph:P593:end

% output: EX-XII-XIII-47-06-O001
\par\begin{minipage}{\linewidth}\raggedright
% mdparagraph:P594:begin
\subsection*{Cook 47.7}
% mdparagraph:P594:end

% item: EX-XII-XIII-47-07
% mdparagraph:P595:begin
And he said to (tō) them, Why would ye do so (verb last)?
% mdparagraph:P595:end

% mdparagraph:P596:begin
\textbf{Hint:} tō.
% mdparagraph:P596:end

% mdparagraph:P597:begin
\textbf{(verb last)} Put the finite verb last in this specified clause. This local composition target is not a universal Old English order rule.
% mdparagraph:P597:end

% output: EX-XII-XIII-47-07-O001
\end{minipage}\par\addvspace{4pt}

\par\begin{minipage}{\linewidth}\raggedright
% mdparagraph:P598:begin
\subsection*{Cook 47.8}
% mdparagraph:P598:end

% item: EX-XII-XIII-47-08
% mdparagraph:P599:begin
Let him be my servant who stole the cup, and go ye free to your father.
% mdparagraph:P599:end

\end{minipage}\par\addvspace{4pt}

% teaching: U18-SUP-FRIGE-NOM-M-PL
% mdparagraph:P600:begin
\textbf{Supplied:} frige — free; nominative, masculine, plural.
% mdparagraph:P600:end

% output: EX-XII-XIII-47-08-O001
\par\begin{minipage}{\linewidth}\raggedright
% mdparagraph:P601:begin
\subsection*{Cook 47.9}
% mdparagraph:P601:end

% item: EX-XII-XIII-47-09
% mdparagraph:P602:begin
Then said Judah, Thou askedst us before whether we had father or brother.
% mdparagraph:P602:end

% mdparagraph:P603:begin
\textbf{Supplied:} Iudas — Judah; nominative, masculine, singular.
% mdparagraph:P603:end

% output: EX-XII-XIII-47-09-O001
\end{minipage}\par\addvspace{4pt}

\par\begin{minipage}{\linewidth}\raggedright
% mdparagraph:P604:begin
\subsection*{Cook 47.10}
% mdparagraph:P604:end

% item: EX-XII-XIII-47-10
% mdparagraph:P605:begin
And we said, We have (om. an) old father, and he hath with him our youngest brother, and he loveth him alone (ænne) above (ofer) us all.
% mdparagraph:P605:end

% mdparagraph:P606:begin
\textbf{Hint:} ænne.
% mdparagraph:P606:end

% mdparagraph:P607:begin
\textbf{Hint:} ofer.
% mdparagraph:P607:end

\end{minipage}\par\addvspace{4pt}

% teaching: NEW-U18-GRAM-U18-HELP-OFER-PREFERENCE
% teaching: U18-SUP-PREFERENCE
% mdparagraph:P608:begin
\textbf{Supplied:} ofer ūs ealle. — Above us all /\allowbreak{} more than all of us, describing the father’s preference. Ūs and ealle are accusative plural.
% mdparagraph:P608:end

% mdparagraph:P609:begin
\textbf{Supplied:} gingstan — youngest; accusative, masculine, singular.
% mdparagraph:P609:end

% teaching: U18-SUP-AENNE-ACC-M-SG-SOURCE
% mdparagraph:P610:begin
\textbf{Supplied:} ænne — one; single; accusative, masculine, singular.
% mdparagraph:P610:end

% mdparagraph:P611:begin
\textbf{Supplied:} ofer — above; more than, in preference; governs accusative.
% mdparagraph:P611:end

% mdparagraph:P612:begin
\textbf{(om. an)} Omit a separate Old English word for the marked English expression an. Preserve its grammatical relation through the surrounding selected construction.
% mdparagraph:P612:end

% output: EX-XII-XIII-47-10-O001
\par\begin{minipage}{\linewidth}\raggedright
% mdparagraph:P613:begin
\subsection*{Cook 47.11}
% mdparagraph:P613:end

% item: EX-XII-XIII-47-11
% mdparagraph:P614:begin
And thou commandest us that we (om. should) bring him to thee, that thou (om. mightest) see him, and know by (be) that that we were not spies (scēawere, 44).
% mdparagraph:P614:end

% mdparagraph:P615:begin
\textbf{Hint:} be.
% mdparagraph:P615:end

% mdparagraph:P616:begin
\textbf{Hint:} scēawere.
% mdparagraph:P616:end

\end{minipage}\par\addvspace{4pt}

% teaching: NEW-U18-GRAM-U18-H-47-11
% teaching: U18-SUP-INFERENCE
% mdparagraph:P617:begin
\textbf{Supplied:} ðū bīetst ūs ðæt wē hine ðē bringen, ðæt ðū hine gesēo, and wite be ðǣm ðæt wē ne wǣren scēaweras. — Retain actual present commandest: bīetst is present second singular, not an editorial backshift. The commanded bringing has first-plural present subjunctive; seeing and knowing are coordinated second-singular present subjunctives under purpose. The final reported status preserves original were as first-plural past subjunctive. Be ðǣm means by that evidence, and its following ðæt introduces the content known. The complete tense and proof chain is supplied.
% mdparagraph:P617:end

% mdparagraph:P618:begin
\textbf{Supplied:} be — by/\allowbreak{}from that; on that evidence.
% mdparagraph:P618:end

% mdparagraph:P619:begin
\textbf{(om. should)} Omit a separate Old English word for the marked English expression should. Preserve its grammatical relation through the surrounding selected construction.
% mdparagraph:P619:end

% mdparagraph:P620:begin
\textbf{(om. mightest)} Omit a separate Old English word for the marked English expression mightest. Preserve its grammatical relation through the surrounding selected construction.
% mdparagraph:P620:end

% mdparagraph:P621:begin
\textbf{(scēawere, 44)} The number refers to the noun’s declension. Use the complete scēawere table in section 2 and its stated stem rule.
% mdparagraph:P621:end

% output: EX-XII-XIII-47-11-O001
\reserve{7\baselineskip}
% mdparagraph:P622:begin
\section*{Cook 48. Recognition and invitation}
% mdparagraph:P622:end

\par\begin{minipage}{\linewidth}\raggedright
% mdparagraph:P623:begin
\subsection*{Cook 48.1}
% mdparagraph:P623:end

% item: EX-XII-XIII-48-01
% mdparagraph:P624:begin
And I swore oaths (om. to) my father that I (om. would) bring him home again to (tō) him, and said, Except I bring him again back to thee, let me be ever guilty (scyldig) toward (wið) thee.
% mdparagraph:P624:end

% mdparagraph:P625:begin
\textbf{Hint:} tō.
% mdparagraph:P625:end

% mdparagraph:P626:begin
\textbf{Hint:} scyldig.
% mdparagraph:P626:end

% mdparagraph:P627:begin
\textbf{Hint:} wið.
% mdparagraph:P627:end

% mdparagraph:P628:begin
\textbf{Supplied:} Būtan — unless; except if.
% mdparagraph:P628:end

% mdparagraph:P629:begin
\textbf{Supplied:} ā — ever; always.
% mdparagraph:P629:end

% mdparagraph:P630:begin
\textbf{Supplied:} wið — toward; in relation to a person.
% mdparagraph:P630:end

% mdparagraph:P631:begin
\textbf{(om. to)} Omit a separate Old English word for the marked English expression to. Preserve its grammatical relation through the surrounding selected construction.
% mdparagraph:P631:end

% mdparagraph:P632:begin
\textbf{(om. would)} Omit a separate Old English word for the marked English expression would. Preserve its grammatical relation through the surrounding selected construction.
% mdparagraph:P632:end

% output: EX-XII-XIII-48-01-O001
\end{minipage}\par\addvspace{4pt}

\par\begin{minipage}{\linewidth}\raggedright
% mdparagraph:P633:begin
\subsection*{Cook 48.2}
% mdparagraph:P633:end

% item: EX-XII-XIII-48-02
% mdparagraph:P634:begin
My lord, let (lǣtan) the lad go home with his brethren, and (om. let) me be thy servant instead of (for, 172) him.
% mdparagraph:P634:end

% mdparagraph:P635:begin
\textbf{Hint:} lǣtan.
% mdparagraph:P635:end

% mdparagraph:P636:begin
\textbf{Hint:} for.
% mdparagraph:P636:end

\end{minipage}\par\addvspace{4pt}

% teaching: NEW-U18-GRAM-U18-H-48-02
% mdparagraph:P637:begin
\textbf{Supplied:} lǣt ðone cnapan mid his brōðrum hām gān, — Complete permission frame: singular lǣt addresses Joseph; accusative ðone cnapan is the actor of plain infinitive gān. The boy goes home with his brothers. The accusative actor is not the recipient of a finite command and gān has no person ending. This local frame does not infer unrestricted causative/\allowbreak{}permissive infinitives from the familiar verb alone.
% mdparagraph:P637:end

% teaching: U18-SUP-SUBSTITUTION
% mdparagraph:P638:begin
\textbf{Supplied:} for — instead of; in place of.
% mdparagraph:P638:end

% mdparagraph:P639:begin
\textbf{(om. let)} Omit a separate Old English word for the marked English expression let. Preserve its grammatical relation through the surrounding selected construction.
% mdparagraph:P639:end

% mdparagraph:P640:begin
\textbf{(for, 172)} For plus accusative marks substitution: for hine means instead of him/\allowbreak{}in his place. It is distinct from the causal for + dative used of the money in 45.6.
% mdparagraph:P640:end

% output: EX-XII-XIII-48-02-O001
\par\begin{minipage}{\linewidth}\raggedright
% mdparagraph:P641:begin
\subsection*{Cook 48.3}
% mdparagraph:P641:end

% item: EX-XII-XIII-48-03
% mdparagraph:P642:begin
I dare not go home without the lad, lest I see my father’s grief.
% mdparagraph:P642:end

\end{minipage}\par\addvspace{4pt}

% teaching: NEW-U18-GRAM-U18-H-48-03
% mdparagraph:P643:begin
\textbf{Supplied:} Ic ne dearr hām gān būtan ðǣm cnapan, ðȳ-lǣs-ðe ic mīnes fæder sārnesse gesēo. — Negative present dare plus plain go has the same speaker subject. Būtan with dative denotes absence of the boy, not the unless conjunction used elsewhere. The following negative-purpose clause lest I see uses first-singular present subjunctive gesēo; it does not assert that the grief was seen. Father is genitive possessor within accusative grief.
% mdparagraph:P643:end

% output: EX-XII-XIII-48-03-O001
\par\begin{minipage}{\linewidth}\raggedright
% mdparagraph:P644:begin
\subsection*{Cook 48.4}
% mdparagraph:P644:end

% item: EX-XII-XIII-48-04
% mdparagraph:P645:begin
Then could (mugan) not Joseph conceal (diernan, 113) himself longer, but he drove all the others out, that no foreign man (om. should) be among them.
% mdparagraph:P645:end

% mdparagraph:P646:begin
\textbf{Hint:} mugan.
% mdparagraph:P646:end

% mdparagraph:P647:begin
\textbf{Hint:} diernan.
% mdparagraph:P647:end

% mdparagraph:P648:begin
\textbf{Supplied:} Iosep — Joseph; nominative, masculine, singular.
% mdparagraph:P648:end

\end{minipage}\par\addvspace{4pt}

% teaching: U18-SUP-LENG-INVARIANT
% mdparagraph:P649:begin
\textbf{Supplied:} lęng — longer.
% mdparagraph:P649:end

% mdparagraph:P650:begin
\textbf{Supplied:} betweox — among; between.
% mdparagraph:P650:end

% mdparagraph:P651:begin
\textbf{(diernan, 113)} Diernan is a long weak class I verb; the assigned system here is the complete plain/\allowbreak{}inflected infinitive pair diernan/\allowbreak{}diernanne, not a silently expanded finite paradigm. The plain infinitive follows past ability meahte, and accusative hine is reflexive.
% mdparagraph:P651:end

% mdparagraph:P652:begin
\textbf{(om. should)} Omit a separate Old English word for the marked English expression should. Preserve its grammatical relation through the surrounding selected construction.
% mdparagraph:P652:end

% output: EX-XII-XIII-48-04-O001
\par\begin{minipage}{\linewidth}\raggedright
% mdparagraph:P653:begin
\subsection*{Cook 48.5}
% mdparagraph:P653:end

% item: EX-XII-XIII-48-05
% mdparagraph:P654:begin
And he wept and cried (clēopian) with a loud (hlūd) voice (160. 1); and he said to (tō) his brethren, I am Joseph; (om. doth) our father live yet (nū gīet)?
% mdparagraph:P654:end

% mdparagraph:P655:begin
\textbf{Hint:} clēopian.
% mdparagraph:P655:end

% mdparagraph:P656:begin
\textbf{Hint:} hlūd.
% mdparagraph:P656:end

% mdparagraph:P657:begin
\textbf{Hint:} tō.
% mdparagraph:P657:end

% mdparagraph:P658:begin
\textbf{Hint:} nū gīet.
% mdparagraph:P658:end

\end{minipage}\par\addvspace{4pt}

% teaching: NEW-U18-GRAM-U18-H-48-05
% mdparagraph:P659:begin
\textbf{Supplied:} clēopode hlūdre stefne; — Exact 48.5 hint has marked clēopian. Supply the derived weak class II past third singular clēopode and the dative/\allowbreak{}instrumental voice phrase meaning cried with a loud voice. This marked source variant is supplied locally; retain the earlier spelling in its own paradigm.
% mdparagraph:P659:end

% mdparagraph:P660:begin
\textbf{Supplied:} clēopode — cry aloud; indicative, preterite, third person, singular.
% mdparagraph:P660:end

% mdparagraph:P661:begin
\textbf{Supplied:} Iosep — Joseph; nominative, masculine, singular.
% mdparagraph:P661:end

% mdparagraph:P662:begin
\textbf{(160. 1)} Hlūdre stefne is a dative/\allowbreak{}instrumental manner phrase, with a loud voice; no extra mid is required. The feminine adjective agrees with singular stefne.
% mdparagraph:P662:end

% mdparagraph:P663:begin
\textbf{(om. doth)} Omit a separate Old English word for the marked English expression doth. Preserve its grammatical relation through the surrounding selected construction.
% mdparagraph:P663:end

% output: EX-XII-XIII-48-05-O001
\par\begin{minipage}{\linewidth}\raggedright
% mdparagraph:P664:begin
\subsection*{Cook 48.6}
% mdparagraph:P664:end

% item: EX-XII-XIII-48-06
% mdparagraph:P665:begin
Fear ye (184) nothing, nor (om. let) it grieve (190) you that ye sold (sęllan) me into this kingdom; indeed (om. it was) for your benefit God sent me into the land of the Egyptians.
% mdparagraph:P665:end

% mdparagraph:P666:begin
\textbf{Hint:} sęllan.
% mdparagraph:P666:end

\end{minipage}\par\addvspace{4pt}

% teaching: NEW-U18-GRAM-U18-HELP-FOR-BENEFIT
% teaching: U18-SUP-BENEFIT
% mdparagraph:P667:begin
\textbf{Supplied:} for ēowerre ðearfe — For your good /\allowbreak{} for your benefit. Ēowerre and ðearfe are dative feminine singular. Use the whole phrase visibly.
% mdparagraph:P667:end

% teaching: U18-SUP-SELLING-GOAL
% mdparagraph:P668:begin
\textbf{Supplied:} on — into; destination of sale.
% mdparagraph:P668:end

% mdparagraph:P669:begin
\textbf{Supplied:} for — for; to the benefit of; governs dative.
% mdparagraph:P669:end

% mdparagraph:P670:begin
\textbf{Supplied:} Egypta — Egyptians; genitive, masculine, plural.
% mdparagraph:P670:end

% mdparagraph:P671:begin
\textbf{(184)} Reflexive dative ēow belongs with fear; nānes ðinges is the separately licensed genitive feared object. The negation is concordant: do not fear anything, not an affirmative double-negative reading.
% mdparagraph:P671:end

% mdparagraph:P672:begin
\textbf{(om. let)} Omit a separate Old English word for the marked English expression let. Preserve its grammatical relation through the surrounding selected construction.
% mdparagraph:P672:end

% mdparagraph:P673:begin
\textbf{(190)} Ofðynce is third-singular impersonal subjunctive, let it not grieve/\allowbreak{}displease you. Dative ēow gives the experiencers; the following ðæt clause gives what troubles them. Plural experiencers do not cause plural verbal agreement.
% mdparagraph:P673:end

% mdparagraph:P674:begin
\textbf{(om. it was)} Omit a separate Old English word for the marked English expression it was. Preserve its grammatical relation through the surrounding selected construction.
% mdparagraph:P674:end

% output: EX-XII-XIII-48-06-O001
\par\begin{minipage}{\linewidth}\raggedright
% mdparagraph:P675:begin
\subsection*{Cook 48.7}
% mdparagraph:P675:end

% item: EX-XII-XIII-48-07
% mdparagraph:P676:begin
Now two years has been [was] famine over (172) all (om. the) earth, and (om. there) yet shall (om. be) five, in which one can (mugan) neither (89. a) plow (ęrian, 116) nor reap (rīpan, 102).
% mdparagraph:P676:end

% mdparagraph:P677:begin
\textbf{Hint:} mugan.
% mdparagraph:P677:end

% mdparagraph:P678:begin
\textbf{Hint:} ęrian.
% mdparagraph:P678:end

% mdparagraph:P679:begin
\textbf{Hint:} rīpan.
% mdparagraph:P679:end

\end{minipage}\par\addvspace{4pt}

% teaching: NEW-U18-GRAM-U18-H-48-07
% teaching: NEW-U18-VOC-SCULAN-U18-FUTURE-ELLIPSIS
% teaching: U18-SUP-NAHWAETHER-INVARIANT
% mdparagraph:P680:begin
\textbf{Supplied:} Nū twā gēar hunger wæs ofer ealle eorðan, and gīet sculon fīf, on ðǣm mann nāhwæðer ne ęrian ne rīpan mæg. — Two years is accusative duration; singular hunger governs wæs. Ofer takes accusative extent ealle eorðan, as 172 requires. Fīf stands for five further years; the source explicitly omits there and be, so sculon has an understood existential complement, not a missing translation item. On ðǣm refers to those years. Indefinite mann is singular subject of mæg; nāhwæðer plus ne...ne excludes both infinitive activities. Cook 89 a supplies the full pronoun; All ellipsis, future use and negative coordination are supplied together.
% mdparagraph:P680:end

% mdparagraph:P681:begin
\textbf{Supplied:} sculon — shall; be yet to come (with an omitted be); indicative, present, third person, plural.
% mdparagraph:P681:end

% mdparagraph:P682:begin
\textbf{Supplied:} nāhwæðer — neither.
% mdparagraph:P682:end

% mdparagraph:P683:begin
\textbf{[was]} Use the bracketed English as the explicit target wording, preserving the complete sentence and every participant: was.
% mdparagraph:P683:end

% mdparagraph:P684:begin
\textbf{(172)} Ofer with accusative ealle eorðan marks extent over all the earth. The noun’s identical weak ending does not make it a dative; the full phrase has the accusative relation required by 172.
% mdparagraph:P684:end

% mdparagraph:P685:begin
\textbf{(om. the)} Omit a separate Old English word for the marked English expression the. Preserve its grammatical relation through the surrounding selected construction.
% mdparagraph:P685:end

% mdparagraph:P686:begin
\textbf{(om. there)} Omit a separate Old English word for the marked English expression there. Preserve its grammatical relation through the surrounding selected construction.
% mdparagraph:P686:end

% mdparagraph:P687:begin
\textbf{(om. be)} Omit a separate Old English word for the marked English expression be. Preserve its grammatical relation through the surrounding selected construction.
% mdparagraph:P687:end

% mdparagraph:P688:begin
\textbf{(89. a)} Nāhwæðer means neither. In the supplied whole clause, nāhwæðer with paired ne ... ne excludes plowing and reaping under singular mæg; mann is the singular indefinite subject. The lexical hint 89 a alone does not teach the coordination, so the complete negative frame is visible.
% mdparagraph:P688:end

% mdparagraph:P689:begin
\textbf{(ęrian, 116)} The source spells ęrian, with ogonek, and assigns short weak class I 116. The complete assigned nonfinite pair is ęrian/\allowbreak{}ęrianne; only plain ęrian occurs after mæg. No new full finite weak class I system is inferred.
% mdparagraph:P689:end

% mdparagraph:P690:begin
\textbf{(rīpan, 102)} Rīpan is strongI 102. The complete assigned infinitive pair rīpan/\allowbreak{}rīpanne preserves the source’s long ī; only plain rīpan occurs after mæg. A class reference does not itself assign all finite forms.
% mdparagraph:P690:end

% output: EX-XII-XIII-48-07-O001
\par\begin{minipage}{\linewidth}\raggedright
% mdparagraph:P691:begin
\subsection*{Cook 48.8}
% mdparagraph:P691:end

% item: EX-XII-XIII-48-08
% mdparagraph:P692:begin
Go quickly to my father and tell him that God has set me as a (tō) lord (om. to) all (om. the) Egyptians; bid (bēodan) him (164) that he come to me and dwell in this land.
% mdparagraph:P692:end

% mdparagraph:P693:begin
\textbf{Hint:} tō.
% mdparagraph:P693:end

% mdparagraph:P694:begin
\textbf{Hint:} bēodan.
% mdparagraph:P694:end

\end{minipage}\par\addvspace{4pt}

% teaching: NEW-U18-GRAM-U18-H-48-08
% teaching: NEW-U18-VOC-GESETTAN-U18-APPOINT-LORD
% teaching: U18-SUP-EGYPTUM-DAT-PL
% mdparagraph:P695:begin
\textbf{Supplied:} God mē tō hlāforde eallum Egyptum gesętte; — God appointed me as lord to all Egyptians. Mē is the accusative appointee, tō hlāforde the role complement, and eallum Egyptum the dative people governed; use the bare dative for the people, following the instruction to omit a separate equivalent of English to. The exact name phrase is supplied.
% mdparagraph:P695:end

% mdparagraph:P696:begin
\textbf{Supplied:} hrædlīce — quickly.
% mdparagraph:P696:end

% mdparagraph:P697:begin
\textbf{Supplied:} Egyptum — Egyptians; dative, masculine, plural.
% mdparagraph:P697:end

% mdparagraph:P698:begin
\textbf{Supplied:} gesętte — appoint as lord; indicative, preterite, third person, singular.
% mdparagraph:P698:end

% mdparagraph:P699:begin
\textbf{(om. to)} Omit a separate Old English word for the marked English expression to. Preserve its grammatical relation through the surrounding selected construction.
% mdparagraph:P699:end

% mdparagraph:P700:begin
\textbf{(om. the)} Omit a separate Old English word for the marked English expression the. Preserve its grammatical relation through the surrounding selected construction.
% mdparagraph:P700:end

% mdparagraph:P701:begin
\textbf{(164)} Him is the dative addressee of plural command bēodað; the father is subject of cume and wunie in the commanded ðæt clause. Keep the commanding brothers distinct from the father who is to come and dwell.
% mdparagraph:P701:end

% output: EX-XII-XIII-48-08-O001
\par\begin{minipage}{\linewidth}\raggedright
% mdparagraph:P702:begin
\subsection*{Cook 48.9}
% mdparagraph:P702:end

% item: EX-XII-XIII-48-09
% mdparagraph:P703:begin
Do thus, that ye perish not.
% mdparagraph:P703:end

% output: EX-XII-XIII-48-09-O001
\end{minipage}\par\addvspace{4pt}

\par\begin{minipage}{\linewidth}\raggedright
% mdparagraph:P704:begin
\subsection*{Cook 48.10}
% mdparagraph:P704:end

% item: EX-XII-XIII-48-10
% mdparagraph:P705:begin
Now ye see how it is with me, and ye hear what I say (sprecan) to (tō) you.
% mdparagraph:P705:end

% mdparagraph:P706:begin
\textbf{Hint:} sprecan.
% mdparagraph:P706:end

% mdparagraph:P707:begin
\textbf{Hint:} tō.
% mdparagraph:P707:end

\end{minipage}\par\addvspace{4pt}

% teaching: NEW-U18-GRAM-U18-H-48-10
% teaching: NEW-U18-VOC-GESEON-U18-STATE-CONTENT
% mdparagraph:P708:begin
\textbf{Supplied:} Nū gē gesēoð hū hit mid mē is, and gē hīerað hwæt ic tō ēow sprece. — Hū introduces how things stand with the speaker; this is a state-of-affairs expression, not physical accompaniment. Hwæt is what he says, heard by plural gē. Sprece retains the first-singular speaker, and tō ēow the plural addressee.
% mdparagraph:P708:end

% mdparagraph:P709:begin
\textbf{Supplied:} gesēoð — see how a state of affairs stands; indicative, present, second person, plural.
% mdparagraph:P709:end

% output: EX-XII-XIII-48-10-O001
\reserve{7\baselineskip}
% mdparagraph:P710:begin
\section*{7. The gifts and the end of the banquet}
% mdparagraph:P710:end

% mdparagraph:P711:begin
The king’s daughter has heard and watched Apollonius. The continuation moves through her gifts, their delivery and the farewells. The exact passage is Cook 175.18–177.2, ending before the following teacher scene. Read it with the original notes and modern help below.
% mdparagraph:P711:end

% mdparagraph:P712:begin
Sōðlice, mid-ðȳ-ðe ðæs cyninges dohtor geseah ðæt Apollonius on eallum gōdum cræftum swā wel wæs getogen,\textsuperscript{P175N10} ðā gefēoll hiere mōd on his lufe. Ðā, æfter ðæs bēorscipes geęndunge, cwæð ðæt mǣden tō ðǣm cyninge: “Lēofa\textsuperscript{P175N11} fæder, ðū līefdest mē, lȳtle\textsuperscript{P176N1} ǣr,\textsuperscript{P176N1} ðæt\textsuperscript{P176N2} ic\textsuperscript{P176N2} mōste\textsuperscript{P176N2} giefan Apollonio swā-hwæt-swā ic wolde of ðīnum goldhorde.” Arcestrates se cyning cwæð tō hiere: “Gief him swā-hwæt-swā ðū wille.\textsuperscript{P176N3}” Hēo ðā swīðe blīðe ūtēode,\textsuperscript{P176N4} and cwæð: “Lārēow\textsuperscript{P176N5} Apolloni, ic giefe ðē, be mīnes fæder lēafe, twā hund punda\textsuperscript{P176N6} goldes,\textsuperscript{P176N7} and fēower hund punda\textsuperscript{P176N6} gewihte\textsuperscript{P176N8} seolfres,\textsuperscript{P176N7} and ðone mǣstan dǣl dēorwurðes\textsuperscript{P176N9} rēafes, and twēntig ðēowa\textsuperscript{P176N10} manna.\textsuperscript{P176N10}” And hēo ðā ðus cwæð tō ðǣm ðēowum mannum: “Berað ðās ðing mid ēow, ðe ic behēt Apollonio mīnum lārēowe, and lęcgeað innan būre\textsuperscript{P176N11} beforan mīnum frēondum.” Ðis wearð ðā ðus gedōn, æfter ðǣre cwēne hǣse\textsuperscript{P176N12}; and ealle ðā męnn hiere giefa hęredon ðe\textsuperscript{P176N13} hīe gesāwon. Ðā sōðlice geęndode se gebēorscipe, and ðā męnn ealle ārison,\textsuperscript{P176N14} and grētton ðone cyning and ðā cwēne, and bǣdon hīe gesunde bēon,\textsuperscript{P176N15} and hām gewęndon. Ēac swilce Apollonius cwæð: “Ðū gōda cyning, and earmra\textsuperscript{P176N16} gemiltsiend, and ðū cwēn, lāre\textsuperscript{P176N16} lufiend, bēon gē gesunde.\textsuperscript{P176N17}” Hē beseah ēac tō ðǣm ðēowum mannum, ðe ðæt mǣden him forgiefen hæfde,\textsuperscript{P176N18} and him cwæð tō: “Nimað ðās ðing mid ēow, ðe mē sēo cwēn forgeaf, and gān\textsuperscript{P177N1} wē sēcean ūre giesthūs, ðæt wē mægen ūs\textsuperscript{P177N2} geręstan.”
% mdparagraph:P712:end

\reserve{18\baselineskip}
% mdparagraph:P713:begin
\subsection*{Finding the cited lines}
% mdparagraph:P713:end

% mdparagraph:P714:begin
These locators retain the source’s physical line divisions, including divided words. The running text joins those words.
% mdparagraph:P714:end

% mdtable:T65:begin
\par\begingroup
\fontsize{10.5}{13.2}\selectfont
\setlength{\LTleft}{0pt}\setlength{\LTright}{\fill}
\begin{longtable}{@{}>{\raggedright\arraybackslash}p{35.03mm}>{\raggedright\arraybackslash}p{86.16mm}@{}}
\toprule
\textbf{Page and line} & \textbf{Beginning of the printed line} \\
\midrule
\endfirsthead
\toprule
\textbf{Page and line} & \textbf{Beginning of the printed line} \\
\midrule
\endhead
\midrule
\multicolumn{2}{r}{\footnotesize Continued on the following page.}\\
\endfoot
\bottomrule
\endlastfoot
175.18 & Sōðlice, mid-ðȳ-ðe ðæs cyninges dohtor … \\
175.19 & lonius on eallum gōdum cræftum … \\
175.20 & ðā gefēoll hiere mōd on … \\
175.21 & geęndunge, cwæð ðæt mǣden tō … \\
176.1 & fæder, ðū līefdest mē, lȳtle … \\
176.2 & Apollonio swā-hwæt-swā ic wolde of … \\
176.3 & Arcestrates se cyning cwæð tō … \\
176.4 & hwæt-swā ðū wille.” Hēo ðā … \\
176.5 & cwæð: “Lārēow Apolloni, ic giefe … \\
176.6 & lēafe, twā hund punda goldes, … \\
176.7 & gewihte seolfres, and ðone mǣstan … \\
176.8 & rēafes, and twēntig ðēowa manna.” … \\
176.9 & cwæð tō ðǣm ðēowum mannum: … \\
176.10 & ēow, ðe ic behēt Apollonio … \\
176.11 & innan būre beforan mīnum frēondum.” … \\
176.12 & ðus gedōn, æfter ðǣre cwēne … \\
176.13 & hiere giefa hęredon ðe hīe … \\
176.14 & ode se gebēorscipe, and ðā … \\
176.15 & grētton ðone cyning and ðā … \\
176.16 & gesunde bēon, and hām gewęndon. … \\
176.17 & lonius cwæð: “Ðū gōda cyning, … \\
176.18 & and ðū cwēn, lāre lufiend, … \\
176.19 & beseah ēac tō ðǣm ðēowum … \\
176.20 & forgiefen hæfde, and him cwæð … \\
177.1 & ēow, ðe mē sēo cwēn … \\
177.2 & giesthūs, ðæt wē mægen ūs … \\
\end{longtable}
\par\endgroup

% mdtable:T65:end
\reserve{7\baselineskip}
% mdparagraph:P715:begin
\subsection*{Modern help for the reading}
% mdparagraph:P715:end

% mdparagraph:P716:begin
The explanations supply unfamiliar senses and whole constructions. Where an original note already identifies a referent, the later question asks for an explanation with that aid visible. Source spelling is retained.
% mdparagraph:P716:end

\reserve{7\baselineskip}
% mdparagraph:P717:begin
\subsection*{175.18–20: seeing and falling in love}
% mdparagraph:P717:end

% teaching: NEW-U18-GRAM-U18-RSUP-01
% mdparagraph:P718:begin
The daughter sees that a fact holds; indicative wæs belongs to the content clause. Wæs getogen means was educated or accomplished in good skills. In the following idiom, singular mōd is the subject of gefēoll, hiere its possessor, and on his lufe the state into which her mind falls: she falls in love with him.
% mdparagraph:P718:end

\reserve{7\baselineskip}
% mdparagraph:P719:begin
\subsection*{175.20–176.2: the earlier permission}
% mdparagraph:P719:end

% teaching: NEW-U18-GRAM-U18-RSUP-02
% mdparagraph:P720:begin
Līefdest means permitted, with mē as dative person permitted. Ðæt ic mōste giefan is naturally translated by permitted me to give. The free-choice object swā-hwæt-swā means whatever; wolde is first-singular preterite subjunctive and leaves giving understood. Of with dative identifies the treasury from which she may take the gift.
% mdparagraph:P720:end

\reserve{7\baselineskip}
% mdparagraph:P721:begin
\subsection*{176.3–4: renewed permission}
% mdparagraph:P721:end

% teaching: NEW-U18-GRAM-U18-RSUP-03
% mdparagraph:P722:begin
Gief is a singular command to the daughter, with him as recipient. Wille is second-singular present subjunctive within whatever you wish. Its spelling also occurs in other persons and moods; this context fixes the use.
% mdparagraph:P722:end

\reserve{7\baselineskip}
% mdparagraph:P723:begin
\subsection*{176.4: going out}
% mdparagraph:P723:end

% teaching: NEW-U18-GRAM-U18-RSUP-04
% mdparagraph:P724:begin
Hēo … ūtēode means she went out. Swīðe blīðe means very gladly. The exact prefixed preterite is supplied.
% mdparagraph:P724:end

\reserve{7\baselineskip}
% mdparagraph:P725:begin
\subsection*{176.5–8: the gifts}
% mdparagraph:P725:end

% teaching: NEW-U18-GRAM-U18-RSUP-05
% mdparagraph:P726:begin
Giefe is first-singular present indicative; ðē is the recipient. Be mīnes fæder lēafe means with my father’s permission. The following quantities all name what she gives. Their complete weight and servant expressions are supplied below.
% mdparagraph:P726:end

\reserve{7\baselineskip}
% mdparagraph:P727:begin
\subsection*{176.9–11: the instructions}
% mdparagraph:P727:end

% teaching: NEW-U18-GRAM-U18-RSUP-06
% mdparagraph:P728:begin
Berað and lęcgeað are plural commands to the enslaved men: carry the gifts along and place them in the chamber before her friends. In the promise relative, behēt is first-singular preterite, ðe represents the promised things, and Apollonio mīnum lārēowe is the dative recipient. The source lęcgeað is imperative despite matching the newly learned indicative plural in spelling.
% mdparagraph:P728:end

\reserve{7\baselineskip}
% mdparagraph:P729:begin
\subsection*{176.12: the command is carried out}
% mdparagraph:P729:end

% teaching: NEW-U18-GRAM-U18-RSUP-07
% mdparagraph:P730:begin
Ðis wearð … gedōn means this was done. Gedōn is a past participle in the singular passive event, not a present subjunctive plural of the same spelling. The action follows the princess’s command.
% mdparagraph:P730:end

\reserve{7\baselineskip}
% mdparagraph:P731:begin
\subsection*{176.12–13: praise and a separated relative}
% mdparagraph:P731:end

% teaching: NEW-U18-GRAM-U18-RSUP-08
% mdparagraph:P732:begin
All the men who saw the gifts praised her gifts. Original note 13 explicitly identifies męnn as antecedent of ðe. The relative is subject of its clause; hīe is accusative plural, the gifts. The relative stands after the main clause’s object, separated from its own antecedent. Both finite verbs have plural subjects.
% mdparagraph:P732:end

\reserve{7\baselineskip}
% mdparagraph:P733:begin
\subsection*{176.13–16: the guests leave}
% mdparagraph:P733:end

% teaching: NEW-U18-GRAM-U18-RSUP-09
% mdparagraph:P734:begin
Geęndode is intransitive: the banquet ended. The men rose, took leave of king and princess, wished them well and returned home. In bǣdon hīe gesunde bēon, hīe is the accusative actor of bēon and gesunde its agreeing accusative plural predicate. This complete farewell wish is supplied.
% mdparagraph:P734:end

\reserve{7\baselineskip}
% mdparagraph:P735:begin
\subsection*{176.16–18: Apollonius’s farewell}
% mdparagraph:P735:end

% teaching: NEW-U18-GRAM-U18-RSUP-10
% mdparagraph:P736:begin
The king and princess are addressed separately and then resumed together by plural gē. Bēon is second-plural present subjunctive, may you be; gesunde is nominative plural agreeing with gē. The agent nouns and their genitive complements are supplied below.
% mdparagraph:P736:end

\reserve{7\baselineskip}
% mdparagraph:P737:begin
\subsection*{176.18–20: the servants}
% mdparagraph:P737:end

% teaching: NEW-U18-GRAM-U18-RSUP-11
% mdparagraph:P738:begin
Beseah tō means looked toward, with a dative target. In the relative, the maiden is subject, the enslaved men the object represented by ðe, and singular him the recipient Apollonius. Forgiefen hæfde is active past perfect. The later him is plural and refers to the servants; tō follows the verb but still belongs to that dative addressee.
% mdparagraph:P738:end

\reserve{7\baselineskip}
% mdparagraph:P739:begin
\subsection*{177.1–2: the proposed departure}
% mdparagraph:P739:end

% teaching: NEW-U18-GRAM-U18-RSUP-12
% mdparagraph:P740:begin
Nimað is a plural command to take the gifts along. Gān wē is a first-plural subjunctive exhortation, let us go; supplied sēcean states the purpose of going. In the gift relative, the princess is subject and mē recipient. Final ðæt expresses purpose, with plural mægen and accusative reflexive ūs before geręstan. The words propose seeking the guesthouse and resting; they do not narrate a completed arrival.
% mdparagraph:P740:end

\reserve{7\baselineskip}
% mdparagraph:P741:begin
\subsection*{Supplied expressions and exact word help}
% mdparagraph:P741:end

% teaching: NEW-U18-GRAM-U18-HELP-LITTLE-EARLIER
% mdparagraph:P742:begin
\textbf{lȳtle ǣr} — A little earlier: lȳtle is instrumental of degree/\allowbreak{}difference with the comparative adverb ǣr. Cook’s original note points to 178. It is not a feminine accusative object.
% mdparagraph:P742:end

% teaching: NEW-U18-GRAM-U18-HELP-QUANTITIES
% mdparagraph:P743:begin
\textbf{twā hund punda goldes and fēower hund punda gewihte seolfres} — Two hundred pounds of gold and four hundred pounds by weight of silver. Hund means hundred. Punda is genitive plural after the hundreds; goldes and seolfres are material genitives. Gewihte is a dative expression of measure: by weight.
% mdparagraph:P743:end

% teaching: NEW-U18-GRAM-U18-HELP-PRECIOUS-CLOTHING
% mdparagraph:P744:begin
\textbf{ðone mǣstan dǣl dēorwurðes rēafes} — The greatest part of costly clothing. Mǣstan is weak accusative masculine singular with dǣl; dēorwurðes agrees with the genitive neuter singular rēafes.
% mdparagraph:P744:end

% teaching: NEW-U18-GRAM-U18-HELP-SERVANTS
% mdparagraph:P745:begin
\textbf{twēntig ðēowa manna /\allowbreak{} tō ðǣm ðēowum mannum} — Twenty enslaved men; to the enslaved men. Preserve ðēowa without r. The selected local explanation treats ðēowa and manna as appositional genitive plurals, and ðēowum and mannum as appositional datives. The phrase’s meaning and case are supplied without requiring independent historical classification.
% mdparagraph:P745:end

% teaching: NEW-U18-GRAM-U18-HELP-DELAYED-RELATIVE
% mdparagraph:P746:begin
\textbf{ealle ðā męnn hiere giefa hęredon ðe hīe gesāwon} — All the men who saw them praised her gifts. The relative ðe is subject and reaches back to męnn; hīe is the accusative plural gifts. The relative clause is separated from its antecedent. Keep the people and gifts distinct.
% mdparagraph:P746:end

% teaching: NEW-U18-GRAM-U18-HELP-FAREWELL-ACC
% mdparagraph:P747:begin
\textbf{bǣdon hīe gesunde bēon} — Wished them to be well: hīe refers to the king and princess and is the accusative actor of bēon; gesunde is accusative plural predicate. The complete construction is supplied.
% mdparagraph:P747:end

% teaching: NEW-U18-GRAM-U18-HELP-FAREWELL-NOM
% mdparagraph:P748:begin
\textbf{bēon gē gesunde} — May you be well: gē is nominative plural and gesunde nominative plural predicate. This is a direct wish with a plural subjunctive, distinct from the preceding accusative-infinitive construction despite identical gesunde.
% mdparagraph:P748:end

% teaching: NEW-U18-GRAM-U18-HELP-AGENT-ADDRESSES
% mdparagraph:P749:begin
\textbf{earmra gemiltsiend /\allowbreak{} lāre lufiend} — Pitier of the poor /\allowbreak{} lover of learning. Supply both agent nouns and their object genitives. These phrases address the king and princess respectively; historical or Latin parallels in the original notes are explained in English.
% mdparagraph:P749:end

% teaching: NEW-U18-GRAM-U18-HELP-SERVANT-RELATIVE
% mdparagraph:P750:begin
\textbf{ðēowum mannum ðe ðæt mǣden him forgiefen hæfde} — The enslaved men whom the girl had given to him. Relative ðe is accusative plural; mǣden is the subject; him is Apollonius.
% mdparagraph:P750:end

% teaching: NEW-U18-GRAM-U18-HELP-POSTPOSED-TO
% mdparagraph:P751:begin
\textbf{him cwæð tō} — Said to them: him is dative plural referring to the servants; tō follows the verb while governing him.
% mdparagraph:P751:end

% teaching: NEW-U18-GRAM-U18-HELP-FINAL-PLAN
% mdparagraph:P752:begin
\textbf{gān wē sēcean ūre giesthūs, ðæt wē mægen ūs geręstan} — Let us go seek our guesthouse, so that we may rest. Gān wē is an exhortation, ðæt marks purpose, and ūs is accusative reflexive with geręstan.
% mdparagraph:P752:end

% teaching: U18-SUP-APOLLONIUS-NOM
% mdparagraph:P753:begin
\textbf{Apollonius} — Apollonius; nominative, masculine, singular.
% mdparagraph:P753:end

% teaching: U18-SUP-APOLLONIUS-DAT
% mdparagraph:P754:begin
\textbf{Apollonio} — Apollonius; dative, masculine, singular.
% mdparagraph:P754:end

% teaching: U18-SUP-APOLLONIUS-VOC
% mdparagraph:P755:begin
\textbf{Apolloni} — Apollonius; direct address, masculine, singular.
% mdparagraph:P755:end

% teaching: U18-SUP-ARCESTRATES-NOM
% mdparagraph:P756:begin
\textbf{Arcestrates} — Arcestrates, the king; nominative, masculine, singular.
% mdparagraph:P756:end

% teaching: U18-SUP-LUFE-ACC-SG
% mdparagraph:P757:begin
\textbf{lufe} — love; accusative, feminine, singular.
% mdparagraph:P757:end

% teaching: U18-SUP-BEORSCIPES-GEN-SG
% mdparagraph:P758:begin
\textbf{bēorscipes} — feast; banquet; genitive, masculine, singular.
% mdparagraph:P758:end

% teaching: U18-SUP-GEENDUNGE-DAT-SG
% mdparagraph:P759:begin
\textbf{geęndunge} — end; close; dative, feminine, singular.
% mdparagraph:P759:end

% teaching: U18-SUP-LEAFE-DAT-SG
% mdparagraph:P760:begin
\textbf{lēafe} — permission; dative, feminine, singular.
% mdparagraph:P760:end

% teaching: U18-SUP-GEWIHTE-DAT-SG
% mdparagraph:P761:begin
\textbf{gewihte} — weight; dative, neuter, singular.
% mdparagraph:P761:end

% teaching: U18-SUP-DEORWURTHES-GEN-N-SG
% mdparagraph:P762:begin
\textbf{dēorwurðes} — precious; costly; genitive, neuter, singular.
% mdparagraph:P762:end

% teaching: U18-SUP-REAFES-GEN-SG
% mdparagraph:P763:begin
\textbf{rēafes} — clothing; garments; genitive, neuter, singular.
% mdparagraph:P763:end

% teaching: U18-SUP-BURE-DAT-SG
% mdparagraph:P764:begin
\textbf{būre} — chamber; room; dative, masculine, singular.
% mdparagraph:P764:end

% teaching: U18-SUP-HAESE-DAT-SG
% mdparagraph:P765:begin
\textbf{hǣse} — command; order; dative, feminine, singular.
% mdparagraph:P765:end

% teaching: U18-SUP-GIEFA-ACC-PL
% mdparagraph:P766:begin
\textbf{giefa} — gifts; accusative, feminine, plural.
% mdparagraph:P766:end

% teaching: U18-SUP-CWENE-ACC-SG-VARIANT
% mdparagraph:P767:begin
\textbf{cwēne} — princess; accusative, feminine, singular.
% mdparagraph:P767:end

% teaching: U18-SUP-GEMILTSIEND-ADDRESS
% mdparagraph:P768:begin
\textbf{gemiltsiend} — one who shows mercy; pitier; nominative, masculine, singular.
% mdparagraph:P768:end

% teaching: U18-SUP-LARE-GEN-SG
% mdparagraph:P769:begin
\textbf{lāre} — learning; study; genitive, feminine, singular.
% mdparagraph:P769:end

% teaching: U18-SUP-LUFIEND-ADDRESS
% mdparagraph:P770:begin
\textbf{lufiend} — lover; nominative, masculine, singular.
% mdparagraph:P770:end

% teaching: U18-SUP-BLITHE-JOYOUSLY
% mdparagraph:P771:begin
\textbf{blīðe} — joyfully; gladly.
% mdparagraph:P771:end

% teaching: U18-SUP-THEOWA-GEN-PL-PHRASE
% mdparagraph:P772:begin
\textbf{ðēowa} — enslaved men; servants; genitive, masculine, plural.
% mdparagraph:P772:end

% teaching: U18-SUP-THEOWUM-DAT-PL-PHRASE
% mdparagraph:P773:begin
\textbf{ðēowum} — enslaved men; servants; dative, masculine, plural.
% mdparagraph:P773:end

% teaching: U18-SUP-MOD-NOM-SG
% mdparagraph:P774:begin
\textbf{mōd} — heart; mind; nominative, neuter, singular.
% mdparagraph:P774:end

% teaching: U18-SUP-EARMRA-GEN-PL-SUBSTANTIVE
% mdparagraph:P775:begin
\textbf{earmra} — the poor; unfortunate people; genitive, masculine, plural.
% mdparagraph:P775:end

% teaching: U18-SUP-READING-ADDRESSEE
% mdparagraph:P776:begin
\textbf{tō} — to a speech addressee; governs dative.
% mdparagraph:P776:end

% teaching: U18-SUP-READING-DIRECTED-LOOK
% mdparagraph:P777:begin
\textbf{tō} — toward; target of a look; governs dative.
% mdparagraph:P777:end

% teaching: U18-SUP-READING-PERMISSION
% mdparagraph:P778:begin
\textbf{be} — by; with permission; governs dative.
% mdparagraph:P778:end

% teaching: U18-SUP-READING-LOVE-GOAL
% mdparagraph:P779:begin
\textbf{on} — into love for; governs accusative.
% mdparagraph:P779:end

% teaching: U18-SUP-READING-CONNECTIVE
% mdparagraph:P780:begin
\textbf{sōðlice} — now; indeed; narrative connective.
% mdparagraph:P780:end

% teaching: U18-SUP-READING-ACCOMPLISHMENTS
% mdparagraph:P781:begin
\textbf{cræftum} — skills; accomplishments; dative, masculine, plural.
% mdparagraph:P781:end

% mdparagraph:P782:begin
\textbf{mid-ðȳ-ðe} — when; while.
% mdparagraph:P782:end

% mdparagraph:P783:begin
\textbf{geseah} — see that a stated circumstance holds; indicative, preterite, third person, singular.
% mdparagraph:P783:end

% teaching: NEW-U18-VOC-GETEON-U18-READING
% mdparagraph:P784:begin
\textbf{getogen} — train; educate; past, participle.
% mdparagraph:P784:end

% teaching: NEW-U18-VOC-GEFEALLAN-U18-READING
% mdparagraph:P785:begin
\textbf{gefēoll} — fall, figuratively of the heart falling in love; indicative, preterite, third person, singular.
% mdparagraph:P785:end

% mdparagraph:P786:begin
\textbf{æfter} — after, temporal with dative.
% mdparagraph:P786:end

% teaching: NEW-U18-VOC-LIEFAN-PERMIT-U18-READING
% mdparagraph:P787:begin
\textbf{līefdest} — allow; permit; indicative, preterite, second person, singular.
% mdparagraph:P787:end

% mdparagraph:P788:begin
\textbf{swā-hwæt-swā} — whatsoever; whatever.
% mdparagraph:P788:end

% teaching: NEW-U18-VOC-UTGAN-U18-READING
% mdparagraph:P789:begin
\textbf{ūtēode} — go out; indicative, preterite, third person, singular.
% mdparagraph:P789:end

% mdparagraph:P790:begin
\textbf{hund} — hundred, in the supplied gold and silver weight expressions.
% mdparagraph:P790:end

% mdparagraph:P791:begin
\textbf{twēntig} — twenty.
% mdparagraph:P791:end

% teaching: NEW-U18-VOC-BEHATAN-U18-READING
% mdparagraph:P792:begin
\textbf{behēt} — promise; indicative, preterite, first person, singular.
% mdparagraph:P792:end

% mdparagraph:P793:begin
\textbf{lęcgeað} — lay; place; imperative, second person, plural.
% mdparagraph:P793:end

% mdparagraph:P794:begin
\textbf{innan} — within; inside.
% mdparagraph:P794:end

% teaching: NEW-U18-VOC-GEDON-U18-READING
% mdparagraph:P795:begin
\textbf{gedōn} — do; perform; past, participle.
% mdparagraph:P795:end

% mdparagraph:P796:begin
\textbf{ðe} — that; which; who; nominative, masculine, plural.
% mdparagraph:P796:end

% teaching: NEW-U18-VOC-GEENDIAN-U18-READING
% mdparagraph:P797:begin
\textbf{geęndode} — end; come to an end; indicative, preterite, third person, singular.
% mdparagraph:P797:end

% teaching: NEW-U18-VOC-BIDDAN-U18-FAREWELL
% mdparagraph:P798:begin
\textbf{bǣdon} — wish someone well, in the supplied farewell formula; indicative, preterite, third person, plural.
% mdparagraph:P798:end

% teaching: NEW-U18-VOC-BESEON-U18-READING
% mdparagraph:P799:begin
\textbf{beseah} — look toward; indicative, preterite, third person, singular.
% mdparagraph:P799:end

% mdparagraph:P800:begin
\textbf{sēcean} — seek; plain infinitive.
% mdparagraph:P800:end

% mdparagraph:P801:begin
\textbf{geręstan} — rest; repose; plain infinitive.
% mdparagraph:P801:end

\reserve{7\baselineskip}
% mdparagraph:P802:begin
\subsection*{Reading tasks}
% mdparagraph:P802:end

% mdparagraph:P803:begin
Keep the complete text, original notes and modern help available. Translate the whole assigned passage; the analysis questions explain the supplied relations without testing recall of unfamiliar source forms.
% mdparagraph:P803:end

\reserve{8\baselineskip}
% mdparagraph:P804:begin
\subsection*{R01}
% mdparagraph:P804:end

% item: NEW-U18-R01
% mdparagraph:P805:begin
Translate the complete assigned passage in the following three consecutive divisions. All local help and original notes remain available.
% mdparagraph:P805:end

% output: NEW-U18-R01-O001
% mdparagraph:P806:begin
\itempara{1}{Translate 175.18–176.2, through goldhorde.}
% mdparagraph:P806:end

% output: NEW-U18-R01-O002
% mdparagraph:P807:begin
\itempara{2}{Translate 176.3–176.13, through gesāwon.}
% mdparagraph:P807:end

% output: NEW-U18-R01-O003
% mdparagraph:P808:begin
\itempara{3}{Translate from Ðā sōðlice geęndode in 176.13 through 177.2 geręstan.}
% mdparagraph:P808:end

\reserve{8\baselineskip}
% mdparagraph:P809:begin
\subsection*{R02}
% mdparagraph:P809:end

% item: NEW-U18-R02
% mdparagraph:P810:begin
For each cited form, name the referent and explain its grammatical role. Keep the whole passage available.
% mdparagraph:P810:end

% output: NEW-U18-R02-O001
% mdparagraph:P811:begin
\itempara{1}{Explain hiere and his in 175.20.}
% mdparagraph:P811:end

% output: NEW-U18-R02-O002
% mdparagraph:P812:begin
\itempara{2}{Identify the speaker and addressee of ic and ðē in 176.5.}
% mdparagraph:P812:end

% output: NEW-U18-R02-O003
% mdparagraph:P813:begin
\itempara{3}{Distinguish him in 176.19 from him in 176.20.}
% mdparagraph:P813:end

\reserve{8\baselineskip}
% mdparagraph:P814:begin
\subsection*{R03}
% mdparagraph:P814:end

% item: NEW-U18-R03
% mdparagraph:P815:begin
Parse the following selected source expressions and explain their functions. Supplied lexical and quantity explanations remain visible.
% mdparagraph:P815:end

% output: NEW-U18-R03-O001
% mdparagraph:P816:begin
\itempara{1}{Compare the grammatical subjects of geseah 175.18 and gefēoll 175.20.}
% mdparagraph:P816:end

% output: NEW-U18-R03-O002
% mdparagraph:P817:begin
\itempara{2}{Explain the relationships expressed by punda and goldes in 176.6.}
% mdparagraph:P817:end

% output: NEW-U18-R03-O003
% mdparagraph:P818:begin
\itempara{3}{Use original note 13 to explain ðe hīe gesāwon in 176.13 and the role of giefa.}
% mdparagraph:P818:end

% output: NEW-U18-R03-O004
% mdparagraph:P819:begin
\itempara{4}{Compare the case, number and agreement of gesunde in 176.16 and 176.18.}
% mdparagraph:P819:end

\reserve{8\baselineskip}
% mdparagraph:P820:begin
\subsection*{R04}
% mdparagraph:P820:end

% item: NEW-U18-R04
% mdparagraph:P821:begin
Describe the clause relations in the following spans using the complete source and its local support.
% mdparagraph:P821:end

% output: NEW-U18-R04-O001
% mdparagraph:P822:begin
\itempara{1}{Analyse the opening sentence 175.18–20: identify the temporal frame, the seeing-content clause and the main event.}
% mdparagraph:P822:end

% output: NEW-U18-R04-O002
% mdparagraph:P823:begin
\itempara{2}{Compare ðæt ic mōste 176.1 with swā-hwæt-swā 176.2 and 176.3–4.}
% mdparagraph:P823:end

% output: NEW-U18-R04-O003
% mdparagraph:P824:begin
\itempara{3}{Analyse the final ðæt clause 177.2 and explain how it relates to gān wē sēcean.}
% mdparagraph:P824:end

\reserve{8\baselineskip}
% mdparagraph:P825:begin
\subsection*{R05}
% mdparagraph:P825:end

% item: NEW-U18-R05
% mdparagraph:P826:begin
Explain the passage’s sequence from the actual statements, distinguishing completed actions from proposals.
% mdparagraph:P826:end

% output: NEW-U18-R05-O001
% mdparagraph:P827:begin
\itempara{1}{What authority does the princess give for her gifts, and who selects them?}
% mdparagraph:P827:end

% output: NEW-U18-R05-O002
% mdparagraph:P828:begin
\itempara{2}{What does the final speech establish about Apollonius’s next action, and what remains unnarrated at this boundary?}
% mdparagraph:P828:end

\reserve{8\baselineskip}
% mdparagraph:P829:begin
\section*{Source and editorial notes}
% mdparagraph:P829:end

\begingroup\fontsize{10}{12.8}\selectfont
\setlength{\parskip}{1.5pt}
% mdparagraph:P830:begin
The English prompts are Cook, Exercises in Old English, joint Group XII, XIII, Exercises 44–48, printed pages 29–33 (PDF pages 37–41). The original prose is Cook, A First Book in Old English, second edition, Reader XII, printed 175.18–177.2 (PDF 191–193). Modern explanations and lettered tasks are separate teaching material. Note numbers restart by printed page; P175, P176 and P177 identify the page.
% mdparagraph:P830:end

% mdparagraph:P831:begin
Preserve the selected spellings, including unmarked æcern, læfel and feoh, reading lęcgeað, and ðēowa without added r. The strong servant noun and the acorn plural have separately checked dictionary and manuscript-edition support. Derived paradigm cells apply the given stem and ending rules; their display does not assert that every cell occurs in the passage.
% mdparagraph:P831:end

% mdparagraph:P832:begin
The note on getogen raises an etymological question about modern wanton; here getogen means educated or accomplished. The hǣse note points to modern behest; hǣs here means command. These historical prompts are not additional exercises. The Beowulf comparison describes the company rising and one man greeting another and wishing him health; it illustrates farewell without assigning a poetry reading.
% mdparagraph:P832:end

% mdparagraph:P833:begin
The Latin notes mean a little earlier (176 note 1), teacher or master (note 5), enslaved people (note 10), dining room (note 11), saying farewell (note 15), and farewell or be well (note 17). Note 2 recommends an English infinitive for the finite Old English permission clause. Note 4 observes that the going-out detail is absent from the Latin. Note 9 records a manuscript adjective variant; the learner text retains Cook’s edited dēorwurðes. The other numbered grammatical references are explained in the modern reading help.
% mdparagraph:P833:end

% mdparagraph:P834:begin
\textsuperscript{P175N10}: See getēon. What relation has getogen to Mod. Eng. wanton?
% mdparagraph:P834:end

% mdparagraph:P835:begin
\textsuperscript{P175N11}: See 55.
% mdparagraph:P835:end

% mdparagraph:P836:begin
\textsuperscript{P176N1}: Lat. paulo ante. See 178.
% mdparagraph:P836:end

% mdparagraph:P837:begin
\textsuperscript{P176N2}: Translate by the infinitive sign, to. The OE. follows the Latin.
% mdparagraph:P837:end

% mdparagraph:P838:begin
\textsuperscript{P176N3}: See 197.
% mdparagraph:P838:end

% mdparagraph:P839:begin
\textsuperscript{P176N4}: Not in Latin.
% mdparagraph:P839:end

% mdparagraph:P840:begin
\textsuperscript{P176N5}: Lat. magister.
% mdparagraph:P840:end

% mdparagraph:P841:begin
\textsuperscript{P176N6}: See 154. c.
% mdparagraph:P841:end

% mdparagraph:P842:begin
\textsuperscript{P176N7}: See 153. f.
% mdparagraph:P842:end

% mdparagraph:P843:begin
\textsuperscript{P176N8}: See 174.
% mdparagraph:P843:end

% mdparagraph:P844:begin
\textsuperscript{P176N9}: MS. deorwurðan.
% mdparagraph:P844:end

% mdparagraph:P845:begin
\textsuperscript{P176N10}: Lat. servos.
% mdparagraph:P845:end

% mdparagraph:P846:begin
\textsuperscript{P176N11}: Lat. triclinio.
% mdparagraph:P846:end

% mdparagraph:P847:begin
\textsuperscript{P176N12}: See the derivation of Mod. Eng. behest.
% mdparagraph:P847:end

% mdparagraph:P848:begin
\textsuperscript{P176N13}: Refers to męnn.
% mdparagraph:P848:end

% mdparagraph:P849:begin
\textsuperscript{P176N14}: So in Beowulf (653–655): “Werod eall ārās; grētte ðā . . . guma ōðerne, . . . and him hǣl ābēad.”
% mdparagraph:P849:end

% mdparagraph:P850:begin
\textsuperscript{P176N15}: Lat. vale dicentes.
% mdparagraph:P850:end

% mdparagraph:P851:begin
\textsuperscript{P176N16}: See 153. d.
% mdparagraph:P851:end

% mdparagraph:P852:begin
\textsuperscript{P176N17}: Lat. valete.
% mdparagraph:P852:end

% mdparagraph:P853:begin
\textsuperscript{P176N18}: See 188.
% mdparagraph:P853:end

% mdparagraph:P854:begin
\textsuperscript{P177N1}: See 193. a.
% mdparagraph:P854:end

% mdparagraph:P855:begin
\textsuperscript{P177N2}: See 184. b.
% mdparagraph:P855:end

\endgroup
\end{document}
